% Acides forts et bases fortes — Chimie Tle E — Cours signé OnBuch+
% Programme officiel MINESEC, Terminale C, D, E. Compiler avec : tectonic L07-acides-forts-et-bases-fortes.tex
\newcommand{\DOCMATIERE}{Chimie}
\newcommand{\DOCNIVEAU}{Tle E}
\newcommand{\DOCMODULE}{Module 2 · Acides et bases}
\newcommand{\DOCLECON}{07}
\newcommand{\DOCTITRE}{Acides forts et bases fortes}
% OnBuch+ — préambule « cours signés » ENRICHI (compilé avec tectonic / XeLaTeX)
% Système visuel « Pop » d'OnBuch+ : fond crème (jamais blanc), bordures encre
% épaisses, ombres « dures » décalées, titres Archivo Black en capitales.
% Version enrichie pour les cours de chimie : chimie (mhchem, chemfig),
% graphes (pgfplots), boîtes pédagogiques supplémentaires (définition,
% propriété, méthode, attention, le savais-tu, exercice résolu, exercice,
% TP, bilan), page de garde refaite et signature OnBuch+ ancrée en bas.
%
% Macros attendues dans main.tex AVANT \input{preamble} :
%   \DOCMATIERE  \DOCNIVEAU  \DOCMODULE  \DOCLECON (numéro)  \DOCTITRE
\documentclass[11pt]{article}

\usepackage{fontspec}
\usepackage[french]{babel}
\usepackage[a4paper,margin=2cm,top=3.4cm,bottom=2.8cm,headheight=1.4cm,footskip=1.3cm]{geometry}
\usepackage{amsmath,amssymb,amsfonts}
\usepackage[table]{xcolor} % [table] : fournit aussi \rowcolors (alternance de lignes)
\usepackage{pagecolor}
\usepackage{tikz}
\usetikzlibrary{arrows.meta,positioning,calc,shapes.geometric,shapes.misc,shapes.arrows,decorations.pathmorphing,decorations.markings,patterns,shadows,mindmap,trees,fit,backgrounds,matrix,intersections}
\usepackage{pgfplots}
\pgfplotsset{compat=1.18}
\usepackage[version=4]{mhchem}
\usepackage{chemfig}
\usepackage{enumitem}
\usepackage{fancyhdr}
% [most] charge toutes les bibliothèques tcolorbox (listings, minted, raster,
% documentation...) et gonfle inutilement la mémoire TeX sur un document long
% (~300 boîtes) : on ne charge que ce qui est réellement utilisé.
\usepackage[skins,breakable]{tcolorbox}
\usepackage{graphicx}
\usepackage{colortbl}
\usepackage{booktabs}
\usepackage{tabularx}
\usepackage{array}
\usepackage{multicol}
\usepackage{siunitx}
% parse-numbers=false : accepte aussi \SI{2,5 \times 10^{-3}}{...}
\sisetup{locale=FR,output-decimal-marker={,},group-minimum-digits=4,parse-numbers=false}
\usepackage{qrcode}
% \degree : commande fréquemment utilisée par les modèles pour un angle ou une
% température en dehors de \SI{}{} (non fournie par les paquets chargés).
\providecommand{\degree}{^{\circ}}
\usepackage{lastpage}
\usepackage{hyperref}
\hypersetup{hidelinks}

% ============================================================
% Palette « Pop » OnBuch+ — mêmes hex que la classe P (plus_ui.dart)
% ============================================================
\definecolor{poporange}{HTML}{F59321}
\definecolor{popdark}{HTML}{E07A0C}
\definecolor{popcream}{HTML}{FFF6E8}
\definecolor{popink}{HTML}{1C1714}
\definecolor{poppurple}{HTML}{7A5AE0}
\definecolor{popgreen}{HTML}{2E9E6A}
\definecolor{poppink}{HTML}{E0457B}
\definecolor{popblue}{HTML}{2F7DE1}
\definecolor{popgold}{HTML}{C98A12}
\definecolor{popmuted}{HTML}{8A7F76}
\definecolor{popline}{HTML}{EADFD2}
\definecolor{popgray}{HTML}{8A7F76}
\colorlet{popgrey}{popgray}
% Alias tolérants (le modèle nomme parfois les couleurs différemment)
\colorlet{popwhite}{white}
\colorlet{popblack}{popink}
\colorlet{popred}{poppink}
\colorlet{popyellow}{popgold}
% Teintes claires pour fonds de tableaux / zones
\colorlet{popblueL}{popblue!10!white}
\colorlet{poporangeL}{poporange!14!white}
\colorlet{popgreenL}{popgreen!12!white}
\colorlet{poppurpleL}{poppurple!12!white}
\colorlet{poppinkL}{poppink!12!white}
\colorlet{popgoldL}{popgold!14!white}

\pagecolor{popcream}

% ============================================================
% Typographie
% ============================================================
\defaultfontfeatures{Ligatures=TeX}
\setmainfont{PlusJakartaSans-Medium.ttf}[Path=../../../fonts/,BoldFont=PlusJakartaSans-Bold.ttf]
% Maths en sans-serif (Fira Math) avec chiffres et lettres droites de Plus
% Jakarta, pour que les formules se fondent dans le texte.
\usepackage[math-style=ISO,bold-style=ISO]{unicode-math}
\setmathfont{FiraMath-Regular.otf}[Path=../../../fonts/]
\setmathfont{PlusJakartaSans-Medium.ttf}[Path=../../../fonts/,range={up/num,up/{Latin,latin},\mathup}]
\usepackage{icomma} % virgule décimale sans espace en mode math
% Caractères absents de la police de texte (sinon carré vide) : rendus par
% leur équivalent mathématique, en texte comme en formule.
\usepackage{newunicodechar}
\newunicodechar{α}{\ensuremath{\alpha}}
\newunicodechar{Α}{\ensuremath{\alpha}}
\newunicodechar{β}{\ensuremath{\beta}}
\newunicodechar{δ}{\ensuremath{\delta}}
\newunicodechar{✓}{\popcheck{popgreen}}
\newfontfamily\popdisplay{ArchivoBlack-Regular.ttf}[Path=../../../fonts/]
\newfontfamily\poplogo{SpaceGrotesk-Bold.ttf}[Path=../../../fonts/]
\newfontfamily\popsemi{PlusJakartaSans-SemiBold.ttf}[Path=../../../fonts/]
\newfontfamily\popxbold{PlusJakartaSans-ExtraBold.ttf}[Path=../../../fonts/]
\newfontfamily\popxboldls{PlusJakartaSans-ExtraBold.ttf}[Path=../../../fonts/,LetterSpace=3.0]

\setlist[enumerate]{leftmargin=1.7em,itemsep=5pt,topsep=4pt}
\setlist[itemize]{leftmargin=1.7em,itemsep=4pt,topsep=4pt,label={\color{poporange}\textbullet}}
\setlength{\parindent}{0pt}
\setlength{\parskip}{5pt}
\renewcommand{\arraystretch}{1.25}

% chemfig : liaisons un peu plus épaisses, encre Pop
\setchemfig{atom sep=2em,bond style={line width=0.9pt,popink},cram width=3pt}

% pgfplots : style Pop par défaut
\pgfplotsset{
  every axis/.append style={axis line style={popink,line width=1pt},tick style={popink},
    grid style={popline,line width=0.5pt},label style={font=\small},tick label style={font=\footnotesize}},
  popcurve/.style={poporange,line width=1.8pt,no markers,smooth},
}

% ============================================================
% Pill / squiggle
% ============================================================
\newcommand{\poppill}[3]{%
  \tikz[baseline=-0.55ex]\node[draw=popink,line width=1.2pt,rounded corners=2.6pt,%
    fill=#2,inner xsep=6.5pt,inner ysep=3pt]{%
    \popxboldls\fontsize{7}{7}\selectfont\color{#3}\MakeUppercase{#1}};%
}
\newcommand{\popsquiggle}[1]{%
  \tikz[baseline=-0.4ex]{
    \draw[#1,line width=2.6pt,line cap=round]
      plot [smooth] coordinates {(0,0.09) (0.34,0.01) (0.68,0.21) (1.02,0.01) (1.36,0.10)};
  }%
}
% Mot-clé surligné (marqueur orange sous le texte)
\newcommand{\cle}[1]{{\popxbold\color{popdark}#1}}

% ============================================================
% En-tête / pied
% ============================================================
\pagestyle{fancy}
\fancyhf{}
\renewcommand{\headrulewidth}{0pt}
\renewcommand{\footrulewidth}{0pt}
\lhead{\raisebox{-0.15ex}{\poplogo\fontsize{13}{13}\selectfont\color{popink}OnBuch\color{poporange}+}}
\rhead{\poppill{\DOCMATIERE\ \textbullet\ \DOCNIVEAU\ \textbullet\ Leçon \DOCLECON}{white}{popink}}
\lfoot{\popxbold\fontsize{7.6}{7.6}\selectfont\color{popmuted}\MakeUppercase{Cours signé OnBuch+}\ \textbullet\ {\popsemi\DOCTITRE}}
\rfoot{\tikz[baseline=-0.6ex]\node[rounded corners=8.5pt,fill=popink,minimum height=17pt,minimum width=17pt,inner xsep=5pt,inner sep=0]{\popxbold\fontsize{8}{8}\selectfont\color{popcream}\thepage\,/\,\pageref*{LastPage}};}

% ============================================================
% Titres
% ============================================================
\newcommand{\chaptitre}[2]{%
  \begin{tcolorbox}[enhanced,colback=poporange,colframe=popink,boxrule=2.6pt,arc=11pt,
      drop shadow=popink,left=15pt,right=15pt,top=12pt,bottom=11pt,before skip=3pt,after skip=18pt]
    \poppill{#2}{popink}{popcream}\par\vspace{9pt}
    {\raggedright\hyphenpenalty=10000\exhyphenpenalty=10000\popdisplay\fontsize{22}{24}\selectfont\color{popcream}\MakeUppercase{#1}\par}
  \end{tcolorbox}
}
% Section numérotée : pastille orange avec le numéro + titre Archivo
\newcounter{popsec}
\newcommand{\coursec}[1]{%
  \stepcounter{popsec}\setcounter{popsub}{0}%
  \par\vspace{16pt}\noindent
  \begin{minipage}{\linewidth}
  \tikz[baseline=-0.7ex]\node[draw=popink,line width=1.6pt,fill=poporange,rounded corners=4pt,minimum size=22pt,inner sep=0]{\popdisplay\fontsize{12}{12}\selectfont\color{popcream}\thepopsec};%
  \hspace{8pt}{\popdisplay\fontsize{15}{17}\selectfont\color{popink}\MakeUppercase{#1}}\par
  \vspace{4pt}\noindent{\color{poporange}\rule{36pt}{3.6pt}}
  \end{minipage}\par\nopagebreak\vspace{8pt}
}
\newcounter{popsub}
\newcommand{\courssub}[1]{%
  \stepcounter{popsub}%
  \par\vspace{9pt}\noindent{\popxbold\fontsize{11.5}{13}\selectfont\color{popink}{\color{poporange}\thepopsec.\thepopsub}\hspace{6pt}#1}\par\nopagebreak\vspace{4pt}
}

% ============================================================
% Boîtes pédagogiques — même recette : fond blanc, bordure épaisse colorée,
% ombre dure encre, badge plein. Titre optionnel [#1].
% ============================================================
\tcbset{popbase/.style={breakable,enhanced,colback=white,boxrule=2.3pt,arc=9pt,
  shadow={3pt}{-3pt}{0pt}{popink},left=14pt,right=14pt,top=11pt,bottom=11pt,before skip=11pt,after skip=15pt,
  fontupper=\popsemi\fontsize{10.6}{14.5}\selectfont\color{popink},
  fontlower=\popsemi\fontsize{10.6}{14.5}\selectfont\color{popink}}}
% Coche dessinée (le glyphe ✓ n'existe pas dans les polices du document)
\newcommand{\popcheck}[1]{\tikz[baseline=-0.5ex]\draw[#1,line width=1.6pt,line cap=round,line join=round](-0.13,0.0)--(-0.04,-0.09)--(0.13,0.1);}
\newcommand{\popboxhead}[3]{\poppill{#1}{#2}{white}\ifx\relax#3\relax\else\hspace{6pt}{\popxbold\fontsize{10.6}{12}\selectfont\color{popink}#3}\fi\par\vspace{7pt}}

\newtcolorbox{aretenir}[1][]{popbase,colframe=popgreen,before upper={\popboxhead{À retenir}{popgreen}{#1}}}
\newtcolorbox{exemplebox}[1][]{popbase,colframe=poporange,before upper={\popboxhead{Exemple}{poporange}{#1}}}
\newtcolorbox{definition}[1][]{popbase,colframe=popblue,before upper={\popboxhead{Définition}{popblue}{#1}}}
\newtcolorbox{propriete}[1][]{popbase,colframe=poppurple,before upper={\popboxhead{Propriété}{poppurple}{#1}}}
\newtcolorbox{methode}[1][]{popbase,colframe=popgold,before upper={\popboxhead{Méthode}{popgold}{#1}}}
\newtcolorbox{attention}[1][]{popbase,colframe=poppink,before upper={\popboxhead{Attention}{poppink}{#1}}}
\newtcolorbox{experience}[1][]{popbase,colframe=popblue,colback=popblueL,before upper={\popboxhead{Expérience}{popblue}{#1}}}
% « Le savais-tu ? » : carte violette pleine (contexte, culture, Cameroun)
\newtcolorbox{savaistu}[1][]{popbase,colback=poppurple,colframe=popink,
  fontupper=\popsemi\fontsize{10.4}{14.2}\selectfont\color{white},
  before upper={\poppill{Le savais-tu\,?}{white}{poppurple}\ifx\relax#1\relax\else\hspace{6pt}{\popxbold\color{white}#1}\fi\par\vspace{7pt}}}
% Exercice résolu : énoncé en haut, \tcblower puis la solution sur fond crème
\newcounter{popexo}\newcounter{popexr}
\newtcolorbox{exoresolu}[1][]{popbase,colframe=poporange,colbacklower=popcream,
  segmentation style={popink,line width=1.2pt,dashed},
  before upper={\stepcounter{popexr}\popboxhead{Exercice résolu \thepopexr}{poporange}{#1}},
  before lower={\poppill{Solution}{popink}{popcream}\par\vspace{6pt}}}
% Exercice (non résolu en place) — la correction est dans la partie Corrigés
\newtcolorbox{exercice}[1][]{popbase,colframe=popink,
  before upper={\stepcounter{popexo}\popboxhead{Exercice \thepopexo}{popink}{#1}}}
% Corrigé
\newtcolorbox{corrige}[1][]{popbase,colframe=popgreen,colback=popgreenL,
  before upper={\popboxhead{Corrigé}{popgreen}{#1}}}
% Objectifs / prérequis / bilan
\newtcolorbox{objectifs}{popbase,colframe=popink,colback=poporangeL,
  before upper={\popboxhead{Objectifs de la leçon}{popink}{}}}
\newtcolorbox{prerequis}{popbase,colframe=popink,colback=popblueL,
  before upper={\popboxhead{Prérequis}{popblue}{}}}
\newtcolorbox{bilan}{popbase,colframe=popink,colback=white,boxrule=2.8pt,
  before upper={\popboxhead{Fiche bilan}{poporange}{L'essentiel à connaître pour l'examen}}}
% Figure encadrée avec légende
\newtcolorbox{popfigure}[1][]{enhanced,breakable=false,colback=white,colframe=popline,boxrule=1.4pt,arc=8pt,
  left=10pt,right=10pt,top=10pt,bottom=8pt,before skip=10pt,after skip=12pt,halign=center,
  lower separated=false,halign lower=center,
  fontlower=\popsemi\fontsize{9}{11.5}\selectfont\color{popmuted},
  #1}
\newcommand{\legende}[1]{\par\vspace{6pt}{\popsemi\fontsize{9}{11.5}\selectfont\color{popmuted}#1}}

% ============================================================
% Signature OnBuch+
% ============================================================
\newcommand{\poplogoinline}[1]{{\poplogo\fontsize{#1}{#1}\selectfont\color{popink}OnBuch\color{poporange}+}}
\newcommand{\signatureonbuch}{%
  \begin{tcolorbox}[enhanced,colback=popink,colframe=popink,boxrule=2.2pt,arc=10pt,
      drop shadow=poporange,left=14pt,right=14pt,top=10pt,bottom=10pt,before skip=0pt,after skip=0pt]
    \tikz[baseline=-0.7ex]\node[circle,fill=popgreen,draw=popcream,line width=1.3pt,minimum size=22pt,inner sep=0]{\popcheck{white}};%
    \hspace{9pt}\begin{minipage}[c]{0.62\linewidth}
      {\popxbold\fontsize{12}{13}\selectfont\color{popcream}Signé\ \textbullet\ {\poplogo OnBuch\color{poporange}+}}\par\vspace{2pt}
      {\raggedright\popsemi\fontsize{8}{10.5}\selectfont\color{popline}\MakeUppercase{Équipe pédagogique OnBuch+}\\\MakeUppercase{Conforme au programme officiel MINESEC}\par}
    \end{minipage}\hfill
    {\poplogo\fontsize{22}{22}\selectfont\color{popcream}OnBuch\color{poporange}+}
  \end{tcolorbox}%
}

% ============================================================
% Page de garde
% #1 titre · #2 matière · #3 niveau · #4 module · #5 n° leçon · #6 durée ·
% #7 description / objectifs (texte ou itemize)
% ============================================================
\newcommand{\pagegarde}[7]{%
  \thispagestyle{empty}
  \newgeometry{a4paper,margin=1.7cm,top=1.6cm,bottom=1.4cm}%
  \begin{tikzpicture}[remember picture,overlay]
    \fill[poporange!18] ([xshift=-3.2cm,yshift=-3.6cm]current page.north east) circle (4.6cm);
    \fill[poppurple!14] ([xshift=2.4cm,yshift=7.6cm]current page.south west) circle (3.8cm);
  \end{tikzpicture}%
  {\poplogo\fontsize{30}{30}\selectfont\color{popink}OnBuch\color{poporange}+}\hfill
  \raisebox{0.6ex}{\poppill{Cours signé \textbullet\ Premium}{popink}{popcream}}\par
  \vspace{1pt}\popsquiggle{poppurple}\par
  \vspace{8pt}
  \begin{tcolorbox}[enhanced,colback=poporange,colframe=popink,boxrule=2.8pt,arc=13pt,
      drop shadow=popink,left=18pt,right=18pt,top=12pt,bottom=13pt,after skip=10pt]
    \poppill{#2 \textbullet\ #3}{popink}{popcream}\hspace{5pt}\poppill{#4}{white}{popink}\par\vspace{8pt}
    {\popdisplay\fontsize{14}{14}\selectfont\color{popink}LEÇON #5}\par\vspace{4pt}
    {\raggedright\hyphenpenalty=10000\exhyphenpenalty=10000\popdisplay\fontsize{25}{27}\selectfont\color{popcream}\MakeUppercase{#1}\par}
    \vspace{8pt}
    {\popxbold\fontsize{9.5}{9.5}\selectfont\color{popink}\MakeUppercase{Durée conseillée : #6}}
  \end{tcolorbox}
  \begin{tcolorbox}[enhanced,colback=white,colframe=popink,boxrule=2.2pt,arc=10pt,
      drop shadow=popink,left=16pt,right=16pt,top=10pt,bottom=10pt,after skip=8pt]
    \poppill{Dans cette leçon}{poppurple}{white}\par\vspace{5pt}
    \popsemi\fontsize{9.3}{12.6}\selectfont\color{popink}#7
  \end{tcolorbox}
  \noindent\begin{minipage}[t]{0.487\linewidth}
    \begin{tcolorbox}[enhanced,colback=popgreen,colframe=popink,boxrule=2.2pt,arc=10pt,
        drop shadow=popink,left=11pt,right=11pt,top=10pt,bottom=10pt,height=3.1cm,valign=center]
      \tcbox[colback=white,colframe=popink,boxrule=1.3pt,arc=5pt,boxsep=0pt,nobeforeafter,
        left=3pt,right=3pt,top=3pt,bottom=3pt,valign=center]{\qrcode[height=1.9cm]{https://play.google.com/store/apps/details?id=com.luvvix.onbuch&hl=fr}}%
      \hspace{8pt}\begin{minipage}[c]{0.5\linewidth}
        {\popdisplay\fontsize{11}{12.5}\selectfont\color{white}\MakeUppercase{Télécharge OnBuch}}\par\vspace{4pt}
        {\raggedright\popsemi\fontsize{8.4}{11}\selectfont\color{white}Gratuit \textbullet\ Google Play\\Tuteur IA, annales, concours, résultats\par}
      \end{minipage}
    \end{tcolorbox}
  \end{minipage}\hfill
  \begin{minipage}[t]{0.487\linewidth}
    \begin{tcolorbox}[enhanced,colback=poppurple,colframe=popink,boxrule=2.2pt,arc=10pt,
        drop shadow=popink,left=11pt,right=11pt,top=10pt,bottom=10pt,height=3.1cm,valign=center]
      \tikz[baseline=-0.7ex]\node[circle,fill=white,draw=popink,line width=1.3pt,minimum size=1.6cm,inner sep=0]{\popdisplay\fontsize{24}{24}\selectfont\color{poppurple}+};%
      \hspace{8pt}\begin{minipage}[c]{0.6\linewidth}
        {\popdisplay\fontsize{11}{12.5}\selectfont\color{white}\MakeUppercase{Passe à OnBuch+}}\par\vspace{4pt}
        {\raggedright\popsemi\fontsize{8.4}{11}\selectfont\color{white}Tous les cours signés, Tuteur IA illimité, fiches PDF — onglet OnBuch+ de l'app\par}
      \end{minipage}
    \end{tcolorbox}
  \end{minipage}
  \vspace{8pt}
  \popcontenu
  \vfill
  \signatureonbuch
  \restoregeometry
  \newpage
}

% Rangée « Ce cours contient » (page de garde)
\newcommand{\poptile}[3]{% #1 couleur #2 gros texte #3 légende
  \begin{tcolorbox}[enhanced,colback=#1,colframe=popink,boxrule=2pt,arc=9pt,shadow={2.5pt}{-2.5pt}{0pt}{popink},
      left=6pt,right=6pt,top=9pt,bottom=9pt,halign=center,valign=center,height=1.75cm,width=\linewidth,nobeforeafter]
    {\popdisplay\fontsize{12.5}{13}\selectfont\color{white}\MakeUppercase{#2}}\par\vspace{3pt}
    {\centering\popsemi\fontsize{7.8}{9.5}\selectfont\color{white}#3\par}
  \end{tcolorbox}}
\newcommand{\popcontenu}{%
  \par\noindent{\popxboldls\fontsize{7.5}{7.5}\selectfont\color{popmuted}CE COURS SIGNÉ CONTIENT}\par\vspace{7pt}
  \noindent\begin{minipage}[t]{0.235\linewidth}\poptile{popblue}{Cours}{complet, illustré, pas à pas}\end{minipage}\hfill
  \begin{minipage}[t]{0.235\linewidth}\poptile{poporange}{Méthodes}{et exercices résolus}\end{minipage}\hfill
  \begin{minipage}[t]{0.235\linewidth}\poptile{poppink}{Exercices}{type examen + corrigés}\end{minipage}\hfill
  \begin{minipage}[t]{0.235\linewidth}\poptile{popgreen}{Fiche bilan}{l'essentiel à retenir}\end{minipage}\par}

% Sommaire
\newtcolorbox{sommaire}{enhanced,colback=white,colframe=popink,boxrule=2.2pt,arc=10pt,
  drop shadow=popink,left=18pt,right=18pt,top=13pt,bottom=13pt,before skip=6pt,after skip=20pt,
  before upper={\poppill{Sommaire}{poppurple}{white}\par\vspace{9pt}\popsemi\fontsize{10.6}{14.5}\selectfont\color{popink}}}

% Signature de fin de cours — poussée tout en bas de la dernière page
\newcommand{\findecours}{%
  \par\vfill\nopagebreak
  \begin{center}\popsquiggle{poporange}\end{center}\vspace{4pt}
  \signatureonbuch
}

%% FIGFIX-BEGIN v4 — correctifs globaux des figures (docs/onbuch-generation/tools/figfix)
\usepackage{adjustbox}\usepackage{etoolbox}
% toute figure TikZ/pgfplots plus large que la ligne est réduite à la largeur disponible
\BeforeBeginEnvironment{tikzpicture}{\begin{adjustbox}{max width=\linewidth}}
\AfterEndEnvironment{tikzpicture}{\end{adjustbox}}
% légendes pgfplots lisibles par-dessus les courbes
\pgfplotsset{every axis/.append style={legend style={fill=white,fill opacity=0.92,text opacity=1,draw=black!15,inner sep=3pt}}}
% étiquettes d'arêtes (syntaxe edge["..."]) avec fond blanc
\tikzset{every edge quotes/.append style={fill=white,inner sep=1.2pt}}
%% FIGFIX-END
\begin{document}

\pagegarde{Acides forts et bases fortes}{Chimie}{Tle E}{Module 2 · Acides et bases}{07}{3 h (cours + TP)}{Cette leçon introduit la notion d'acide fort et de base forte à travers l'étude de l'acide chlorhydrique et de l'hydroxyde de sodium. Elle fournit les outils de calcul du pH, l'inventaire des espèces en solution et l'étude de l'effet de la dilution, compétences indispensables pour le Baccalauréat C, D et E.
\vspace{4pt}
\begin{multicols}{2}\small
\begin{itemize}
\item À la fin de cette leçon, je sais écrire l'équation de la réaction d'un acide fort ou d'une base forte avec l'eau et justifier son caractère total
\item À la fin de cette leçon, je sais montrer qu'un acide (ou une base) est fort(e) en comparant son pH à -log C (ou à pKe + log C)
\item À la fin de cette leçon, je sais calculer le pH d'une solution d'acide fort ou de base forte de concentration connue
\item À la fin de cette leçon, je sais dresser l'inventaire des espèces chimiques présentes en solution et calculer leurs concentrations par électroneutralité et conservation de la matière
\item À la fin de cette leçon, je sais prévoir l'effet d'une dilution au dixième sur le pH d'un acide fort ou d'une base forte
\item À la fin de cette leçon, je sais calculer le pH d'un mélange d'une solution d'acide fort et d'une solution de base forte
\end{itemize}\end{multicols}}

\begin{sommaire}
\begin{multicols}{2}
\begin{enumerate}
\item L'acide chlorhydrique, modèle d'acide fort
\item L'hydroxyde de sodium, modèle de base forte
\item Espèces chimiques en solution : inventaire et calcul des concentrations
\item Effet de la dilution sur le pH
\item Mélange d'un acide fort et d'une base forte
\item Autres acides forts et bases fortes usuels ; sécurité
\item Travaux pratiques
\item Méthodes et exercices résolus
\item Exercices
\item Corrigés des exercices
\item Fiche bilan
\end{enumerate}
\end{multicols}
\end{sommaire}

\begin{objectifs}
\begin{itemize}
\item À la fin de cette leçon, je sais écrire l'équation de la réaction d'un acide fort ou d'une base forte avec l'eau et justifier son caractère total
\item À la fin de cette leçon, je sais montrer qu'un acide (ou une base) est fort(e) en comparant son pH à -log C (ou à pKe + log C)
\item À la fin de cette leçon, je sais calculer le pH d'une solution d'acide fort ou de base forte de concentration connue
\item À la fin de cette leçon, je sais dresser l'inventaire des espèces chimiques présentes en solution et calculer leurs concentrations par électroneutralité et conservation de la matière
\item À la fin de cette leçon, je sais prévoir l'effet d'une dilution au dixième sur le pH d'un acide fort ou d'une base forte
\item À la fin de cette leçon, je sais calculer le pH d'un mélange d'une solution d'acide fort et d'une solution de base forte
\item À la fin de cette leçon, je sais préparer une solution par dilution et prévoir son pH
\item À la fin de cette leçon, je sais citer d'autres acides forts et bases fortes usuels et appliquer les règles de sécurité associées
\end{itemize}
\end{objectifs}

\begin{prerequis}
\begin{itemize}
\item Notion de pH et sa mesure (papier pH, pH-mètre) vue en Première
\item Fonction logarithme décimal : propriétés log(a×b), log(10\^{}x) = x
\item Produit ionique de l'eau : Ke = [H3O+][HO-] = 10\^{}-14 à 25 °C, pKe = 14
\item Concentration molaire, dilution et facteur de dilution (CiVi = CfVf)
\item Écriture des équations de dissolution et d'ionisation en solution aqueuse
\end{itemize}
\end{prerequis}

\newpage

\coursec{L'acide chlorhydrique, modèle d'acide fort}

Au marché Mokolo à Yaoundé, une vendeuse de jus de gingembre ajoute de la soude caustique pour nettoyer ses calebasses : quelques gouttes suffisent à rendre l'eau très corrosive. Dans le même quartier, un mécanicien manipule l'acide d'une batterie de voiture qui ronge le tissu en quelques secondes. Pourquoi certaines solutions acides ou basiques sont-elles si agressives alors que d'autres, comme le jus de citron dilué, sont inoffensives ? Comment les techniciens du laboratoire de l'hôpital central prévoient-ils le pH des solutions qu'ils préparent ? Cette leçon donne les clés pour comprendre et calculer le pH des acides et des bases dits « forts ».

\textbf{Problématique.} Comment prévoir le pH d'une solution d'acide chlorhydrique à partir de sa seule concentration ? Que devient exactement le chlorure d'hydrogène lorsqu'on le dissout dans l'eau ?

\courssub{Préparation de l'acide chlorhydrique}

\textbf{Intuition.} L'acide chlorhydrique n'est pas un liquide que l'on trouve tel quel dans la nature : c'est une \emph{solution aqueuse}, c'est-à-dire un mélange d'eau et d'une espèce dissoute. Cette espèce dissoute est le \cle{chlorure d'hydrogène} \ce{HCl}, un gaz incolore, âcre et très soluble dans l'eau (environ \SI{400}{L} de gaz dissous par litre d'eau à \SI{20}{\celsius}).

Pour préparer l'acide chlorhydrique, on fait barboter le gaz chlorure d'hydrogène dans de l'eau distillée. La dissolution s'accompagne d'une réaction chimique avec l'eau :

\[ \ce{HCl(g) + H2O(l) -> H3O+(aq) + Cl-(aq)} \]

Cette équation se lit ainsi : une molécule de chlorure d'hydrogène cède son proton \ce{H+} à une molécule d'eau ; il se forme un ion oxonium \ce{H3O+} (responsable de l'acidité) et un ion chlorure \ce{Cl-} (spectateur, sans propriété acido-basique notable). La réaction est \emph{exothermique} (elle dégage de la chaleur) et, nous allons le montrer, \emph{totale}.

\begin{attention}[Manipulation du gaz \ce{HCl} et pictogrammes]
Le chlorure d'hydrogène est un gaz très irritant pour les voies respiratoires et corrosif. Sa dissolution doit être réalisée sous hotte, avec un dispositif anti-retour (entonnoir à tulipe affleurant la surface de l'eau), car sa très grande solubilité peut provoquer une aspiration brutale du liquide dans le circuit de gaz.
Les pictogrammes de danger (SGH) associés à l'acide chlorhydrique concentré sont :
\begin{itemize}
\item \textbf{GHS05 (Corrosif)} : provoque de graves brûlures de la peau et des lésions oculaires.
\item \textbf{GHS07 (Danger)} : irritant pour les voies respiratoires.
\end{itemize}
Porter systématiquement gants, lunettes de protection et blouse.
\end{attention}

\courssub{Activité expérimentale : mesure du pH de solutions d'acide chlorhydrique}

\textbf{Question.} Le pH d'une solution d'acide chlorhydrique est-il lié simplement à sa concentration $C$ en soluté apporté ?

\begin{experience}[pH de solutions d'acide chlorhydrique]
\textbf{Protocole.} On prépare par dilutions successives trois solutions d'acide chlorhydrique de concentrations $C_1 = \SI{1,0e-1}{mol.L^{-1}}$, $C_2 = \SI{1,0e-2}{mol.L^{-1}}$ et $C_3 = \SI{1,0e-3}{mol.L^{-1}}$. On mesure le pH de chacune à l'aide d'un pH-mètre préalablement étalonné (étalons pH 4,01 ; 7,00 ; 10,01), à \SI{25}{\celsius}.

\textbf{Résultats.}
\begin{center}
\begin{tabularx}{\linewidth}{|l|X|X|X|}
\hline
\rowcolor{popblueL} Concentration $C$ (\si{mol.L^{-1}}) & \num{1,0e-1} & \num{1,0e-2} & \num{1,0e-3} \\
\hline
$-\log C$ & 1,0 & 2,0 & 3,0 \\
\hline
pH mesuré & 1,0 & 2,0 & 3,0 \\
\hline
\end{tabularx}
\end{center}

\textbf{Interprétation.} Dans les trois cas, le pH mesuré est exactement égal à $-\log C$. Or, par définition, $\text{pH} = -\log[\ce{H3O+}]$. On en déduit que $[\ce{H3O+}] = C$ : chaque molécule de \ce{HCl} apportée a produit exactement un ion \ce{H3O+}. La réaction de \ce{HCl} avec l'eau est donc \emph{totale}.
\end{experience}

\begin{popfigure}
\begin{center}
\begin{tikzpicture}
\begin{axis}[
  width=0.85\linewidth, height=6.2cm,
  xlabel={$-\log C$}, ylabel={pH mesuré},
  xmin=0, xmax=4.5, ymin=0, ymax=4.5,
  grid=both, grid style={popline!40},
  legend pos=north west,
  axis line style={popdark},
  tick style={popdark},
  label style={popdark},
  tick label style={popdark},
]
\addplot[popcurve, domain=0:4.5, samples=50] {x};
\addlegendentry{pH $= -\log C$}
\addplot[only marks, mark=*, mark size=2.5pt, poporange] coordinates {(1,1) (2,2) (3,3)};
\addlegendentry{Points expérimentaux}
\end{axis}
\end{tikzpicture}
\end{center}
\legende{pH mesuré en fonction de $-\log C$ pour l'acide chlorhydrique : les points expérimentaux s'alignent sur la droite de pente 1, preuve que la réaction de \ce{HCl} avec l'eau est totale.}
\end{popfigure}

\courssub{Un acide fort : définition et critère expérimental}

\begin{definition}[Acide fort]
Un \cle{acide fort} est un acide dont la réaction avec l'eau est \textbf{totale} : son taux d'avancement final vaut $\tau = 1$ (ou avancement final $x_f = n_{\text{apporté}}$). Toutes les molécules d'acide apportées sont transformées en ions. On écrit la réaction avec une \emph{flèche simple} :
\[ \ce{HCl + H2O -> H3O+ + Cl-} \]
\end{definition}

\textbf{Conséquence immédiate.} Si la réaction est totale, l'avancement final vaut $x_f = n_{\text{apporté}} = C \cdot V$, donc :
\[ [\ce{H3O+}] = \frac{x_f}{V} = C \qquad\text{et}\qquad [\ce{Cl-}] = C. \]
En prenant l'opposé du logarithme décimal :
\[ \text{pH} = -\log[\ce{H3O+}] = -\log C. \]

\begin{aretenir}[Relation fondamentale de l'acide fort]
Pour une solution d'acide fort de concentration $C$ (en \si{mol.L^{-1}}) :
\[ \boxed{\ \text{pH} = -\log C \ } \qquad\text{valable si } C > \SI{e-6}{mol.L^{-1}}. \]
En dessous de \SI{e-6}{mol.L^{-1}}, les ions \ce{H3O+} issus de l'autoprotolyse de l'eau (environ \SI{e-7}{mol.L^{-1}} à \SI{25}{\celsius}) ne sont plus négligeables et la formule donne des résultats absurdes (un pH supérieur à 7 pour un acide !).
\end{aretenir}

\begin{methode}[Vérifier expérimentalement qu'un acide est fort]
\begin{enumerate}
\item Mesurer le pH de la solution avec un pH-mètre étalonné (à température connue, idéalement \SI{25}{\celsius}).
\item Calculer $-\log C$ à partir de la concentration $C$ en soluté apporté (concentration formelle).
\item Comparer : si $\text{pH} = -\log C$ (à l'incertitude de mesure près, généralement $\pm 0,05$ unité de pH), alors $[\ce{H3O+}] = C$ et l'acide est \textbf{fort}. Si $\text{pH} > -\log C$, la réaction est partielle : l'acide est \textbf{faible}.
\end{enumerate}
\end{methode}

\begin{exemplebox}[Calcul direct du pH]
Une solution d'acide chlorhydrique a pour concentration $C = \SI{5,0e-3}{mol.L^{-1}}$. Comme $C > \SI{e-6}{mol.L^{-1}}$, on applique la relation :
\[ \text{pH} = -\log(\num{5,0e-3}) = -(\log 5{,}0 + \log 10^{-3}) = 3 - 0{,}70 = 2{,}30. \]
Vérification inverse : $[\ce{H3O+}] = 10^{-\text{pH}} = 10^{-2{,}30} \approx \SI{5,0e-3}{mol.L^{-1}} = C$. L'acide est bien fort.
\end{exemplebox}

\courssub{Que contient réellement la solution ?}

Puisque la réaction est totale, il ne reste \emph{aucune molécule} \ce{HCl} dans la solution. L'inventaire des espèces dissoutes est : ions oxonium \ce{H3O+}, ions chlorure \ce{Cl-}, molécules d'eau (solvant, en très large excès), et des traces d'ions hydroxyde \ce{HO-} provenant de l'autoprotolyse de l'eau.

\begin{popfigure}
\begin{center}
\begin{tikzpicture}[scale=0.95]
% Bécher
\draw[thick, popdark] (-3.2,0) -- (-3.2,3.4) (3.2,0) -- (3.2,3.4) (-3.2,0) -- (3.2,0);
% Niveau liquide
\draw[popblue!50, fill=popblue!10] (-3.15,0.05) rectangle (3.15,2.6);
% Ions H3O+ (orange)
\foreach \x/\y in {-2.2/0.8, -0.6/1.9, 1.0/0.7, 2.3/1.8, -1.0/2.2, 0.5/1.1}{
  \node[circle, fill=poporange, text=white, font=\scriptsize\bfseries, minimum size=0.8cm] at (\x,\y) {\ce{H3O+}};
}
% Ions Cl- (vert)
\foreach \x/\y in {-1.4/1.4, 0.2/0.6, 1.8/1.2, -2.6/2.1, 2.0/2.3, -0.8/0.4}{
  \node[circle, fill=popgreen, text=white, font=\scriptsize\bfseries, minimum size=0.75cm] at (\x,\y) {\ce{Cl-}};
}
% Molécules d'eau (bleu clair, plus petites)
\foreach \x/\y in {-0.9/2.35, 0.8/2.25, 2.6/0.55, -1.9/0.45, -2.5/1.0, 1.5/1.5}{
  \node[circle, fill=popblue!60, text=white, font=\tiny\bfseries, minimum size=0.45cm] at (\x,\y) {\ce{H2O}};
}
\node[popdark, font=\small\itshape] at (0,-0.6) {Solution : uniquement des ions \ce{H3O+} et \ce{Cl-} hydratés dans l'eau};
% Flèche dissolution
\draw[-{Stealth[length=4mm]}, very thick, poppurple] (-6.8,1.6) -- (-4.2,1.6);
\node[above, poppurple, font=\small\bfseries] at (-5.5,1.7) {dissolution + réaction};
% Molécule HCl gazeuse
\node[circle, fill=popmuted, text=white, font=\small\bfseries, minimum size=1.1cm] at (-8.2,1.6) {\ce{HCl}};
\node[below, popmuted, font=\small\itshape] at (-8.2,0.9) {gaz};
\end{tikzpicture}
\end{center}
\legende{Dissolution du chlorure d'hydrogène : chaque molécule \ce{HCl} réagit totalement avec l'eau pour donner un ion oxonium \ce{H3O+} et un ion chlorure \ce{Cl-}, entourés (hydratés) par des molécules d'eau.}
\end{popfigure}

\begin{attention}[Erreur fréquente : écrire \ce{HCl} dans l'inventaire des espèces]
Dans une solution d'acide chlorhydrique, il est \textbf{interdit} d'écrire \ce{HCl(aq)} parmi les espèces présentes : le chlorure d'hydrogène n'existe plus sous forme moléculaire, il est \emph{totalement ionisé}. On écrira $[\ce{H3O+}] = [\ce{Cl-}] = C$. Écrire « la solution contient des molécules \ce{HCl} » est une erreur classique sanctionnée au Baccalauréat.
\end{attention}

\begin{exoresolu}[pH et inventaire d'une solution d'acide chlorhydrique]
On dissout \SI{0,365}{g} de chlorure d'hydrogène gazeux dans de l'eau distillée de façon à obtenir $V = \SI{500}{mL}$ de solution. On donne $M(\ce{HCl}) = \SI{36,5}{g.mol^{-1}}$.
\begin{enumerate}
\item Écrire l'équation de la réaction de \ce{HCl} avec l'eau.
\item Calculer la concentration $C$ de la solution, puis son pH.
\item Donner les concentrations de toutes les espèces ioniques présentes à \SI{25}{\celsius}.
\end{enumerate}
\tcblower
\textbf{1.} Le chlorure d'hydrogène est un acide fort, sa réaction avec l'eau est totale :
\[ \ce{HCl(g) + H2O(l) -> H3O+(aq) + Cl-(aq)} \]

\textbf{2.} Quantité de matière apportée :
\[ n = \frac{m}{M} = \frac{0{,}365}{36{,}5} = \SI{1,00e-2}{mol}. \]
Concentration de la solution (concentration formelle) :
\[ C = \frac{n}{V} = \frac{\num{1,00e-2}}{\num{0,500}} = \SI{2,0e-2}{mol.L^{-1}}. \]
Comme $C > \SI{e-6}{mol.L^{-1}}$, la relation de l'acide fort s'applique :
\[ \text{pH} = -\log C = -\log(\num{2,0e-2}) = 2 - \log 2{,}0 = 1{,}70. \]

\textbf{3.} \textbf{Tableau d'avancement} (réaction totale, $x_f = n_{\text{apporté}}$) :
\begin{center}
\begin{tabular}{|c|c|c|c|c|}
\hline
\rowcolor{popblueL} & \ce{HCl} & \ce{H2O} & \ce{H3O+} & \ce{Cl-} \\
\hline
État initial & \num{1,00e-2} mol & excès & 0 & 0 \\
\hline
État final & 0 & excès & \num{1,00e-2} mol & \num{1,00e-2} mol \\
\hline
\end{tabular}
\end{center}
Concentrations à l'équilibre (volume $V = \SI{0,500}{L}$) :
\begin{itemize}
\item $[\ce{H3O+}] = \dfrac{\num{1,00e-2}}{\num{0,500}} = \SI{2,0e-2}{mol.L^{-1}}$
\item $[\ce{Cl-}] = \dfrac{\num{1,00e-2}}{\num{0,500}} = \SI{2,0e-2}{mol.L^{-1}}$
\end{itemize}
Ions hydroxyde issus de l'autoprotolyse de l'eau ($K_e = \num{1,0e-14}$ à \SI{25}{\celsius}) :
\[ [\ce{HO-}] = \frac{K_e}{[\ce{H3O+}]} = \frac{\num{1,0e-14}}{\num{2,0e-2}} = \SI{5,0e-13}{mol.L^{-1}}, \]
concentration totalement négligeable devant $C$, ce qui confirme a posteriori notre hypothèse.

\textbf{Vérification de l'électroneutralité de la solution :}
\[ [\ce{H3O+}] = [\ce{Cl-}] + [\ce{HO-}] \]
\[ \num{2,0e-2} \approx \num{2,0e-2} + \num{5,0e-13} \quad \text{(relation satisfaite)}. \]
\end{exoresolu}

\begin{savaistu}[L'acide chlorhydrique dans ton estomac]
L'estomac humain sécrète en permanence de l'acide chlorhydrique : le suc gastrique a un pH voisin de 1,5, soit $[\ce{H3O+}] \approx \SI{3e-2}{mol.L^{-1}}$. Cette acidité extrême active les enzymes digestives (pepsine) et détruit la plupart des microbes avalés avec les aliments. La paroi de l'estomac se protège grâce à un mucus riche en ions hydrogénocarbonate \ce{HCO3-} ; lorsque cette protection faiblit, l'acide attaque la muqueuse : c'est le reflux gastrique, que l'on soulage avec des antiacides contenant des bases (hydrogénocarbonate de sodium, hydroxyde de magnésium).
\end{savaistu}

\begin{aretenir}[L'essentiel de la section]
\begin{itemize}
\item L'acide chlorhydrique est une solution aqueuse obtenue par dissolution du gaz \ce{HCl} dans l'eau : \ce{HCl + H2O -> H3O+ + Cl-} (réaction totale, flèche simple).
\item Un \cle{acide fort} réagit totalement avec l'eau ($\tau = 1$) : il n'existe pas sous forme moléculaire en solution.
\item Pour un acide fort de concentration $C$ : $[\ce{H3O+}] = C$ et $\text{pH} = -\log C$, valable si $C > \SI{e-6}{mol.L^{-1}}$.
\item Critère expérimental : un acide est fort si le pH mesuré est égal à $-\log C$.
\item Inventaire des espèces : \ce{H3O+}, \ce{Cl-}, \ce{H2O}, traces \ce{HO-}. Pas de \ce{HCl} moléculaire.
\item Respecter les consignes de sécurité (pictogrammes GHS05, GHS07) lors de la manipulation.
\end{itemize}
\end{aretenir}

\coursec{L'hydroxyde de sodium, modèle de base forte}

Dans la section précédente, nous avons découvert l'acide chlorhydrique, modèle de l'\cle{acide fort} : sa réaction avec l'eau est totale, et son pH obéit à la loi simple $pH = -\log C$. Mais la chimie des solutions aqueuses ne serait pas complète sans son pendant : les \cle{bases}. Dans les savonneries artisanales de Maroua ou de Bafoussam, on utilise depuis des générations une substance blanche, vendue en pastilles ou en perles : la \emph{soude}, ou hydroxyde de sodium \ce{NaOH}. Une solution de soude « monte » très haut sur l'échelle des pH. Comment relier quantitativement son pH à sa concentration ? C'est tout l'objet de cette section : nous allons montrer que la soude est une \cle{base forte}, établir la relation $pH = pK_e + \log C$, et apprendre à manipuler ce produit en toute sécurité.

\courssub{Dissolution de la soude dans l'eau}

\textbf{Situation concrète.}\par
On dissout quelques pastilles de soude dans de l'eau distillée contenue dans un bécher. Deux observations frappent immédiatement l'expérimentateur : le bécher chauffe (la dissolution est fortement \emph{exothermique}) et le papier pH trempé dans la solution vire au bleu foncé, indice d'un pH très élevé. Que s'est-il passé à l'échelle microscopique ?

L'hydroxyde de sodium solide \ce{NaOH(s)} est un \cle{cristal ionique} : il est déjà constitué d'ions sodium \ce{Na+} et d'ions hydroxyde \ce{HO-}, régulièrement empilés. Lors de la dissolution, les molécules d'eau polaires arrachent ces ions au cristal et les entourent (on dit qu'elles les \emph{solvatent}). L'équation de dissolution s'écrit :

\[ \ce{NaOH(s) ->[eau] Na+(aq) + HO-(aq)} \]

Cette équation se lit : une formule unité de soude solide libère dans l'eau \emph{un} ion sodium et \emph{un} ion hydroxyde. La dissolution est \textbf{totale} : il ne reste aucun solide tant qu'on ne dépasse pas la solubilité.

\begin{attention}[Dissolution très exothermique]
La dissolution de la soude dans l'eau dégage beaucoup de chaleur : la température du bécher peut dépasser \SI{60}{\celsius} pour une solution concentrée. Ne jamais tenir le bécher à pleine main, agiter doucement, et ajouter \emph{toujours} les pastilles de soude \textbf{dans} l'eau (jamais l'inverse en grande quantité), par petites portions.
\end{attention}

\textbf{Les ions hydroxyde et l'eau : l'effet de nivellement.}\par
Les ions \ce{HO-} ainsi libérés sont la base active de la solution. Selon la théorie de Brønsted, une base est une espèce capable de capter un proton \ce{H+}. L'ion hydroxyde est la base \emph{la plus forte qui puisse exister dans l'eau} (effet de nivellement du solvant) : toute base théorique plus forte que \ce{HO-} (comme l'éthanolate \ce{C2H5O-} ou l'amide \ce{NH2-}) réagit \textbf{totalement} avec l'eau pour \emph{former} des ions \ce{HO-}. 

Pour la soude, la dissolution fournit directement des ions \ce{HO-}. Ceux-ci ne réagissent pas davantage avec l'eau (l'équilibre \ce{HO- + H2O <=> H2O + HO-} est une identité). La réaction de dissolution est donc \textbf{totale} : tous les ions hydroxyde apportés par la soude restent présents en solution, sous forme \ce{HO-} libre.

\courssub{Étude expérimentale : mesures de pH}

\begin{experience}[Mesure du pH de solutions de soude]
\textbf{Protocole.} On prépare par dilutions successives trois solutions de soude de concentrations $C_1 = \SI{1,0e-1}{mol.L^{-1}}$, $C_2 = \SI{1,0e-2}{mol.L^{-1}}$ et $C_3 = \SI{1,0e-3}{mol.L^{-1}}$. On mesure le pH de chacune avec un pH-mètre étalonné, à \SI{25}{\celsius}.

\textbf{Observations.} Les résultats sont consignés dans le tableau ci-dessous.

\begin{center}
\begin{tabularx}{\linewidth}{|l|X|X|X|}
\hline
\rowcolor{popblueL} Solution & $C$ (\si{mol.L^{-1}}) & pH mesuré & $-\log C$ \\
\hline
$S_1$ & $1,0\times 10^{-1}$ & $13,0$ & $1,0$ \\
\hline
$S_2$ & $1,0\times 10^{-2}$ & $12,0$ & $2,0$ \\
\hline
$S_3$ & $1,0\times 10^{-3}$ & $11,0$ & $3,0$ \\
\hline
\end{tabularx}
\end{center}

\textbf{Interprétation.} Contrairement à l'acide chlorhydrique, le pH ne vaut \emph{pas} $-\log C$ : il vaut $14 - (-\log C)$, c'est-à-dire $14 + \log C$. Chaque dilution au dixième fait \emph{baisser} le pH d'une unité. De plus, la concentration en ions hydroxyde mesurée (par dosage) est exactement égale à $C$ : la soude est \textbf{totalement ionisée}.
\end{experience}

\courssub{Définition d'une base forte}

\begin{definition}[Base forte]
Une \cle{base forte} est une base dont la réaction avec l'eau (ou la dissolution pour un hydroxyde ionique) est \textbf{totale}. En solution aqueuse, elle est \textbf{totalement ionisée} : chaque formule unité dissoute libère quantitativement des ions hydroxyde \ce{HO-}. Il ne subsiste en solution \emph{aucune} molécule (ou formule unité) de base non dissociée.
\end{definition}

Pour la soude, la conséquence immédiate est la relation fondamentale :

\begin{aretenir}[Concentration en ions hydroxyde]
Pour une solution de soude (base forte) de concentration apportée $C$ :
\[ [\ce{HO-}] = C \]
Tous les ions hydroxyde apportés se retrouvent intacts en solution. Les ions sodium \ce{Na+} sont des \emph{ions spectateurs} : $[\ce{Na+}] = C$ également.
\end{aretenir}

\courssub{Inventaire des espèces, électroneutralité et conservation de la matière}

Avant tout calcul de pH, le chimiste dresse l'inventaire rigoureux des espèces présentes. C'est une compétence clé du programme.

\begin{methode}[Inventaire et relations pour une solution de base forte]
Pour une solution de soude de concentration apportée $C$ :
\begin{enumerate}
\item \textbf{Espèces présentes :} \ce{H2O}, \ce{Na+}, \ce{HO-}, \ce{H3O+} (issus de l'autoprotolyse de l'eau).
\item \textbf{Électroneutralité :} La solution est globalement neutre.
      \[ [\ce{Na+}] + [\ce{H3O+}] = [\ce{HO-}] \]
\item \textbf{Conservation de la matière (pour le sodium) :} Tout le sodium vient de la soude dissoute.
      \[ [\ce{Na+}] = C \]
\item \textbf{Conséquence (approximation usuelle) :} Pour $C > 10^{-6}$ \si{mol.L^{-1}}, $[\ce{H3O+}] \ll [\ce{HO-}]$, donc $[\ce{HO-}] \approx [\ce{Na+}] = C$.
\end{enumerate}
\end{methode}

\courssub{Relation fondamentale : pH = pK\textsubscript{e} + log C}

\textbf{Démonstration pas à pas.}\par
Reprenons le produit ionique de l'eau, valable pour \emph{toute} solution aqueuse à \SI{25}{\celsius} :

\[ K_e = [\ce{H3O+}][\ce{HO-}] = 10^{-14} \quad \text{donc} \quad pK_e = 14 \]

Pour une solution de base forte de concentration $C$, nous venons d'établir que $[\ce{HO-}] = C$ (pour $C > 10^{-6}$ \si{mol.L^{-1}}). En injectant dans l'expression de $K_e$ :

\[ [\ce{H3O+}] = \frac{K_e}{[\ce{HO-}]} = \frac{10^{-14}}{C} \]

Appliquons maintenant la définition du pH, $pH = -\log[\ce{H3O+}]$ :

\begin{align*}
pH &= -\log\left(\frac{10^{-14}}{C}\right) \\
&= -\log(10^{-14}) + \log C \\
&= 14 + \log C
\end{align*}

La relation générale s'écrit donc :

\begin{aretenir}[pH d'une base forte]
Pour une solution de base forte de concentration $C$, à \SI{25}{\celsius} :
\[ pH = pK_e + \log C = 14 + \log C \]
\textbf{Condition de validité :} $C > 10^{-6}$ \si{mol.L^{-1}}. En dessous, les ions \ce{HO-} provenant de l'autoprotolyse de l'eau ne sont plus négligeables, et le pH tend vers $7$ (une dilution infinie d'une base ne rend jamais la solution acide !).
\end{aretenir}

\begin{popfigure}
\begin{center}
\begin{tikzpicture}
\begin{axis}[
    width=0.85\linewidth, height=6.5cm,
    xlabel={$\log C$}, ylabel={pH},
    xmin=-6, xmax=0, ymin=7, ymax=14.5,
    grid=major, grid style={popline!40},
    axis lines=left,
    tick label style={font=\small},
    label style={font=\small},
]
\addplot[popcurve,domain=-6:-1,samples=50]{14+x};
\addplot[poporange, dashed, domain=-6:-1, samples=2]{7};
\node[font=\small, popdark, anchor=south west] at (axis cs:-5.8,13.1) {$pH = 14 + \log C$};
\node[font=\footnotesize, poporange, anchor=south] at (axis cs:-3.5,7.1) {zone de non-validité ($C < 10^{-6}$)};
\addplot[only marks, mark=*, poppurple, mark size=2pt] coordinates {(-1,13) (-2,12) (-3,11)};
\end{axis}
\end{tikzpicture}
\end{center}
\legende{pH d'une solution de soude en fonction de $\log C$ à \SI{25}{\celsius} : droite de pente $+1$. Les points correspondent aux mesures de l'expérience. Pour $C < 10^{-6}$ \si{mol.L^{-1}}, le pH plafonne vers 7.}
\end{popfigure}

\begin{methode}[Calculer le pH d'une base forte]
\begin{enumerate}
\item Écrire l'équation de dissolution et vérifier qu'il s'agit bien d'une \textbf{base forte} (soude \ce{NaOH}, potasse \ce{KOH}, éthanolate de sodium \ce{C2H5ONa}...).
\item En déduire $[\ce{HO-}] = C$ (vérifier $C > 10^{-6}$ \si{mol.L^{-1}}).
\item Calculer le $pOH = -\log[\ce{HO-}] = -\log C$.
\item Conclure : $pH = pK_e - pOH = 14 + \log C$ (à \SI{25}{\celsius}).
\item Vérifier la cohérence : une base doit donner $pH > 7$.
\end{enumerate}
\end{methode}

\begin{exemplebox}[Soude à $C = \SI{2,0e-2}{mol.L^{-1}}$]
On dissout des pastilles de soude pour obtenir une solution de concentration $C = \SI{2,0e-2}{mol.L^{-1}}$ à \SI{25}{\celsius}.

Étape 1 : la soude est une base forte, donc $[\ce{HO-}] = C = \SI{2,0e-2}{mol.L^{-1}}$.

Étape 2 : $pOH = -\log(2,0\times 10^{-2}) = 2 - \log 2 = 2 - 0,30 = 1,70$.

Étape 3 : $pH = 14 - pOH = 14 - 1,70 = 12,3$.

Vérification directe : $pH = 14 + \log(2,0\times 10^{-2}) = 14 - 2 + 0,30 = 12,3$. Le pH est bien supérieur à 7 : cohérent pour une base.
\end{exemplebox}

\begin{attention}[Erreur fréquente : la mauvaise formule]
\textbf{Ne jamais appliquer $pH = -\log C$ à une base forte !} Cette formule est réservée aux \emph{acides forts}. Pour la soude à $C = \SI{1,0e-2}{mol.L^{-1}}$, elle donnerait $pH = 2$, alors que la solution est \emph{basique} ($pH = 12$). Avant tout calcul, pose-toi la question : « ai-je un acide ou une base ? ». Pour une base forte, la bonne réponse est $pH = 14 + \log C$.
\end{attention}

\begin{exoresolu}[pH et concentration d'une solution de déboucheur diluée]
Un déboucheur de canalisation contient de la soude. On prélève \SI{10,0}{mL} d'une solution commerciale que l'on dilue 100 fois. Le pH de la solution diluée vaut $11,7$ à \SI{25}{\celsius}. Déterminer la concentration $C_d$ de la solution diluée, puis la concentration $C_0$ de la solution commerciale.
\tcblower
La soude est une base forte : $pH = 14 + \log C_d$, donc :
\[ \log C_d = pH - 14 = 11,7 - 14 = -2,3 \]
\[ C_d = 10^{-2,3} = \SI{5,0e-3}{mol.L^{-1}} \]
La solution commerciale est 100 fois plus concentrée (facteur de dilution $F = 100$) :
\[ C_0 = 100 \times C_d = \SI{5,0e-1}{mol.L^{-1}} \]
Vérification : $C_0 > 10^{-6}$ \si{mol.L^{-1}}, la formule est valable. Le pH de la solution commerciale vaudrait $14 + \log(0,5) = 13,7$ : un milieu extrêmement corrosif, d'où les pictogrammes de danger sur le flacon !
\end{exoresolu}

\courssub{La soude, un produit corrosif : précautions de manipulation}

Les ions hydroxyde attaquent les protéines et les graisses de la peau (réaction de saponification des graisses cutanées) : la soude provoque des \textbf{brûlures chimiques graves}, d'autant plus dangereuses qu'elles sont peu douloureuses au début (la soude « savonne » les graisses de l'épiderme, d'où une sensation glissante caractéristique). Les projections dans les yeux peuvent entraîner la cécité.

\begin{attention}[Règles de sécurité avec la soude]
\begin{itemize}
\item Porter obligatoirement des \textbf{gants}, des \textbf{lunettes de protection} et une \textbf{blouse}.
\item Manipuler les pastilles avec une spatule, jamais à mains nues.
\item En cas de contact avec la peau : rincer abondamment à l'eau claire pendant plusieurs minutes, puis prévenir l'enseignant.
\item Ne jamais goûter ni sentir directement une solution de soude.
\end{itemize}
\end{attention}

\begin{popfigure}
\begin{center}
\begin{tikzpicture}[scale=1.0]
% Flacon
\draw[thick, popdark, fill=popcream] (-1.1,-2.2) rectangle (1.1,0.6);
\draw[thick, popdark, fill=popblueL] (-0.35,0.6) rectangle (0.35,1.0);
\draw[thick, popdark, fill=popmuted!40] (-0.45,1.0) rectangle (0.45,1.35);
\draw[popblue!50, fill=popblue!15] (-1.0,-2.1) rectangle (1.0,-1.1);
\node[font=\footnotesize\bfseries, popdark] at (0,0.25) {SOUDE};
\node[font=\scriptsize, popdark] at (0,-0.05) {\ce{NaOH}};
% Pictogramme GHS05
\begin{scope}[shift={(0,-1.55)}, scale=0.42]
\draw[line width=2.2pt, red, rotate=45] (-1.05,-1.05) rectangle (1.05,1.05);
\fill[white, rotate=45] (-0.98,-0.98) rectangle (0.98,0.98);
\draw[line width=1.6pt, red, rotate=45] (-1.05,-1.05) rectangle (1.05,1.05);
% deux tubes qui coulent
\draw[thick, black] (-0.55,0.55) -- (0.05,0.25) -- (0.05,0.05);
\draw[thick, black] (-0.55,0.75) -- (0.15,0.4) -- (0.15,0.2);
\fill[black] (0.10,-0.05) circle (0.05);
\fill[black] (0.10,-0.22) circle (0.04);
% surface attaquée
\draw[thick, black] (-0.6,-0.45) -- (0.6,-0.45);
\draw[thick, black] (-0.15,-0.45) -- (-0.05,-0.30) -- (0.05,-0.45);
\end{scope}
\node[font=\scriptsize\bfseries, popdark] at (0,-2.55) {GHS05 -- Corrosif};
% Mention de danger
\node[font=\scriptsize, poporange, anchor=west] at (1.5,-1.0) {\textbf{H314} :};
\node[font=\scriptsize, popdark, anchor=west, text width=4.2cm] at (1.5,-1.55) {Provoque de graves brûlures de la peau et des lésions oculaires.};
\end{tikzpicture}
\end{center}
\legende{Flacon de soude portant le pictogramme GHS05 (corrosif) : les gouttes qui tombent des deux tubes rongent la surface et la main. Ce symbole impose gants et lunettes.}
\end{popfigure}

\begin{savaistu}[La soude au Cameroun : du savon au déboucheur]
Dans de nombreuses localités du Cameroun, la fabrication artisanale du savon repose sur la \emph{saponification} : on fait chauffer de l'huile de palme rouge avec une lessive de soude (ou de potasse obtenue à partir de cendres de bois). Les corps gras réagissent avec les ions \ce{HO-} pour former du savon et du glycérol. La soude est aussi le principe actif des déboucheurs de canalisations vendus dans les marchés de Douala et Yaoundé : elle dissout les graisses et les matières organiques qui obstruent les tuyaux. Enfin, à l'échelle industrielle, la soude intervient dans le raffinage des huiles et la fabrication de la pâte à papier. Un produit dangereux, certes, mais indispensable -- à condition de respecter les règles de sécurité !
\end{savaistu}

\begin{aretenir}[L'essentiel de la section]
\begin{itemize}
\item La soude se dissout totalement dans l'eau : \ce{NaOH(s) -> Na+(aq) + HO-(aq)} (dissolution exothermique).
\item Une \cle{base forte} est totalement ionisée : $[\ce{HO-}] = C$ (pour $C > 10^{-6}$ \si{mol.L^{-1}}).
\item Inventaire : \ce{H2O}, \ce{Na+}, \ce{HO-}, \ce{H3O+} ; Électroneutralité : $[\ce{Na+}] + [\ce{H3O+}] = [\ce{HO-}]$ ; Conservation : $[\ce{Na+}] = C$.
\item À \SI{25}{\celsius} : $pH = pK_e + \log C = 14 + \log C$, valable pour $C > 10^{-6}$ \si{mol.L^{-1}}.
\item Diluer dix fois une base forte fait \emph{diminuer} son pH d'une unité.
\item La soude est corrosive (GHS05, H314) : gants, lunettes et blouse obligatoires.
\end{itemize}
\end{aretenir}

\coursec{Espèces chimiques en solution : inventaire et calcul des concentrations}

Dans les sections précédentes, tu as découvert deux modèles fondamentaux : l'acide chlorhydrique, \emph{acide fort} totalement ionisé en \ce{H3O+} et \ce{Cl-}, et l'hydroxyde de sodium, \emph{base forte} totalement dissociée en \ce{Na+} et \ce{HO-}. Mais une solution aqueuse n'est jamais aussi simple qu'elle n'y paraît : même dans un acide, il reste des ions \ce{HO-}, et même dans une base, il reste des ions \ce{H3O+}. L'objectif de cette section est de t'apprendre à \cle{inventorier} toutes les espèces présentes et à \cle{calculer} chacune de leurs concentrations. C'est une compétence clé du Baccalauréat : presque chaque sujet de chimie en Terminale C, D ou E commence par la question « faire l'inventaire des espèces chimiques présentes et calculer leurs concentrations ».

\courssub{L'inventaire des espèces : ne rien oublier}

\textbf{L'intuition d'abord.} Imagine que tu verses une pastille d'acide dans un verre d'eau de la fontaine du quartier. Ce que tu obtiens n'est pas seulement « de l'acide » : c'est une véritable population microscopique. Il y a les molécules d'eau, de très loin majoritaires ; les ions apportés par le soluté ; et, discrètement, les ions issus de l'eau elle-même. Oublier l'une de ces populations, c'est fausser tous les calculs qui suivent.

\begin{methode}[Règle d'or de l'inventaire]
Pour dresser l'inventaire des espèces chimiques présentes dans une solution aqueuse, procède toujours dans le même ordre :
\begin{enumerate}
\item \textbf{Le solvant} : l'eau \ce{H2O}, toujours présente et toujours \emph{très majoritaire} (environ \SI{55,5}{mol.L^{-1}} de molécules non ionisées) ;
\item \textbf{Les ions du soluté} : ceux produits par la dissolution et la réaction (totale) avec l'eau, par exemple \ce{H3O+} et \ce{Cl-} pour l'acide chlorhydrique ;
\item \textbf{Les ions de l'eau} : \emph{toute} solution aqueuse contient \ce{H3O+} \textbf{et} \ce{HO-}, issus de l'autoprotolyse de l'eau :
\[ \ce{2 H2O <=> H3O+ + HO-} \qquad K_e = [\ce{H3O+}][\ce{HO-}] = 10^{-14} \text{ à } \SI{25}{\celsius} \]
Cette réaction, très limitée, est cependant universelle : elle a lieu dans l'eau pure, dans les acides, dans les bases.
\end{enumerate}
\end{methode}

\begin{attention}[Le piège classique de l'inventaire]
L'erreur la plus fréquente au Baccalauréat est d'écrire, pour une solution d'acide chlorhydrique : « espèces présentes : \ce{H3O+} et \ce{Cl-} ». \textbf{C'est incomplet !} Il faut toujours ajouter \ce{HO-} et \ce{H2O}. Les ions \ce{HO-} existent dans une solution acide, certes en quantité infime, mais ils existent. Les oublier dans l'inventaire coûte des points ; les oublier dans l'équation d'électroneutralité fausse le raisonnement. Retiens : \emph{acide ou base, \ce{H3O+} et \ce{HO-} sont toujours là, tous les deux}.
\end{attention}

\courssub{Exemple 1 : solution d'acide chlorhydrique de concentration C}

Considérons une solution d'acide chlorhydrique de concentration apportée $C$. Le chlorure d'hydrogène \ce{HCl} réagit totalement avec l'eau (c'est précisément la définition d'un acide fort) :

\[ \ce{HCl + H2O -> H3O+ + Cl-} \qquad \text{(réaction totale, rapide, exothermique)} \]

\textbf{Inventaire complet :} \ce{H2O} (solvant, majoritaire), \ce{H3O+}, \ce{Cl-}, \ce{HO-}. Il ne reste \emph{aucune} molécule \ce{HCl} en solution, car la réaction est totale.

Pour trouver les concentrations, on dispose de trois outils qui, utilisés ensemble, permettent de tout calculer.

\begin{definition}[Équation d'électroneutralité]
Une solution est macroscopiquement \cle{électriquement neutre} : elle ne peut pas accumuler de charge. L'\cle{équation d'électroneutralité} traduit ce fait : la somme des concentrations des charges positives est égale à la somme des concentrations des charges négatives (chaque concentration étant multipliée par la valeur absolue de la charge de l'ion). Pour l'acide chlorhydrique :
\[ [\ce{H3O+}] = [\ce{Cl-}] + [\ce{HO-}] \]
\end{definition}

\begin{definition}[Équation de conservation de la matière]
La \cle{conservation de la matière} exprime qu'un élément chimique ne disparaît pas : tout le chlore apporté par \ce{HCl} se retrouve sous forme d'ions \ce{Cl-}. Comme la réaction est totale :
\[ [\ce{Cl-}] = C \]
\end{definition}

Le troisième outil est le produit ionique de l'eau, $K_e = [\ce{H3O+}][\ce{HO-}] = 10^{-14}$ à \SI{25}{\celsius}, qui relie les deux ions de l'eau.

\begin{methode}[Calculer les concentrations à partir du pH]
Toute solution aqueuse contient \ce{H3O+} et \ce{HO-}. Si le pH est connu, leurs concentrations s'obtiennent immédiatement, à \SI{25}{\celsius} :
\[ [\ce{H3O+}] = 10^{-\mathrm{pH}} \qquad\text{et}\qquad [\ce{HO-}] = \frac{K_e}{[\ce{H3O+}]} = 10^{\mathrm{pH}-14} \]
Ensuite, l'électroneutralité et la conservation de la matière donnent les concentrations des ions du soluté.
\end{methode}

\begin{exemplebox}[Acide chlorhydrique à $C = \SI{1,0e-2}{mol.L^{-1}}$]
Une solution d'acide chlorhydrique de concentration $C = \SI{1,0e-2}{mol.L^{-1}}$ a un pH mesuré égal à $2,0$.
\begin{itemize}
\item Conservation de la matière : $[\ce{Cl-}] = C = \SI{1,0e-2}{mol.L^{-1}}$ ;
\item $[\ce{H3O+}] = 10^{-\mathrm{pH}} = 10^{-2} = \SI{1,0e-2}{mol.L^{-1}}$ ;
\item $[\ce{HO-}] = 10^{\mathrm{pH}-14} = 10^{-12} = \SI{1,0e-12}{mol.L^{-1}}$.
\end{itemize}
\textbf{Vérification de l'électroneutralité :} $[\ce{Cl-}] + [\ce{HO-}] = 10^{-2} + 10^{-12} \approx 10^{-2} = [\ce{H3O+}]$. L'égalité est bien vérifiée.
\textbf{Conclusion importante :} $[\ce{HO-}] = 10^{-12} \ll [\ce{Cl-}] = 10^{-2}$ : les ions \ce{HO-} sont \emph{dix milliards de fois} moins nombreux que les ions chlorure. En milieu acide, $[\ce{HO-}]$ est \cle{négligeable} devant $[\ce{Cl-}]$, et l'électroneutralité se simplifie en $[\ce{H3O+}] \approx [\ce{Cl-}] = C$. On retrouve ainsi $\mathrm{pH} = -\log C$, la formule de l'acide fort.
\end{exemplebox}

\begin{popfigure}
\begin{tikzpicture}[scale=1.0]
% Bécher
\draw[thick, popblue] (0,0) -- (0,3.2) (4.6,0) -- (4.6,3.2);
\draw[thick, popblue] (0,0) .. controls (0.3,-0.35) and (4.3,-0.35) .. (4.6,0);
\draw[popblue!60, fill=popblue!12] (0.06,0.05) -- (0.06,2.5) -- (4.54,2.5) -- (4.54,0.05) .. controls (4.2,-0.28) and (0.4,-0.28) .. (0.06,0.05);
\draw[popblue!70] (0.06,2.5) -- (4.54,2.5);
% Ions H3O+ (nombreux)
\foreach \x/\y in {0.7/0.6, 1.6/1.4, 2.5/0.7, 3.4/1.5, 1.0/2.1, 3.9/0.6, 2.0/2.2, 3.0/2.0}{
  \node[circle, fill=poporange, text=white, font=\scriptsize\bfseries, minimum size=0.62cm, inner sep=1pt] at (\x,\y) {\ce{H3O+}};
}
% Ions Cl- (nombreux)
\foreach \x/\y in {1.2/0.8, 2.1/1.6, 3.0/0.9, 0.6/1.6, 3.8/2.2, 1.7/0.4, 2.7/1.9, 4.1/1.2}{
  \node[circle, fill=popgreen, text=white, font=\scriptsize\bfseries, minimum size=0.52cm, inner sep=1pt] at (\x,\y) {\ce{Cl-}};
}
% Un ion HO- rare
\node[circle, fill=poppurple, text=white, font=\scriptsize\bfseries, minimum size=0.52cm, inner sep=1pt] at (0.9,1.1) {\ce{HO-}};
% Légende
\node[anchor=west, font=\small] at (5.1,2.2) {\textcolor{poporange}{\textbf{\ce{H3O+}}} : nombreux ($10^{-2}$ mol/L)};
\node[anchor=west, font=\small] at (5.1,1.6) {\textcolor{popgreen}{\textbf{\ce{Cl-}}} : nombreux ($10^{-2}$ mol/L)};
\node[anchor=west, font=\small] at (5.1,1.0) {\textcolor{poppurple}{\textbf{\ce{HO-}}} : très rares ($10^{-12}$ mol/L)};
\node[anchor=west, font=\small] at (5.1,0.4) {\textcolor{popblue}{\textbf{\ce{H2O}}} : très majoritaire};
\end{tikzpicture}
\legende{Coup d'œil microscopique dans une solution d'acide chlorhydrique à \SI{1,0e-2}{mol.L^{-1}} : pour un ion \ce{HO-}, il y a environ dix milliards d'ions \ce{H3O+} et \ce{Cl-}.}
\end{popfigure}

\courssub{Exemple 2 : solution d'hydroxyde de sodium (soude) de concentration C}

Raisonnons maintenant sur une solution de soude de concentration $C$. L'hydroxyde de sodium est un solide ionique qui se dissout totalement dans l'eau (base forte) :

\[ \ce{NaOH(s) ->[H2O] Na+(aq) + HO-(aq)} \qquad \text{(dissociation totale, exothermique)} \]

\textbf{Inventaire complet :} \ce{H2O} (majoritaire), \ce{Na+}, \ce{HO-}, \ce{H3O+}.

\textbf{Électroneutralité :} $[\ce{H3O+}] + [\ce{Na+}] = [\ce{HO-}]$.

\textbf{Conservation de la matière :} tout le sodium apporté se retrouve en ions \ce{Na+}, d'où $[\ce{Na+}] = C$.

En milieu basique ($\mathrm{pH} > 7$), c'est cette fois $[\ce{H3O+}]$ qui devient négligeable devant $[\ce{HO-}]$ : l'électroneutralité se simplifie en $[\ce{HO-}] \approx [\ce{Na+}] = C$, et par le produit ionique :
\[ [\ce{H3O+}] = \frac{K_e}{[\ce{HO-}]} \approx \frac{K_e}{C} \qquad\text{d'où}\qquad \mathrm{pH} = pK_e + \log C \]
On retrouve la formule de la base forte établie dans la section précédente. Tu vois la cohérence de l'édifice : \emph{toutes} les formules découlent de trois relations simples.

\begin{exoresolu}[Soude à $C = \SI{5,0e-3}{mol.L^{-1}}$]
Une solution d'hydroxyde de sodium a pour concentration $C = \SI{5,0e-3}{mol.L^{-1}}$. On donne $pK_e = 14$ à \SI{25}{\celsius} et $\log 5 = 0,7$.
\begin{enumerate}
\item Faire l'inventaire des espèces chimiques présentes.
\item Écrire l'équation d'électroneutralité et l'équation de conservation de la matière.
\item Calculer le pH de la solution, puis les concentrations de toutes les espèces ioniques.
\end{enumerate}
\tcblower
\textbf{1. Inventaire :} \ce{H2O} (solvant majoritaire), \ce{Na+}, \ce{HO-} (ions du soluté), \ce{H3O+} (ion de l'eau, toujours présent).

\textbf{2. Équations :}
\begin{itemize}
\item Électroneutralité : $[\ce{Na+}] + [\ce{H3O+}] = [\ce{HO-}]$ ;
\item Conservation : $[\ce{Na+}] = C = \SI{5,0e-3}{mol.L^{-1}}$.
\end{itemize}

\textbf{3. Calculs :} la soude est une base forte, donc
\[ \mathrm{pH} = pK_e + \log C = 14 + \log(5{,}0\times 10^{-3}) = 14 + 0{,}7 - 3 = 11{,}7. \]
Alors :
\begin{align*}
[\ce{H3O+}] &= 10^{-\mathrm{pH}} = 10^{-11,7} = \SI{2,0e-12}{mol.L^{-1}} \\
[\ce{HO-}] &= 10^{\mathrm{pH}-14} = 10^{-2,3} = \SI{5,0e-3}{mol.L^{-1}} \\
[\ce{Na+}] &= C = \SI{5,0e-3}{mol.L^{-1}}
\end{align*}
\textbf{Vérification :} $[\ce{Na+}] + [\ce{H3O+}] = 5{,}0\times 10^{-3} + 2{,}0\times 10^{-12} \approx 5{,}0\times 10^{-3} = [\ce{HO-}]$. L'électroneutralité est respectée. On constate que $[\ce{H3O+}] \ll [\ce{Na+}]$ : en milieu basique, les ions \ce{H3O+} sont négligeables, mais ils existent.
\end{exoresolu}

\courssub{Généralisation : la méthode systématique en quatre étapes}

Quelle que soit la solution d'acide fort ou de base forte, la démarche est toujours la même. Apprends-la comme une chorégraphie : elle te servira dans toute la suite du programme (dosages, mélanges, acides faibles).

\begin{methode}[Déterminer toutes les concentrations d'une solution]
\begin{enumerate}
\item \textbf{Inventaire} : écrire la réaction avec l'eau, puis lister toutes les espèces présentes : \ce{H2O}, les ions du soluté, \ce{H3O+} \emph{et} \ce{HO-} ;
\item \textbf{Électroneutralité} : écrire que la somme des charges positives égale la somme des charges négatives ;
\item \textbf{Conservation de la matière} : relier la concentration de chaque ion du soluté à la concentration apportée $C$ (réaction totale) ;
\item \textbf{Calcul via le pH} : $[\ce{H3O+}] = 10^{-\mathrm{pH}}$ ; $[\ce{HO-}] = 10^{\mathrm{pH}-14}$ ; puis conclure sur toutes les concentrations et vérifier l'électroneutralité.
\end{enumerate}
\end{methode}

\begin{popfigure}
\begin{tikzpicture}[
  node distance=0.55cm,
  etape/.style={rectangle, rounded corners=4pt, draw=popblue, fill=popblueL, text width=10.2cm, align=left, font=\small, inner sep=6pt},
  num/.style={circle, fill=poporange, text=white, font=\bfseries, minimum size=0.55cm, inner sep=1pt},
  fleche/.style={-{Stealth[length=2.5mm]}, thick, poporange}
]
\node[etape] (e1) {\textbf{Inventaire} : réaction avec l'eau $\rightarrow$ liste des espèces (\ce{H2O}, ions du soluté, \ce{H3O+}, \ce{HO-})};
\node[num, left=0.15cm of e1] {1};
\node[etape, below=of e1] (e2) {\textbf{Électroneutralité} : $\sum$ charges positives $=$ $\sum$ charges négatives};
\node[num, left=0.15cm of e2] {2};
\node[etape, below=of e2] (e3) {\textbf{Conservation de la matière} : $[\text{ion du soluté}] = C$ (réaction totale)};
\node[num, left=0.15cm of e3] {3};
\node[etape, below=of e3] (e4) {\textbf{Calcul via le pH} : $[\ce{H3O+}] = 10^{-\mathrm{pH}}$ ; $[\ce{HO-}] = 10^{\mathrm{pH}-14}$ ; vérification finale};
\node[num, left=0.15cm of e4] {4};
\draw[fleche] (e1.south) -- (e2.north);
\draw[fleche] (e2.south) -- (e3.north);
\draw[fleche] (e3.south) -- (e4.north);
\end{tikzpicture}
\legende{Organigramme de la méthode systématique en quatre étapes pour déterminer toutes les concentrations d'une solution aqueuse.}
\end{popfigure}

\begin{savaistu}[Sørensen et la naissance du pH]
C'est le chimiste danois \textbf{Søren Peter Lauritz Sørensen} (1868--1939), chef du laboratoire de la brasserie Carlsberg, qui a introduit la notion de pH en 1909. Il cherchait à contrôler l'acidité des bières pour en stabiliser la qualité. Il a défini $\mathrm{pH} = -\log_{10}[\ce{H+}]$ (notation de l'époque, aujourd'hui $[\ce{H3O+}]$). Le « p » vient de l'allemand \emph{Potenz} (puissance) ou du français \emph{puissance} ; le « H » pour l'hydrogène. Au Cameroun, les brasseries (comme les Brasseries du Cameroun) utilisent encore cette rigueur analytique pour garantir la constance de leurs produits, de la bière au jus de bissap.
\end{savaistu}

\begin{aretenir}[L'essentiel de la section]
\begin{itemize}
\item Toute solution aqueuse contient \textbf{toujours} \ce{H2O}, \ce{H3O+} et \ce{HO-}, plus les ions du soluté ;
\item Deux relations universelles : l'\cle{électroneutralité} (neutralité électrique de la solution) et la \cle{conservation de la matière} (l'élément apporté se retrouve intégralement) ;
\item À partir du pH : $[\ce{H3O+}] = 10^{-\mathrm{pH}}$ et $[\ce{HO-}] = 10^{\mathrm{pH}-14}$ à \SI{25}{\celsius} ;
\item En milieu acide, $[\ce{HO-}]$ est négligeable ; en milieu basique, c'est $[\ce{H3O+}]$ qui l'est — mais \emph{négligeable ne veut pas dire absent} ;
\item Acide fort : $[\ce{H3O+}] \approx C$ ; base forte : $[\ce{HO-}] \approx C$.
\end{itemize}
\end{aretenir}

\begin{exercice}[À toi de jouer]
Une solution d'acide nitrique \ce{HNO3} (acide fort) a un pH égal à $3,0$ à \SI{25}{\celsius}.
\begin{enumerate}
\item Faire l'inventaire des espèces chimiques présentes.
\item Calculer $[\ce{H3O+}]$ et $[\ce{HO-}]$.
\item En déduire $[\ce{NO3-}]$ et la concentration $C$ de la solution.
\end{enumerate}
\end{exercice}

\begin{corrige}[Exercice : solution d'acide nitrique]
\textbf{1. Inventaire :} \ce{HNO3 + H2O -> H3O+ + NO3-} (totale). Espèces : \ce{H2O}, \ce{H3O+}, \ce{NO3-}, \ce{HO-}.

\textbf{2.} $[\ce{H3O+}] = 10^{-3} = \SI{1,0e-3}{mol.L^{-1}}$ ; $[\ce{HO-}] = 10^{3-14} = \SI{1,0e-11}{mol.L^{-1}}$.

\textbf{3.} Électroneutralité : $[\ce{NO3-}] = [\ce{H3O+}] - [\ce{HO-}] \approx \SI{1,0e-3}{mol.L^{-1}}$. Conservation : $C = [\ce{NO3-}] = \SI{1,0e-3}{mol.L^{-1}}$.
\end{corrige}

Maintenant que tu sais compter les espèces dans une solution, une question naturelle se pose : que deviennent ces concentrations — et donc le pH — lorsqu'on ajoute de l'eau ? C'est l'objet de la section suivante, consacrée à l'effet de la dilution.

\coursec{Effet de la dilution sur le pH}

Dans la section précédente, nous avons appris à faire l'inventaire complet des espèces présentes dans une solution d'acide fort ou de base forte et à calculer leurs concentrations grâce à l'électroneutralité et à la conservation de la matière. Nous savons donc désormais que, pour une solution d'acide chlorhydrique de concentration $C$, on a $[\ce{H3O+}] = C$ et donc $\mathrm{pH} = -\log C$. Mais que devient ce pH lorsqu'on ajoute de l'eau distillée à la solution ? C'est la question que se pose tout technicien de laboratoire, par exemple à la SONARA ou dans une brasserie, lorsqu'il doit préparer une solution étalon de pH précis à partir d'une solution concentrée. Cette section y répond de manière quantitative.

\courssub{Rappels : la dilution et le facteur de dilution}

\textbf{Intuition.} Diluer, c'est ajouter du solvant (ici de l'eau distillée) à une solution sans ajouter ni retirer de soluté. Le nombre de moles de soluté ne change pas : il est simplement « réparti » dans un volume plus grand. La concentration diminue donc.

\begin{definition}[Dilution et facteur de dilution]
Diluer une solution \emph{mère} de concentration $C_i$, c'est prélever un volume $V_i$ de cette solution et compléter avec de l'eau distillée jusqu'à un volume final $V_f$ : on obtient une solution \emph{fille} de concentration $C_f$. La conservation de la quantité de matière de soluté s'écrit :
\[ C_i \times V_i = C_f \times V_f \]
On définit le \cle{facteur de dilution} :
\[ F = \dfrac{C_i}{C_f} = \dfrac{V_f}{V_i} \]
\end{definition}

\begin{exemplebox}[Dilution au dixième]
On dit qu'on réalise une \cle{dilution au dixième} lorsque $F = 10$ : la concentration est divisée par $10$. Concrètement, on prélève \SI{10,0}{mL} de solution mère à la pipette jaugée et on complète à \SI{100,0}{mL} avec de l'eau distillée dans une fiole jaugée :
\[ C_f = \dfrac{C_i \times V_i}{V_f} = \dfrac{C_i \times 10{,}0}{100{,}0} = \dfrac{C_i}{10} \]
\end{exemplebox}

\begin{attention}[Précision de la verrerie]
Une dilution destinée à fixer un pH précis exige de la verrerie de précision : pipette jaugée (et non éprouvette) pour le prélèvement, fiole jaugée (et non bécher) pour l'ajustage. Le trait de jauge se règle à l'œil, le bas du ménisque affleurant le trait. Une erreur de volume se traduit directement par une erreur de concentration, donc de pH.
\end{attention}

\courssub{Expérience : dilutions successives au dixième d'une solution d'acide chlorhydrique}

\begin{experience}[Dilutions en cascade et mesures de pH]
\textbf{Dispositif et protocole.} On dispose d'une solution mère $S_0$ d'acide chlorhydrique de concentration $C_0 = \SI{1,0e-1}{mol.L^{-1}}$. On prépare par dilutions successives au dixième les solutions $S_1$, $S_2$, $S_3$, $S_4$ : à chaque étape, on prélève \SI{10,0}{mL} de la solution précédente et on complète à \SI{100,0}{mL} dans une fiole jaugée. On mesure le pH de chaque solution avec un pH-mètre étalonné.

\textbf{Observations.} Les résultats sont consignés dans le tableau ci-dessous :

\begin{center}
\begin{tabularx}{\linewidth}{|l|*{5}{>{\centering\arraybackslash}X|}}
\hline
\rowcolor{popblueL} Solution & $S_0$ & $S_1$ & $S_2$ & $S_3$ & $S_4$ \\
\hline
$C$ (\si{mol.L^{-1}}) & $10^{-1}$ & $10^{-2}$ & $10^{-3}$ & $10^{-4}$ & $10^{-5}$ \\
\hline
pH mesuré & $1{,}0$ & $2{,}0$ & $3{,}0$ & $4{,}0$ & $5{,}0$ \\
\hline
\end{tabularx}
\end{center}

\textbf{Interprétation.} À chaque dilution au dixième, la concentration en ions \ce{H3O+} est divisée par $10$ et le pH augmente exactement d'\emph{une unité}.
\end{experience}

\begin{popfigure}
\begin{center}
\begin{tikzpicture}[scale=0.92]
% Fiole S0
\begin{scope}[shift={(0,0)}]
\draw[popdark,thick] (-0.55,0) -- (-0.55,0.9) -- (-0.12,1.5) -- (-0.12,2.1);
\draw[popdark,thick] (0.55,0) -- (0.55,0.9) -- (0.12,1.5) -- (0.12,2.1);
\draw[popdark,thick] (-0.55,0) arc (180:360:0.55);
\draw[popdark] (-0.16,1.9) -- (0.16,1.9);
\fill[poporange!70] (-0.55,0) arc (180:360:0.55) -- (0.55,0.55) -- (-0.55,0.55) -- cycle;
\node[below,font=\small\bfseries] at (0,-0.35) {$S_0$};
\node[below,font=\footnotesize] at (0,-0.75) {$C_0=10^{-1}$};
\node[below,font=\footnotesize\bfseries,poporange!60!black] at (0,-1.12) {pH $=1$};
\end{scope}
% flèche 1
\draw[-{Stealth},popblue,very thick] (1.0,1.0) -- (2.0,1.0) node[midway,above,font=\footnotesize] {$\div 10$};
% Fiole S1
\begin{scope}[shift={(3.0,0)}]
\draw[popdark,thick] (-0.55,0) -- (-0.55,0.9) -- (-0.12,1.5) -- (-0.12,2.1);
\draw[popdark,thick] (0.55,0) -- (0.55,0.9) -- (0.12,1.5) -- (0.12,2.1);
\draw[popdark,thick] (-0.55,0) arc (180:360:0.55);
\draw[popdark] (-0.16,1.9) -- (0.16,1.9);
\fill[poporange!50] (-0.55,0) arc (180:360:0.55) -- (0.55,0.55) -- (-0.55,0.55) -- cycle;
\node[below,font=\small\bfseries] at (0,-0.35) {$S_1$};
\node[below,font=\footnotesize] at (0,-0.75) {$C_1=10^{-2}$};
\node[below,font=\footnotesize\bfseries,poporange!60!black] at (0,-1.12) {pH $=2$};
\end{scope}
\draw[-{Stealth},popblue,very thick] (4.0,1.0) -- (5.0,1.0) node[midway,above,font=\footnotesize] {$\div 10$};
% Fiole S2
\begin{scope}[shift={(6.0,0)}]
\draw[popdark,thick] (-0.55,0) -- (-0.55,0.9) -- (-0.12,1.5) -- (-0.12,2.1);
\draw[popdark,thick] (0.55,0) -- (0.55,0.9) -- (0.12,1.5) -- (0.12,2.1);
\draw[popdark,thick] (-0.55,0) arc (180:360:0.55);
\draw[popdark] (-0.16,1.9) -- (0.16,1.9);
\fill[poporange!30] (-0.55,0) arc (180:360:0.55) -- (0.55,0.55) -- (-0.55,0.55) -- cycle;
\node[below,font=\small\bfseries] at (0,-0.35) {$S_2$};
\node[below,font=\footnotesize] at (0,-0.75) {$C_2=10^{-3}$};
\node[below,font=\footnotesize\bfseries,poporange!60!black] at (0,-1.12) {pH $=3$};
\end{scope}
\draw[-{Stealth},popblue,very thick] (7.0,1.0) -- (8.0,1.0) node[midway,above,font=\footnotesize] {$\div 10$};
% Fiole S3
\begin{scope}[shift={(9.0,0)}]
\draw[popdark,thick] (-0.55,0) -- (-0.55,0.9) -- (-0.12,1.5) -- (-0.12,2.1);
\draw[popdark,thick] (0.55,0) -- (0.55,0.9) -- (0.12,1.5) -- (0.12,2.1);
\draw[popdark,thick] (-0.55,0) arc (180:360:0.55);
\draw[popdark] (-0.16,1.9) -- (0.16,1.9);
\fill[poporange!12] (-0.55,0) arc (180:360:0.55) -- (0.55,0.55) -- (-0.55,0.55) -- cycle;
\node[below,font=\small\bfseries] at (0,-0.35) {$S_3$};
\node[below,font=\footnotesize] at (0,-0.75) {$C_3=10^{-4}$};
\node[below,font=\footnotesize\bfseries,poporange!60!black] at (0,-1.12) {pH $=4$};
\end{scope}
% pointillés vers S4
\draw[-{Stealth},popblue,very thick,dashed] (10.0,1.0) -- (11.0,1.0) node[midway,above,font=\footnotesize] {$\div 10$};
\node[font=\footnotesize,align=center] at (11.9,1.0) {$S_4$ : pH $=5$\\ $S_5$ : pH $=6$};
\end{tikzpicture}
\end{center}
\legende{Dilutions successives au dixième d'une solution d'acide chlorhydrique : la concentration est divisée par 10 à chaque étape et le pH augmente d'une unité. La teinte de la solution symbolise la concentration décroissante en ions \ce{H3O+}.}
\end{popfigure}

\courssub{Justification théorique : pourquoi le pH varie-t-il d'une unité ?}

\textbf{Intuition.} Le pH est un logarithme. Or la fonction logarithme possède une propriété remarquable : multiplier un nombre par $10$ revient à ajouter $1$ à son logarithme. Diviser la concentration par $10$ revient donc à ajouter $1$ au pH. Tout est là.

\textbf{Démonstration rigoureuse.} Pour une solution d'acide fort de concentration $C$ (avec $C > \SI{e-6}{mol.L^{-1}}$), la réaction avec l'eau étant totale :
\[ \ce{HCl + H2O -> H3O+ + Cl-} \quad\text{(réaction totale, rapide)} \]
on a $[\ce{H3O+}] = C$, donc $\mathrm{pH} = -\log C$. Après une dilution au dixième, la nouvelle concentration est $C' = \dfrac{C}{10}$, d'où :
\begin{align*}
\mathrm{pH}' &= -\log C' = -\log\left(\dfrac{C}{10}\right) \\
&= -\log C + \log 10 \\
&= -\log C + 1 = \mathrm{pH} + 1
\end{align*}

\begin{propriete}[Dilution d'un acide fort]
Toute dilution au dixième d'une solution d'acide fort (dans le domaine de validité $C > \SI{e-6}{mol.L^{-1}}$) \cle{augmente le pH d'une unité} : $\Delta\mathrm{pH} = +1$. Plus généralement, une dilution d'un facteur $F = 10^{n}$ augmente le pH de $n$ unités :
\[ \mathrm{pH}_f = \mathrm{pH}_i + \log F \]
\end{propriete}

\textbf{Cas de la base forte.} Pour une solution d'hydroxyde de sodium de concentration $C$, la dissociation est totale :
\[ \ce{NaOH ->[eau] Na+ + HO-} \quad\text{(dissociation totale)} \]
On a $[\ce{HO-}] = C$ et $\mathrm{pH} = \mathrm{p}K_e + \log C$ (avec $\mathrm{p}K_e = 14$ à \SI{25}{\celsius}). Après dilution au dixième :
\begin{align*}
\mathrm{pH}' &= \mathrm{p}K_e + \log\left(\dfrac{C}{10}\right) = \mathrm{p}K_e + \log C - \log 10 \\
&= \mathrm{pH} - 1
\end{align*}

\begin{propriete}[Dilution d'une base forte]
Toute dilution au dixième d'une solution de base forte (dans le domaine de validité) \cle{diminue le pH d'une unité} : $\Delta\mathrm{pH} = -1$. Une dilution d'un facteur $F = 10^{n}$ diminue le pH de $n$ unités :
\[ \mathrm{pH}_f = \mathrm{pH}_i - \log F \]
\end{propriete}

\begin{methode}[Prévoir le pH après dilution]
\begin{enumerate}
\item Identifier la nature de la solution : acide fort ou base forte.
\item Calculer le facteur de dilution $F = \dfrac{C_i}{C_f} = \dfrac{V_f}{V_i}$.
\item Écrire $F$ sous la forme $10^{n}$ (ou calculer $\log F$).
\item Appliquer la règle : acide fort $\Rightarrow \mathrm{pH}_f = \mathrm{pH}_i + \log F$ ; base forte $\Rightarrow \mathrm{pH}_f = \mathrm{pH}_i - \log F$.
\item Vérifier la cohérence : le pH d'un acide dilué doit rester inférieur à 7, celui d'une base diluée supérieur à 7.
\end{enumerate}
\end{methode}

\begin{exoresolu}[Dilution cent fois d'une solution d'acide chlorhydrique]
On dispose d'une solution d'acide chlorhydrique de pH $= 2{,}0$. On prélève \SI{5,0}{mL} de cette solution et on complète à \SI{500,0}{mL} avec de l'eau distillée dans une fiole jaugée. Quel est le pH de la solution obtenue ?
\tcblower
\textbf{Étape 1 : nature de la solution.} L'acide chlorhydrique est un acide fort, totalement ionisé dans l'eau.

\textbf{Étape 2 : facteur de dilution.}
\[ F = \dfrac{V_f}{V_i} = \dfrac{500{,}0}{5{,}0} = 100 = 10^{2} \]

\textbf{Étape 3 : application de la règle.} Diluer $100$ fois, c'est diluer deux fois de suite au dixième : le pH augmente de $2$ unités.
\[ \mathrm{pH}_f = \mathrm{pH}_i + \log F = 2{,}0 + \log 100 = 2{,}0 + 2 = 4{,}0 \]

\textbf{Vérification par les concentrations.} La solution initiale a $C_i = 10^{-\mathrm{pH}} = \SI{1,0e-2}{mol.L^{-1}}$. Après dilution : $C_f = \dfrac{C_i}{100} = \SI{1,0e-4}{mol.L^{-1}}$, d'où $\mathrm{pH}_f = -\log(10^{-4}) = 4{,}0$. Le résultat est confirmé, et $C_f > \SI{e-6}{mol.L^{-1}}$ : la formule est bien applicable.
\end{exoresolu}

\courssub{Limites de validité : le rôle de l'autoprotolyse de l'eau}

\textbf{Le paradoxe apparent.} Appliquons aveuglément la formule $\mathrm{pH} = -\log C$ à une solution d'acide chlorhydrique de concentration $C = \SI{e-8}{mol.L^{-1}}$ : nous obtiendrions pH $= 8$ ! Une solution \emph{acide} aurait un pH \emph{basique}, ce qui est absurde. Où est l'erreur ?

\textbf{L'origine du problème.} Jusqu'ici, nous avons considéré que les ions \ce{H3O+} provenaient exclusivement de l'acide. Or l'eau pure elle-même produit des ions \ce{H3O+} et \ce{HO-} par sa réaction d'\cle{autoprotolyse} :
\[ \ce{2 H2O <=> H3O+ + HO-} \quad\text{(réaction très limitée, } K_e = 10^{-14} \text{ à } \SI{25}{\celsius}\text{)} \]
Dans l'eau pure, cette réaction impose $[\ce{H3O+}] = [\ce{HO-}] = \SI{e-7}{mol.L^{-1}}$. Tant que l'acide apporte une concentration $C \gg \SI{e-7}{mol.L^{-1}}$, cet apport de l'eau est négligeable. Mais lorsque $C$ devient comparable ou inférieur à $\SI{e-7}{mol.L^{-1}}$, les ions \ce{H3O+} de l'autoprotolyse ne sont plus négligeables : ils deviennent même majoritaires.

\textbf{Conséquence.} En diluant indéfiniment un acide fort, la concentration totale en \ce{H3O+} ne descend jamais en dessous de celle de l'eau pure : le pH \cle{tend vers 7 sans jamais l'atteindre ni le dépasser}. Symétriquement, en diluant indéfiniment une base forte, le pH tend vers 7 par valeurs supérieures sans jamais descendre en dessous.

\begin{attention}[L'erreur classique à ne jamais commettre]
\textbf{Erreur fréquente :} croire qu'en diluant suffisamment un acide, son pH devient supérieur à 7. C'est \emph{impossible} ! Une dilution ne fait qu'ajouter de l'eau : elle ne peut pas transformer un acide en base. Le pH d'une solution d'acide fort, aussi diluée soit-elle, reste toujours \emph{légèrement inférieur} à 7 et tend vers 7. De même, le pH d'une base forte diluée reste toujours légèrement supérieur à 7. Au Baccalauréat, tout résultat de calcul donnant pH $> 7$ pour une solution d'acide (ou pH $< 7$ pour une base) doit immédiatement vous alerter : la formule simple a été appliquée hors de son domaine de validité.
\end{attention}

\begin{aretenir}[Domaine de validité des formules]
Les formules $\mathrm{pH} = -\log C$ (acide fort) et $\mathrm{pH} = \mathrm{p}K_e + \log C$ (base forte) ne sont valables que si la concentration est suffisamment grande devant $10^{-7}$ \si{mol.L^{-1}}, en pratique :
\[ C > \SI{e-6}{mol.L^{-1}} \]
En dessous de cette limite, l'autoprotolyse de l'eau n'est plus négligeable et le pH tend vers 7.
\end{aretenir}

\begin{popfigure}
\begin{center}
\begin{tikzpicture}
\begin{axis}[
width=0.98\linewidth, height=7.2cm,
xlabel={$-\log C$},
ylabel={pH},
xmin=0, xmax=9.2, ymin=4, ymax=10.5,
xtick={0,1,2,3,4,5,6,7,8,9},
ytick={4,5,6,7,8,9,10},
grid=major, grid style={popline!40},
legend style={at={(0.03,0.55)},anchor=west,font=\footnotesize,draw=popline},
axis line style={popdark},
]
% Zone de validité
\fill[popgreen!12] (axis cs:0,4) rectangle (axis cs:6,10.5);
% Zone de non validité
\fill[poporange!14] (axis cs:6,4) rectangle (axis cs:9.2,10.5);
% Asymptote pH = 7
\addplot[popmuted,dashed,thick,domain=0:9.2] {7};
\node[font=\footnotesize,popmuted,anchor=south west] at (axis cs:0.15,7.02) {asymptote pH $=7$};
% Acide fort : pH = -log C, puis courbure vers 7
\addplot[poporange,very thick,domain=0:6,samples=100] {x};
\addplot[poporange,very thick,domain=6:9,samples=200] {-log10(10^(-x)+1e-7)};
\addlegendentry{acide fort : pH $= -\log C$}
% Base forte : pH = 14 + log C, puis courbure vers 7
\addplot[popblue,very thick,domain=0:6,samples=100] {14-x};
\addplot[popblue,very thick,domain=6:9,samples=200] {14+log10(10^(-x)+1e-7)};
\addlegendentry{base forte : pH $= \mathrm{p}K_e + \log C$}
% Délimitations
\draw[popgreen!60!black,dashed] (axis cs:6,4) -- (axis cs:6,10.5);
\node[font=\footnotesize,popgreen!50!black,align=center] at (axis cs:3,10.1) {zone de validité\\ $C > 10^{-6}$ \si{mol.L^{-1}}};
\node[font=\footnotesize,poporange!60!black,align=center] at (axis cs:7.7,10.1) {autoprotolyse\\ de l'eau dominante};
\end{axis}
\end{tikzpicture}
\end{center}
\legende{pH en fonction de $-\log C$ pour un acide fort (orange) et une base forte (bleu) à \SI{25}{\celsius}. Dans la zone verte ($C > \SI{e-6}{mol.L^{-1}}$), les courbes sont des droites de pentes $+1$ et $-1$. Aux très faibles concentrations, les deux courbes s'infléchissent et tendent vers l'asymptote pH $= 7$ à cause de l'autoprotolyse de l'eau.}
\end{popfigure}

\begin{savaistu}[Un acide dilué dix millions de fois reste acide]
Si l'on dilue $10^{7}$ fois une solution d'acide chlorhydrique à \SI{0,1}{mol.L^{-1}}, la concentration apportée par l'acide n'est plus que $C = \SI{e-8}{mol.L^{-1}}$. La formule naïve donnerait pH $= 8$... alors que le pH réel mesuré est voisin de $6{,}8$ ! En effet, les ions \ce{H3O+} issus de l'autoprotolyse de l'eau ($\approx \SI{e-7}{mol.L^{-1}}$) sont dix fois plus nombreux que ceux apportés par l'acide. Le calcul exact, qui combine électroneutralité et produit ionique de l'eau, donne $[\ce{H3O+}] \approx \SI{1,6e-7}{mol.L^{-1}}$, soit pH $\approx 6{,}8$ : la solution reste très légèrement acide. Ce phénomène n'est pas qu'une curiosité de laboratoire : il explique pourquoi l'eau de pluie, même non polluée, a un pH voisin de $5{,}6$ (dissolution du \ce{CO2} atmosphérique), et pourquoi les mesures de pH très précises exigent une eau ultra-pure, débarrassée de tout gaz dissous.
\end{savaistu}

\courssub{Application : préparation d'une solution de pH donné par dilution}

\textbf{Problème type.} Au laboratoire, on dispose d'une solution commerciale concentrée et l'on souhaite préparer un volume donné de solution de pH précis. La démarche inverse les calculs précédents : on part du pH visé pour remonter au facteur de dilution, puis aux volumes à utiliser.

\begin{methode}[Préparer une solution de pH imposé]
\begin{enumerate}
\item À partir du pH visé, calculer la concentration fille : $C_f = 10^{-\mathrm{pH}}$ pour un acide fort, ou $C_f = 10^{\mathrm{pH} - \mathrm{p}K_e}$ pour une base forte.
\item Calculer le facteur de dilution $F = \dfrac{C_i}{C_f}$.
\item Choisir le volume final $V_f$ (fiole jaugée disponible) et en déduire le volume à prélever : $V_i = \dfrac{V_f}{F}$.
\item Prélever $V_i$ à la pipette jaugée, introduire dans la fiole jaugée de volume $V_f$, compléter à moitié avec de l'eau distillée, agiter, puis ajuster au trait de jauge et homogénéiser.
\item Contrôler le pH obtenu au pH-mètre étalonné.
\end{enumerate}
\end{methode}

\begin{exoresolu}[Préparation de 200 mL de solution de pH = 3]
Un laborantin du lycée dispose d'une solution mère d'acide chlorhydrique de concentration $C_0 = \SI{1,0e-1}{mol.L^{-1}}$. Il veut préparer \SI{200,0}{mL} de solution de pH $= 3{,}0$. Quel volume de solution mère doit-il prélever ?
\tcblower
\textbf{Étape 1 : concentration de la solution fille.} L'acide chlorhydrique est un acide fort :
\[ C_f = 10^{-\mathrm{pH}} = 10^{-3} = \SI{1,0e-3}{mol.L^{-1}} \]
Comme $C_f > \SI{e-6}{mol.L^{-1}}$, la formule est applicable.

\textbf{Étape 2 : facteur de dilution.}
\[ F = \dfrac{C_0}{C_f} = \dfrac{10^{-1}}{10^{-3}} = 100 \]

\textbf{Étape 3 : volume à prélever.}
\[ V_i = \dfrac{V_f}{F} = \dfrac{200{,}0}{100} = \SI{2,0}{mL} \]

\textbf{Étape 4 : mode opératoire.} On prélève \SI{2,0}{mL} de solution mère à la pipette jaugée, on les introduit dans une fiole jaugée de \SI{200,0}{mL} contenant déjà un peu d'eau distillée, on complète jusqu'au trait de jauge et on homogénéise. Le pH de la solution obtenue est $3{,}0$.

\textbf{Remarque de sécurité :} on verse toujours l'acide dans l'eau, jamais l'inverse, et la solution mère concentrée se manipule avec gants et lunettes.
\end{exoresolu}

\begin{exercice}[Pour s'entraîner]
Une solution d'hydroxyde de sodium a un pH égal à $12{,}0$ à \SI{25}{\celsius}. On prélève \SI{20,0}{mL} de cette solution et on complète à \SI{2,0}{L} avec de l'eau distillée. Déterminer le pH de la solution obtenue. La formule utilisée est-elle dans son domaine de validité ?
\end{exercice}

\begin{corrige}[Exercice 1]
La solution est une base forte : $\mathrm{pH}_i = \mathrm{p}K_e + \log C_i$, donc $C_i = 10^{\mathrm{pH}_i - 14} = 10^{-2}$ \si{mol.L^{-1}}.

Facteur de dilution : $F = \dfrac{V_f}{V_i} = \dfrac{2000}{20{,}0} = 100 = 10^{2}$.

Pour une base forte, le pH diminue de $\log F = 2$ unités :
\[ \mathrm{pH}_f = 12{,}0 - 2 = 10{,}0 \]

Vérification : $C_f = \dfrac{C_i}{100} = \SI{e-4}{mol.L^{-1}} > \SI{e-6}{mol.L^{-1}}$ : la formule est bien dans son domaine de validité. De plus, $\mathrm{pH}_f = 10 > 7$ : la solution reste basique, ce qui est cohérent.
\end{corrige}

\begin{aretenir}[L'essentiel de la section]
\begin{itemize}
\item Diluer au dixième ($F = 10$) une solution d'\textbf{acide fort} : le pH \textbf{augmente de 1}. Diluer au dixième une solution de \textbf{base forte} : le pH \textbf{diminue de 1}.
\item En général : $\mathrm{pH}_f = \mathrm{pH}_i + \log F$ (acide fort) et $\mathrm{pH}_f = \mathrm{pH}_i - \log F$ (base forte).
\item Ces règles ne sont valables que pour $C > \SI{e-6}{mol.L^{-1}}$.
\item Aux très faibles concentrations, l'autoprotolyse de l'eau $\ce{2 H2O <=> H3O+ + HO-}$ devient prépondérante : le pH de toute solution aqueuse diluée \textbf{tend vers 7} sans jamais franchir cette valeur.
\end{itemize}
\end{aretenir}

\coursec{Mélange d'un acide fort et d'une base forte}

Dans la section précédente, nous avons vu que diluer un acide fort ou une base forte modifie le pH de manière prévisible : une dilution au dixième fait varier le pH d'une unité. Mais que se passe-t-il lorsqu'on mélange, non pas une solution avec de l'eau pure, mais un \cle{acide fort} avec une \cle{base forte} ? C'est exactement la situation que vit le laborantin d'une savonnerie artisanale de Foumban qui vérifie la concentration de sa lessive de potasse avec un acide titré, ou le technicien de la SONARA qui neutralise des effluents acides avant rejet. Ce mélange fait intervenir l'une des réactions les plus importantes de toute la chimie des solutions : la réaction entre les ions \ce{H3O+} et \ce{HO-}.

\courssub{La réaction entre les ions \ce{H3O+} et \ce{HO-}}

\textbf{Intuition.} Une solution d'acide fort contient des ions \ce{H3O+} ; une solution de base forte contient des ions \ce{HO-}. Ces deux ions sont « incompatibles » : dès qu'ils se rencontrent, ils réagissent pour former de l'eau. L'ion \ce{H3O+} cède son proton \ce{H+} à l'ion \ce{HO-}, qui le capte avidement. C'est une \emph{réaction acido-basique}, c'est-à-dire un transfert de proton.

\begin{definition}[Réaction entre un acide fort et une base forte]
Lorsqu'on mélange une solution d'acide fort et une solution de base forte, il se produit une réaction entre les ions oxonium et les ions hydroxyde, d'équation :
\[ \ce{H3O+ + HO- -> 2 H2O} \]
Cette réaction est \textbf{instantanée}, \textbf{exothermique} (elle dégage de la chaleur, sensible au toucher sur le bécher) et \textbf{quasi totale} : sa constante d'équilibre vaut
\[ K = \frac{1}{K_e} = \frac{1}{10^{-14}} = 10^{14} \quad \text{à } \SI{25}{\celsius}. \]
Une constante aussi immense signifie que la réaction se poursuit jusqu'à \emph{épuisement total du réactif en défaut}.
\end{definition}

\begin{attention}[Ne pas confondre les équations]
L'équation \ce{H3O+ + HO- -> 2 H2O} est l'équation \emph{nette ionique} : elle ne fait apparaître que les espèces qui réagissent réellement. Les ions \ce{Cl-} (apportés par l'acide chlorhydrique) et \ce{Na+} (apportés par la soude) sont des \textbf{ions spectateurs} : ils restent en solution, inchangés. On écrit parfois l'équation globale \ce{HCl + NaOH -> NaCl + H2O}, mais elle masque la réalité microscopique : \ce{HCl} et \ce{NaOH} n'existent pas en tant que tels dans l'eau, ils sont totalement dissociés. Au Baccalauréat, privilégie toujours l'écriture ionique.
\end{attention}

\courssub{Bilan des quantités de matière avant mélange}

Considérons le cas général : on verse un volume $V_a$ d'une solution d'acide fort de concentration $C_a$ dans un volume $V_b$ d'une solution de base forte de concentration $C_b$. Le volume total du mélange est $V_a + V_b$ (les volumes sont additifs, hypothèse toujours valable pour des solutions aqueuses diluées).

Avant toute réaction, les quantités de matière apportées sont :
\[ n(\ce{H3O+}) = C_a \times V_a \qquad \text{et} \qquad n(\ce{HO-}) = C_b \times V_b. \]

\begin{attention}[Unités : le piège des volumes]
Dans les formules $n = C \times V$, si $C$ est en \si{mol.L^{-1}}, alors $V$ doit être en \textbf{litres}. Or les volumes des énoncés sont presque toujours donnés en millilitres. Convertis systématiquement : $\SI{20}{mL} = \SI{0,020}{L} = \SI{2,0e-2}{L}$. Une erreur de conversion fausse tout l'exercice.
\end{attention}

\courssub{Le tableau d'avancement : l'outil systématique}

Face à un mélange acide fort + base forte, le réflexe du bon élève est de dresser un \cle{tableau d'avancement}. Comme la réaction est totale, l'avancement final $x_f$ est imposé par le réactif en défaut.

\begin{popfigure}
\begin{center}
\begin{tabular}{|l|c|c|c|c|}
\hline
\rowcolor{popblueL}
\textbf{Équation} & \ce{H3O+} & $+$ \ \ce{HO-} & \ce{->} & \ce{2 H2O} \\
\hline
État initial ($x = 0$) & $C_a V_a$ & $C_b V_b$ & & solvant (excès) \\
\hline
État intermédiaire ($x$) & $C_a V_a - x$ & $C_b V_b - x$ & & $2x$ \\
\hline
État final ($x = x_f$) & $C_a V_a - x_f$ & $C_b V_b - x_f$ & & $2x_f$ \\
\hline
\end{tabular}
\end{center}
\legende{Tableau d'avancement (quantités de matière en mol) de la réaction \ce{H3O+ + HO- -> 2 H2O}. Le réactif en défaut impose $x_f$ : sa quantité finale est nulle.}
\end{popfigure}

Trois cas se présentent alors, selon la comparaison de $C_a V_a$ et $C_b V_b$.

\textbf{Cas 1 : l'acide est en excès} ($C_a V_a > C_b V_b$).

Tous les ions \ce{HO-} sont consommés : $x_f = C_b V_b$. Il reste des ions \ce{H3O+} en quantité $C_a V_a - C_b V_b$, dissous dans le volume total $V_a + V_b$. D'où :
\[ [\ce{H3O+}] = \frac{C_a V_a - C_b V_b}{V_a + V_b} \qquad \Longrightarrow \qquad \mathrm{pH} = -\log[\ce{H3O+}] \]
Le mélange final est \textbf{acide} ($\mathrm{pH} < 7$).

\textbf{Cas 2 : la base est en excès} ($C_b V_b > C_a V_a$).

Tous les ions \ce{H3O+} sont consommés : $x_f = C_a V_a$. Il reste des ions \ce{HO-} :
\[ [\ce{HO-}] = \frac{C_b V_b - C_a V_a}{V_a + V_b} \qquad \Longrightarrow \qquad \mathrm{pH} = 14 + \log[\ce{HO-}] \quad (\text{à } \SI{25}{\celsius}) \]
Le mélange final est \textbf{basique} ($\mathrm{pH} > 7$).

\textbf{Cas 3 : l'équivalence} ($C_a V_a = C_b V_b$).

Les deux réactifs sont totalement consommés. La solution ne contient plus que de l'eau et des ions spectateurs (\ce{Na+}, \ce{Cl-}) : c'est une solution de chlorure de sodium, neutre. À \SI{25}{\celsius} :
\[ \boxed{\mathrm{pH} = 7} \]

\begin{definition}[Équivalence acido-basique]
On dit qu'on a atteint l'\cle{équivalence acido-basique} lorsque les réactifs ont été mélangés dans les \textbf{proportions stœchiométriques}, c'est-à-dire lorsque :
\[ C_a V_a = C_b V_b \qquad \text{soit} \qquad n(\ce{H3O+})_{\text{apporté}} = n(\ce{HO-})_{\text{apporté}}. \]
À l'équivalence d'un mélange acide fort -- base forte, le pH vaut $7$ à \SI{25}{\celsius}.
\end{definition}

\begin{popfigure}
\begin{center}
\begin{tikzpicture}[scale=0.92, every node/.style={font=\small}]
% Bécher acide
\draw[thick, popblue] (-0.9,0) -- (-0.9,1.8) (0.9,0) -- (0.9,1.8);
\draw[thick, popblue] (-0.9,0) arc (180:360:0.9);
\fill[popblue!25] (-0.85,0.12) arc (180:360:0.85) -- (0.85,1.1) -- (-0.85,1.1) -- cycle;
\node[popblue, below] at (0,-0.15) {\textbf{Acide fort} : $V_a$, $C_a$};
\node[popdark] at (0,0.55) {\ce{H3O+} : $C_aV_a$ mol};
\node[popmuted, font=\footnotesize] at (0,1.45) {ions \ce{H3O+} et \ce{Cl-}};
% Bécher base
\begin{scope}[xshift=4.2cm]
\draw[thick, poppink] (-0.9,0) -- (-0.9,1.8) (0.9,0) -- (0.9,1.8);
\draw[thick, poppink] (-0.9,0) arc (180:360:0.9);
\fill[poppink!25] (-0.85,0.12) arc (180:360:0.85) -- (0.85,1.1) -- (-0.85,1.1) -- cycle;
\node[poppink, below] at (0,-0.15) {\textbf{Base forte} : $V_b$, $C_b$};
\node[popdark] at (0,0.55) {\ce{HO-} : $C_bV_b$ mol};
\node[popmuted, font=\footnotesize] at (0,1.45) {ions \ce{HO-} et \ce{Na+}};
\end{scope}
% Flèche mélange
\draw[-{Stealth[length=3mm]}, very thick, poporange] (1.6,0.9) -- (2.6,0.9);
\node[poporange, above, font=\footnotesize] at (2.1,0.95) {mélange};
% Bécher résultat
\begin{scope}[xshift=8.6cm]
\draw[thick, popgreen] (-1.1,0) -- (-1.1,2.1) (1.1,0) -- (1.1,2.1);
\draw[thick, popgreen] (-1.1,0) arc (180:360:1.1);
\fill[popgreen!20] (-1.05,0.12) arc (180:360:1.05) -- (1.05,1.5) -- (-1.05,1.5) -- cycle;
\node[popgreen, below] at (0,-0.15) {\textbf{Mélange} : $V_a + V_b$};
\node[popdark, align=center] at (0,0.85) {\footnotesize réaction totale\\ \ce{H3O+ + HO- -> 2 H2O}};
\node[popmuted, font=\footnotesize] at (0,1.8) {reste l'ion en excès + \ce{Na+}, \ce{Cl-}};
\end{scope}
\end{tikzpicture}
\end{center}
\legende{Bilan ionique du mélange : les ions \ce{H3O+} et \ce{HO-} se neutralisent ; seul l'ion apporté en excès survit, dilué dans le volume total $V_a + V_b$.}
\end{popfigure}

\courssub{Méthode de résolution et exemple chiffré}

\begin{methode}[Calculer le pH d'un mélange acide fort + base forte]
\begin{enumerate}
\item \textbf{Calculer les quantités apportées} : $n(\ce{H3O+}) = C_a V_a$ et $n(\ce{HO-}) = C_b V_b$ (volumes en litres !).
\item \textbf{Comparer} ces deux quantités pour identifier le réactif en excès (un tableau d'avancement sécurise le raisonnement).
\item \textbf{Calculer la concentration de l'ion en excès} en divisant la quantité restante par le \emph{volume total} $V_a + V_b$.
\item \textbf{En déduire le pH} : pH $= -\log[\ce{H3O+}]$ si l'acide est en excès ; pH $= 14 + \log[\ce{HO-}]$ si la base est en excès ; pH $= 7$ à l'équivalence.
\end{enumerate}
\end{methode}

\begin{exoresolu}[Mélange acide chlorhydrique -- soude]
\textbf{Énoncé.} On mélange $V_a = \SI{20}{mL}$ d'acide chlorhydrique de concentration $C_a = \SI{0,10}{mol.L^{-1}}$ avec $V_b = \SI{30}{mL}$ d'hydroxyde de sodium de concentration $C_b = \SI{0,050}{mol.L^{-1}}$. Calculer le pH du mélange à \SI{25}{\celsius}.
\tcblower
\textbf{Solution détaillée.}

\emph{Étape 1 : quantités apportées.}
\begin{align*}
n(\ce{H3O+}) &= C_a V_a = 0,10 \times 0,020 = \SI{2,0e-3}{mol} \\
n(\ce{HO-}) &= C_b V_b = 0,050 \times 0,030 = \SI{1,5e-3}{mol}
\end{align*}

\emph{Étape 2 : réactif en excès.} On a $n(\ce{H3O+}) > n(\ce{HO-})$ : l'acide est en excès. Le tableau d'avancement donne $x_f = n(\ce{HO-}) = \SI{1,5e-3}{mol}$.

\emph{Étape 3 : concentration de l'ion en excès.} Il reste :
\[ n_{\text{reste}}(\ce{H3O+}) = 2,0\times 10^{-3} - 1,5\times 10^{-3} = \SI{5,0e-4}{mol} \]
dans un volume total $V_a + V_b = 20 + 30 = \SI{50}{mL} = \SI{0,050}{L}$. Donc :
\[ [\ce{H3O+}] = \frac{5,0\times 10^{-3} - 1,5\times 10^{-3}}{0,050} = \frac{5,0\times 10^{-4}}{0,050} = \SI{1,0e-2}{mol.L^{-1}}. \]

\emph{Étape 4 : pH.}
\[ \mathrm{pH} = -\log(1,0\times 10^{-2}) = 2,0. \]
Le mélange est bien acide, cohérent avec l'excès d'acide.
\end{exoresolu}

\begin{attention}[L'erreur classique : oublier le volume total]
L'erreur la plus fréquente au Baccalauréat consiste à calculer le pH à partir de la \emph{quantité} restante sans la diviser par le volume total $V_a + V_b$. Dans l'exemple précédent, écrire pH $= -\log(5,0\times 10^{-4}) = 3,3$ est \textbf{faux} : le pH dépend d'une \emph{concentration}, pas d'une quantité de matière. Deux mélanges contenant la même quantité d'acide en excès n'ont pas le même pH si leurs volumes totaux diffèrent. Retiens le réflexe : \textbf{quantité restante $\rightarrow$ division par $V_a + V_b$ $\rightarrow$ concentration $\rightarrow$ pH}.
\end{attention}

\begin{exemplebox}[Base en excès]
On mélange $V_a = \SI{10}{mL}$ d'acide nitrique à $C_a = \SI{0,10}{mol.L^{-1}}$ et $V_b = \SI{20}{mL}$ de soude à $C_b = \SI{0,10}{mol.L^{-1}}$.
\[ n(\ce{H3O+}) = 1,0\times 10^{-3}\ \text{mol} < n(\ce{HO-}) = 2,0\times 10^{-3}\ \text{mol} : \text{base en excès.} \]
\[ [\ce{HO-}] = \frac{2,0\times 10^{-3} - 1,0\times 10^{-3}}{0,030} = \SI{3,3e-2}{mol.L^{-1}} \]
\[ \mathrm{pH} = 14 + \log(3,3\times 10^{-2}) = 14 - 1,48 = 12,5. \]
Le mélange est fortement basique, comme attendu.
\end{exemplebox}

\courssub{Vers le dosage acido-basique}

Renversons maintenant la perspective. Dans tout ce qui précède, les concentrations $C_a$ et $C_b$ étaient connues et on calculait le pH. Mais dans la pratique du laboratoire, on utilise souvent la relation d'équivalence dans l'autre sens : pour \emph{déterminer une concentration inconnue}.

\begin{aretenir}[L'essentiel de la section]
\begin{itemize}
\item La réaction \ce{H3O+ + HO- -> 2 H2O} est quasi totale ($K = 10^{14}$), instantanée et exothermique.
\item Le tableau d'avancement permet de traiter méthodiquement les trois cas : acide en excès (pH $< 7$), base en excès (pH $> 7$), équivalence (pH $= 7$ à \SI{25}{\celsius}).
\item Formules clés : $[\ce{H3O+}] = \dfrac{C_aV_a - C_bV_b}{V_a + V_b}$ ou $[\ce{HO-}] = \dfrac{C_bV_b - C_aV_a}{V_a + V_b}$.
\item À l'équivalence : $C_a V_a = C_b V_b$.
\end{itemize}
\end{aretenir}

\begin{savaistu}[Le dosage : une application quotidienne]
La relation $C_a V_a = C_b V_b$ est le cœur du \cle{dosage acido-basique}, technique utilisée chaque jour au Cameroun : contrôle de l'acidité du vin de palme et du vinaigre, vérification des lessives de soude des savonneries artisanales, analyses de qualité à la SONARA ou dans les brasseries. En versant progressivement une base de concentration connue dans un volume précis d'acide de concentration inconnue, et en repérant le volume $V_{b,E}$ pour lequel l'équivalence est atteinte, on remonte à $C_a = \dfrac{C_b V_{b,E}}{V_a}$. C'est précisément l'objet de la leçon suivante : le \textbf{dosage acido-basique}, où nous apprendrons à suivre l'évolution du pH au cours de l'ajout de base et à exploiter la courbe de titrage.
\end{savaistu}

\begin{exercice}[À toi de jouer]
On ajoute $V_b = \SI{25}{mL}$ d'une solution d'hydroxyde de potassium à $C_b = \SI{0,080}{mol.L^{-1}}$ à $V_a = \SI{40}{mL}$ d'acide chlorhydrique à $C_a = \SI{0,050}{mol.L^{-1}}$.
\begin{enumerate}
\item Écrire l'équation de la réaction qui se produit.
\item À l'aide d'un tableau d'avancement, déterminer le réactif en excès.
\item Calculer le pH du mélange à \SI{25}{\celsius}.
\item Quel volume de soude aurait-il fallu verser pour atteindre exactement l'équivalence ?
\end{enumerate}
\end{exercice}

\begin{corrige}[Exercice 1]
\textbf{1.} \ce{H3O+ + HO- -> 2 H2O} (réaction totale, $K = 10^{14}$).

\textbf{2.} $n(\ce{H3O+}) = 0,050 \times 0,040 = \SI{2,0e-3}{mol}$ ; $n(\ce{HO-}) = 0,080 \times 0,025 = \SI{2,0e-3}{mol}$. Les quantités sont \emph{égales} : on est exactement à l'équivalence, aucun réactif n'est en excès.

\textbf{3.} À l'équivalence acide fort -- base forte, pH $= 7$ à \SI{25}{\celsius}. La solution obtenue est une solution neutre de chlorure de potassium.

\textbf{4.} L'équivalence est déjà atteinte avec $V_b = \SI{25}{mL}$ : le volume demandé est donc $V_{b,E} = \dfrac{C_a V_a}{C_b} = \dfrac{0,050 \times 40}{0,080} = \SI{25}{mL}$.
\end{corrige}

\coursec{Autres acides forts et bases fortes usuels ; sécurité}

Dans la section précédente, nous avons appris à calculer le pH d'un mélange d'acide fort et de base forte, en nous appuyant sur nos deux modèles : l'acide chlorhydrique \ce{HCl} et l'hydroxyde de sodium \ce{NaOH}. Une question se pose alors naturellement : existe-t-il d'autres acides forts et d'autres bases fortes ? Les relations $\mathrm{pH} = -\log C$ et $\mathrm{pH} = \mathrm{p}K_e + \log C$ restent-elles valables ? La réponse est oui, et c'est une excellente nouvelle : tout ce que tu as appris dans cette leçon se généralise. Mais attention : ces produits sont puissants, donc dangereux. Cette dernière section dresse l'inventaire des acides forts et bases fortes usuels, puis fixe les règles de sécurité indispensables au laboratoire.

\courssub{Les acides forts usuels}

\begin{definition}[Acide fort]
Un \cle{acide fort} est une espèce chimique qui réagit \textbf{totalement} avec l'eau selon une réaction quantitative (taux d'avancement final $\tau = 1$) :
\[ \ce{AH + H2O -> A- + H3O+} \]
Toute la force d'un acide fort est ainsi « convertie » en ions oxonium \ce{H3O+}. Pour une solution d'acide fort de concentration $C$ (avec $C > \SI{1e-6}{mol.L^{-1}}$) : $\mathrm{pH} = -\log C$.
\end{definition}

L'acide chlorhydrique n'est donc qu'un membre d'une petite famille. Voici les autres acides forts que tu rencontreras au lycée, au laboratoire et dans l'industrie.

\textbf{1. L'acide nitrique \ce{HNO3}.} C'est un liquide incolore, très utilisé pour attaquer les métaux et fabriquer les engrais azotés. Son ionisation dans l'eau est totale :
\[ \ce{HNO3 + H2O -> NO3- + H3O+} \]
L'ion nitrate \ce{NO3-} formé est un ion spectateur sur le plan acido-basique : c'est une base négligeable, il ne réagit pas avec l'eau.

\textbf{2. L'acide bromhydrique \ce{HBr}.} Obtenu par dissolution du bromure d'hydrogène gazeux dans l'eau, il se comporte exactement comme \ce{HCl} :
\[ \ce{HBr + H2O -> Br- + H3O+} \]

\textbf{3. L'acide perchlorique \ce{HClO4}.} C'est l'un des acides les plus forts connus :
\[ \ce{HClO4 + H2O -> ClO4- + H3O+} \]

\textbf{4. L'acide sulfurique \ce{H2SO4}, un cas particulier.} C'est un \emph{diacide} : il possède deux atomes d'hydrogène ionisables. Seule la \textbf{première acidité} est forte (réaction totale avec l'eau) :
\[ \ce{H2SO4 + H2O -> HSO4- + H3O+} \]
L'ion hydrogénosulfate \ce{HSO4-} formé est, lui, un \emph{acide faible} : il ne libère son deuxième proton que partiellement. Au programme de Terminale, pour les solutions diluées usuelles, on retiendra que l'acide sulfurique se comporte comme un acide fort pour sa première acidité.

\begin{attention}[Le piège du diacide]
Ne double pas systématiquement la concentration en \ce{H3O+} pour \ce{H2SO4} ! La deuxième acidité étant faible, on n'a pas $[\ce{H3O+}] = 2C$ en toute rigueur. Au lycée, sauf indication contraire de l'énoncé, on considère seulement la première acidité forte : $[\ce{H3O+}] \approx C$.
\end{attention}

\begin{savaistu}[L'acide sulfurique, roi de l'industrie chimique]
L'acide sulfurique est l'un des produits chimiques les plus fabriqués au monde : plus de 200 millions de tonnes par an ! Sa production sert même d'indicateur du développement industriel d'un pays. Sa principale utilisation est la fabrication des \textbf{engrais phosphatés}, essentiels à l'agriculture intensive — y compris pour les plantations de cacao et de palmier à huile du Cameroun. Il sert aussi dans le raffinage du pétrole (comme à la SONARA), le traitement des minerais et les batteries au plomb des automobiles.
\end{savaistu}

\courssub{Les bases fortes usuelles}

\begin{definition}[Base forte]
Une \cle{base forte} est une espèce chimique dont la mise en solution dans l'eau libère \textbf{quantitativement} des ions hydroxyde \ce{HO-}, soit directement par dissociation totale, soit par réaction totale avec l'eau. Pour une solution de base forte de concentration $C$ (avec $C > \SI{1e-6}{mol.L^{-1}}$) : $[\ce{HO-}] = C$ et $\mathrm{pH} = \mathrm{p}K_e + \log C = 14 + \log C$ à \SI{25}{\celsius}.
\end{definition}

\textbf{1. L'hydroxyde de potassium \ce{KOH} (la potasse).} Solide blanc, très proche de la soude, c'est une base forte : sa dissociation dans l'eau est totale.
\[ \ce{KOH(s) ->[H2O] K+(aq) + HO-(aq)} \]
La potasse est très utilisée dans les savonneries artisanales : c'est elle qui, chauffée avec l'huile de palme, donne le savon noir traditionnel.

\textbf{2. L'hydroxyde de calcium \ce{Ca(OH)2} (chaux éteinte).} Sa dissociation libère \textbf{deux} ions hydroxyde par formule :
\[ \ce{Ca(OH)2(s) ->[H2O] Ca^2+(aq) + 2 HO-(aq)} \]
\begin{attention}[Coefficient stœchiométrique]
Pour l'hydroxyde de calcium, $[\ce{HO-}] = 2C$ et non $C$ ! Le pH d'une solution de \ce{Ca(OH)2} de concentration $C$ vaut donc $\mathrm{pH} = 14 + \log(2C)$. Oublier le facteur 2 est une erreur classique au Baccalauréat. La même remarque vaut pour les diacides forts.
\end{attention}

\textbf{3. L'éthanolate de sodium \ce{C2H5ONa}, un cas remarquable.} Contrairement à la soude ou à la potasse, ce composé \textbf{ne contient aucun ion hydroxyde} dans sa formule ! Pourtant, c'est bien une base forte. Pourquoi ? Parce que sa dissociation libère l'ion \cle{éthanolate} \ce{C2H5O-}, qui est une base si forte qu'elle arrache un proton à la molécule d'eau, de façon \textbf{quasi-totale} :
\[ \ce{C2H5ONa(s) ->[H2O] C2H5O-(aq) + Na+(aq)} \]
\[ \ce{C2H5O- + H2O -> C2H5OH + HO-} \]
Bilan : chaque mole d'éthanolate de sodium mise en solution produit une mole d'ions \ce{HO-}. On a donc bien $[\ce{HO-}] = C$, et la relation $\mathrm{pH} = 14 + \log C$ s'applique normalement.

\begin{definition}[Base forte sans ions \ce{HO-} initiaux]
L'\cle{éthanolate de sodium} \ce{C2H5ONa} est une \textbf{base forte} car l'ion éthanolate (\chemfig{CH_3-CH_2-O^{-}}) réagit \textbf{quasi-totale\-ment} avec l'eau en libérant des ions hydroxyde : \ce{C2H5O- + H2O -> C2H5OH + HO-}. Une base forte n'est donc pas nécessairement un hydroxyde : c'est toute espèce qui génère quantitativement des ions \ce{HO-} dans l'eau.
\end{definition}

\begin{popfigure}
\begin{center}
\begin{tabularx}{\linewidth}{|l|X|X|}
\hline
\textbf{Nom} & \textbf{Formule} & \textbf{Ionisation dans l'eau} \\
\hline
\multicolumn{3}{|l|}{\textbf{Acides forts} ($\mathrm{pH} = -\log C$)} \\
\hline
Acide chlorhydrique & \ce{HCl} & \ce{HCl + H2O -> Cl- + H3O+} \\
Acide nitrique & \ce{HNO3} & \ce{HNO3 + H2O -> NO3- + H3O+} \\
Acide bromhydrique & \ce{HBr} & \ce{HBr + H2O -> Br- + H3O+} \\
Acide perchlorique & \ce{HClO4} & \ce{HClO4 + H2O -> ClO4- + H3O+} \\
Acide sulfurique (1\textsuperscript{re} acidité) & \ce{H2SO4} & \ce{H2SO4 + H2O -> HSO4- + H3O+} \\
\hline
\multicolumn{3}{|l|}{\textbf{Bases fortes} ($\mathrm{pH} = 14 + \log C$)} \\
\hline
Hydroxyde de sodium (soude) & \ce{NaOH} & \ce{NaOH -> Na+ + HO-} \\
Hydroxyde de potassium (potasse) & \ce{KOH} & \ce{KOH -> K+ + HO-} \\
Hydroxyde de calcium (chaux) & \ce{Ca(OH)2} & \ce{Ca(OH)2 -> Ca^2+ + 2 HO-} \\
Éthanolate de sodium & \ce{C2H5ONa} & \ce{C2H5O- + H2O -> C2H5OH + HO-} \\
\hline
\end{tabularx}
\end{center}
\legende{Acides forts et bases fortes usuels : formules et équations d'ionisation (réactions totales avec l'eau).}
\end{popfigure}

\courssub{Généralisation des relations de pH}

\begin{aretenir}[Complément Bac : une seule logique pour tous]
Toutes les relations établies dans cette leçon avec \ce{HCl} et \ce{NaOH} s'appliquent \textbf{de la même façon} quel que soit l'acide fort ou la base forte utilisé(e) :
\begin{itemize}
\item acide fort de concentration $C$ : $[\ce{H3O+}] = C$, donc $\mathrm{pH} = -\log C$ ;
\item base forte de concentration $C$ : $[\ce{HO-}] = C$, donc $\mathrm{pH} = 14 + \log C$ (à \SI{25}{\celsius}) ;
\item dilution au dixième : le pH d'un acide fort augmente de 1 unité, celui d'une base forte diminue de 1 unité ;
\item mélange acide fort + base forte : la réaction \ce{H3O+ + HO- -> 2 H2O} est totale, on dresse un tableau d'avancement.
\end{itemize}
Conditions de validité : $C > \SI{1e-6}{mol.L^{-1}}$ et température de \SI{25}{\celsius} ($\mathrm{p}K_e = 14$).
\end{aretenir}

\begin{exemplebox}[pH d'une solution de potasse]
On dissout de l'hydroxyde de potassium pour obtenir une solution de concentration $C = \SI{5,0e-4}{mol.L^{-1}}$. La dissociation \ce{KOH -> K+ + HO-} étant totale :
\[ [\ce{HO-}] = C = \SI{5,0e-4}{mol.L^{-1}} \]
\[ \mathrm{pH} = 14 + \log C = 14 + \log(5{,}0 \times 10^{-4}) = 14 - 3{,}30 = 10{,}7 \]
La solution, de pH voisin de 11, est bien basique. Vérification : $C > \SI{1e-6}{mol.L^{-1}}$, la formule est valide.
\end{exemplebox}

\begin{exoresolu}[pH d'une solution d'éthanolate de sodium]
\textbf{Énoncé.} On dissout \SI{1,36}{g} d'éthanolate de sodium \ce{C2H5ONa} ($M = \SI{68,0}{g.mol^{-1}}$) dans de l'eau distillée pour obtenir $V = \SI{500}{mL}$ de solution. Calculer le pH de la solution à \SI{25}{\celsius}.
\tcblower
\textbf{Solution.} Étape 1 : quantité de matière dissoute.
\[ n = \frac{m}{M} = \frac{1{,}36}{68{,}0} = \SI{2,0e-2}{mol} \]
Étape 2 : concentration de la solution.
\[ C = \frac{n}{V} = \frac{2{,}0 \times 10^{-2}}{0{,}500} = \SI{4,0e-2}{mol.L^{-1}} \]
Étape 3 : inventaire. La dissociation \ce{C2H5ONa -> C2H5O- + Na+} est totale, puis l'ion éthanolate réagit quasi-totale\-ment avec l'eau : \ce{C2H5O- + H2O -> C2H5OH + HO-}. Donc :
\[ [\ce{HO-}] = C = \SI{4,0e-2}{mol.L^{-1}} \]
Étape 4 : calcul du pH. Comme $C > \SI{1e-6}{mol.L^{-1}}$ :
\[ \mathrm{pH} = 14 + \log(4{,}0 \times 10^{-2}) = 14 - 1{,}40 = 12{,}6 \]
La solution est fortement basique, ce qui confirme que l'éthanolate de sodium est une base forte, bien qu'il ne contienne pas d'ions \ce{HO-} dans sa formule brute.
\end{exoresolu}

\courssub{Sécurité au laboratoire : pictogrammes et conduites à tenir}

Les acides forts concentrés et les bases fortes sont des produits \textbf{corrosifs} : ils détruisent les tissus vivants (peau, yeux, muqueuses) et attaquent même les métaux. Leur manipulation impose une discipline stricte. Les flacons portent des \cle{pictogrammes de danger} normalisés (système GHS) qu'il faut savoir reconnaître instantanément.

\begin{popfigure}
\begin{center}
\begin{tikzpicture}[scale=0.75, every node/.style={transform shape}]
% GHS05 corrosif
\begin{scope}[shift={(0,0)}]
\draw[thick,popdark,rotate=45] (-1.2,-1.2) rectangle (1.2,1.2);
\draw[thick,popink,rotate=45,fill=popink!30] (-1.1,-1.1) rectangle (1.1,1.1);
\draw[very thick,popdark,rotate=45] (-1.2,-1.2) rectangle (1.2,1.2);
\draw[thick,popdark] (-0.55,0.55) -- (-0.15,0.15);
\draw[thick,popdark] (0.55,0.55) -- (0.15,0.15);
\draw[thick,popdark] (-0.15,0.15) -- (-0.35,-0.45);
\draw[thick,popdark] (0.15,0.15) -- (0.35,-0.45);
\draw[thick,popdark] (-0.75,-0.55) -- (0.75,-0.55);
\fill[popink] (-0.45,-0.15) circle (0.06);
\fill[popink] (0.45,-0.15) circle (0.06);
\node[below,font=\bfseries\small] at (0,-1.7) {GHS05 : Corrosif};
\node[below,font=\footnotesize,text width=4cm,align=center] at (0,-2.1) {Acides forts et bases fortes concentrés : brûlures de la peau, lésions oculaires graves};
\end{scope}
% GHS07 nocif/irritant
\begin{scope}[shift={(5.5,0)}]
\draw[thick,popdark,rotate=45] (-1.2,-1.2) rectangle (1.2,1.2);
\draw[thick,poporange,rotate=45,fill=poporange!30] (-1.1,-1.1) rectangle (1.1,1.1);
\draw[very thick,popdark,rotate=45] (-1.2,-1.2) rectangle (1.2,1.2);
\draw[line width=2.5pt,popdark] (-0.5,-0.5) -- (0.5,0.5);
\draw[line width=2.5pt,popdark] (-0.5,0.5) -- (0.5,-0.5);
\node[below,font=\bfseries\small] at (0,-1.7) {GHS07 : Nocif / Irritant};
\node[below,font=\footnotesize,text width=4cm,align=center] at (0,-2.1) {Solutions diluées : irritations de la peau et des yeux, nocivité en cas d'ingestion};
\end{scope}
% Conduites à tenir
\begin{scope}[shift={(11,0)}]
\node[draw=popgreen,thick,fill=popgreen!30,rounded corners,text width=4.2cm,align=left,font=\footnotesize] at (0,0.2)
{\textbf{Conduites à tenir :}\\[2pt]
$\bullet$ Blouse, gants, lunettes\\
$\bullet$ Verser l'acide dans l'eau\\
$\bullet$ Ne jamais pipeter à la bouche\\
$\bullet$ En cas de contact : rincer abondamment à l'eau, prévenir l'enseignant};
\end{scope}
\end{tikzpicture}
\end{center}
\legende{Pictogrammes GHS rencontrés lors de la manipulation des acides forts et bases fortes, et règles de sécurité associées.}
\end{popfigure}

\begin{attention}[La règle d'or : l'acide dans l'eau, jamais l'inverse]
Lorsqu'on dilue un acide fort concentré (notamment l'acide sulfurique commercial à 98 \%), la dissolution est \textbf{très exothermique} : elle dégage énormément de chaleur. Si l'on verse de l'eau dans l'acide concentré, l'eau, moins dense, reste en surface, bout instantanément et provoque des \textbf{projections d'acide bouillant}. La règle de sécurité est donc absolue : \textbf{verser lentement l'acide concentré dans l'eau, jamais l'inverse}, en agitant et éventuellement en refroidissant le récipient. Moyen mnémotechnique : « l'acide plonge dans l'eau, comme le plongeur plonge dans la piscine — et non l'inverse ».
\end{attention}

\textbf{Règles de sécurité essentielles au laboratoire :}
\begin{enumerate}
\item Porter systématiquement une \textbf{blouse} en coton, des \textbf{lunettes de protection} et des \textbf{gants} adaptés dès que l'on manipule un acide fort ou une base forte, même dilué.
\item Travailler sous la \textbf{hotte} pour les produits concentrés ou dégageant des vapeurs (acide chlorhydrique concentré fume à l'air libre).
\item Ne \textbf{jamais pipeter à la bouche} : utiliser une propipette ou une pipette automatique.
\item Verser toujours \textbf{l'acide concentré dans l'eau}, lentement et en agitant.
\item Reboucher immédiatement les flacons et ne jamais reverser un excédent de produit dans le flacon d'origine.
\end{enumerate}

\textbf{Conduite à tenir en cas d'accident :}
\begin{itemize}
\item \textbf{Contact avec la peau} : rincer immédiatement et \textbf{abondamment à l'eau courante} pendant plusieurs minutes, retirer les vêtements souillés, puis prévenir l'enseignant. Ne jamais chercher à « neutraliser » une brûlure acide par une base (ou l'inverse) : la réaction de neutralisation dégage de la chaleur et aggraverait la brûlure.
\item \textbf{Projection dans les yeux} : c'est une urgence absolue. Rincer aussitôt au \textbf{lave-œil} ou à l'eau courante, paupières ouvertes, pendant au moins 10 minutes, et prévenir immédiatement l'enseignant qui alertera les secours. Chaque seconde compte : les lésions oculaires dues aux bases fortes sont particulièrement graves et rapides.
\item \textbf{Ingestion} : ne pas faire vomir, ne rien faire boire, prévenir immédiatement l'enseignant et les secours.
\end{itemize}

\begin{exercice}[Pour vérifier tes acquis]
On dispose d'une solution d'acide nitrique \ce{HNO3} de concentration $C = \SI{2,0e-3}{mol.L^{-1}}$ et d'une solution d'hydroxyde de calcium \ce{Ca(OH)2} de concentration $C' = \SI{1,0e-3}{mol.L^{-1}}$.
\begin{enumerate}
\item Calculer le pH de chacune des deux solutions à \SI{25}{\celsius}.
\item On dilue 10 fois la solution d'acide nitrique. Quel est le nouveau pH ?
\item Citer les deux pictogrammes GHS à apposer sur le flacon d'acide nitrique concentré et la règle à respecter pour le diluer.
\end{enumerate}
\end{exercice}

\begin{corrige}[Exercice]
\begin{enumerate}
\item Acide nitrique, acide fort : $\mathrm{pH} = -\log(2{,}0 \times 10^{-3}) = 2{,}7$. Hydroxyde de calcium, base forte libérant 2 ions \ce{HO-} : $[\ce{HO-}] = 2C' = \SI{2,0e-3}{mol.L^{-1}}$, donc $\mathrm{pH} = 14 + \log(2{,}0 \times 10^{-3}) = 11{,}3$.
\item Dilution au dixième d'un acide fort : le pH augmente d'une unité : $\mathrm{pH} = 3{,}7$.
\item Pictogrammes GHS05 (corrosif) et GHS07 (nocif/irritant). Pour diluer : verser lentement l'acide concentré dans l'eau, jamais l'inverse, avec blouse, gants et lunettes.
\end{enumerate}
\end{corrige}

\begin{aretenir}[L'essentiel de la section]
\begin{itemize}
\item Les acides forts usuels sont \ce{HCl}, \ce{HNO3}, \ce{HBr}, \ce{HClO4} et \ce{H2SO4} (première acidité) ; les bases fortes usuelles sont \ce{NaOH}, \ce{KOH}, \ce{Ca(OH)2} et \ce{C2H5ONa}.
\item L'éthanolate de sodium est une base forte \emph{sans} ions \ce{HO-} initiaux : c'est la réaction quasi-totale \ce{C2H5O- + H2O -> C2H5OH + HO-} qui les génère.
\item Les relations $\mathrm{pH} = -\log C$ et $\mathrm{pH} = 14 + \log C$ s'appliquent à tous les acides forts et toutes les bases fortes (attention aux coefficients stœchiométriques : $2C$ pour \ce{Ca(OH)2}).
\item Pictogrammes à connaître : GHS05 (corrosif) et GHS07 (nocif/irritant).
\item Sécurité : blouse, gants, lunettes ; l'acide concentré se verse dans l'eau, jamais l'inverse ; en cas de contact, rinçage abondant à l'eau et on prévient l'enseignant.
\end{itemize}
\end{aretenir}

\newpage

\coursec{Travaux pratiques}

\courssub{Objectifs du TP}

À l'issue de cette séance de travaux pratiques, tu dois être capable de :

\begin{itemize}
\item préparer par \cle{dilution} des solutions d'acide chlorhydrique de concentrations décroissantes (\SI{1,0e-2}{mol.L^{-1}}, \SI{1,0e-3}{mol.L^{-1}}, \SI{1,0e-4}{mol.L^{-1}}) en utilisant la verrerie jaugée ;
\item mesurer le \cle{pH} de chaque solution au pH-mètre (ou au papier pH de précision) ;
\item comparer le pH mesuré à la valeur théorique $-\log C$ et \textbf{montrer expérimentalement} que l'acide chlorhydrique est un \cle{acide fort}, c'est-à-dire totalement ionisé dans l'eau ;
\item vérifier que chaque dilution au dixième fait \textbf{augmenter le pH d'une unité} ;
\item calculer les concentrations de toutes les espèces chimiques présentes dans une solution d'acide fort et confronter le calcul à l'expérience.
\end{itemize}

\courssub{Matériel et produits}

\begin{center}
\begin{tabularx}{\linewidth}{|l|X|}
\hline
\rowcolor{popblue!40} \textbf{Matériel} & \textbf{Produits} \\
\hline
\rowcolor{popblueL} 1 pH-mètre étalonné (ou papier pH de précision, gamme $1$--$6$) & Solution mère $S_0$ d'acide chlorhydrique à $C_0 = \SI{0,10}{mol.L^{-1}}$ (préparée par l'enseignant à partir de la solution commerciale) \\
\hline
3 fioles jaugées de \SI{100}{mL} & Eau distillée (ou eau de pluie filtrée et bouillie, refroidie) \\
\hline
\rowcolor{popblueL} 1 pipette jaugée de \SI{10,0}{mL} + propipette (poire) & \\
\hline
4 béchers de \SI{100}{mL} propres et secs (ou rincés à la solution) & \\
\hline
\rowcolor{popblueL} 1 agitateur en verre, 1 pissette & \\
\hline
Gants, lunettes de protection, blouse & \\
\hline
\end{tabularx}
\end{center}

\begin{savaistu}[Et si le pH-mètre manque ?]
Dans beaucoup de lycées camerounais, le pH-mètre est rare ou en panne. Le \cle{papier pH} de précision (gamme étroite $1$--$4$ ou $3$--$6$) donne une valeur au demi-unité près : cela suffit pour vérifier le saut d'une unité de pH par dilution. On peut aussi utiliser un \textbf{jus de bissap} (infusion de fleurs d'hibiscus) comme indicateur coloré maison : il vire du rouge vif au rose pâle quand le pH augmente, ce qui illustre qualitativement la dilution, sans remplacer la mesure.
\end{savaistu}

\courssub{Sécurité}

\begin{attention}[L'acide chlorhydrique : un corrosif à manipuler avec respect]
La solution commerciale d'acide chlorhydrique porte le \textbf{pictogramme de corrosion} (liquide attaquant une barre métallique et une main). Même diluée à \SI{0,10}{mol.L^{-1}}, la solution $S_0$ reste irritante.
\begin{itemize}
\item \textbf{Port obligatoire} des lunettes et des gants ; blouse fermée, cheveux attachés.
\item \textbf{Ne jamais pipeter à la bouche} : utiliser systématiquement la propipette.
\item En cas de contact avec la peau : rincer abondamment à l'eau courante pendant plusieurs minutes, puis prévenir l'enseignant.
\item \textbf{Règle d'or de la dilution} : c'est toujours l'enseignant qui a préparé $S_0$ en versant \emph{l'acide dans l'eau}, jamais l'inverse (risque de projections dues à l'échauffement).
\item Ne pas renverser la solution sur la paillasse ; en cas de bris de verre, ne pas ramasser à mains nues.
\item Les solutions diluées usagées peuvent être jetées à l'évier avec beaucoup d'eau courante, selon les consignes de l'enseignant.
\end{itemize}
\end{attention}

\courssub{Schéma du dispositif}

\begin{popfigure}
\begin{center}
\begin{tikzpicture}[scale=0.95, every node/.style={font=\small}]
% Bécher avec pH-mètre
\draw[thick, popdark] (0,0) -- (0,2.6) (2.4,0) -- (2.4,2.6);
\draw[thick, popdark] (0,0) .. controls (0.3,-0.35) and (2.1,-0.35) .. (2.4,0);
\draw[popblue, fill=popblue!25] (0.06,0.1) .. controls (0.35,-0.28) and (2.05,-0.28) .. (2.34,0.1) -- (2.34,1.5) -- (0.06,1.5) -- cycle;
\draw[popblue] (0.06,1.5) -- (2.34,1.5);
\node[popdark] at (1.2,-0.75) {Bécher : solution à mesurer};
% Sonde pH
\draw[thick, poporange, fill=poporange!20] (1.55,3.6) rectangle (1.95,0.9);
\draw[poporange, fill=poporange!60] (1.6,0.9) rectangle (1.9,1.25);
\draw[thick, poporange] (1.75,3.6) -- (1.75,4.3);
% Boîtier pH-mètre
\draw[thick, poppurple, fill=poppurple!15, rounded corners=4pt] (3.4,3.3) rectangle (6.2,4.6);
\draw[thick, poppurple] (1.75,4.3) .. controls (2.6,4.9) .. (3.4,4.1);
\draw[fill=white, draw=poppurple] (3.7,3.75) rectangle (5.1,4.3);
\node[poppurple] at (4.4,4.02) {\textbf{pH = 2,98}};
\node[poppurple] at (4.8,3.0) {pH-mètre étalonné};
% Fiole jaugée
\begin{scope}[xshift=8.2cm]
\draw[thick, popdark] (0,0) .. controls (0.1,-0.3) and (1.9,-0.3) .. (2,0);
\draw[thick, popdark] (0,0) -- (0,1.3) .. controls (0.05,1.9) and (0.75,2.0) .. (0.8,2.5) -- (0.8,3.3);
\draw[thick, popdark] (2,0) -- (2,1.3) .. controls (1.95,1.9) and (1.25,2.0) .. (1.2,2.5) -- (1.2,3.3);
\draw[thick, popdark] (0.72,3.3) -- (1.28,3.3);
\draw[popgreen, fill=popgreen!25] (0.06,0.05) .. controls (0.15,-0.24) and (1.85,-0.24) .. (1.94,0.05) -- (1.94,1.1) -- (0.06,1.1) -- cycle;
\draw[popgreen] (0.06,1.1) -- (1.94,1.1);
\draw[popdark] (0.8,2.9) -- (1.2,2.9);
\node[popdark, right] at (1.3,2.9) {\scriptsize trait de jauge};
\node[popdark] at (1,-0.75) {Fiole jaugée \SI{100}{mL}};
\end{scope}
% Flèche
\draw[-{Stealth[length=3mm]}, very thick, popgold] (2.9,1.4) -- (7.9,1.4) node[midway, above, popdark, align=center] {prélèvement \SI{10,0}{mL}\\ + eau distillée};
\end{tikzpicture}
\end{center}
\legende{Dispositif du TP : mesure du pH au pH-mètre (sonde plongée dans le bécher) et préparation des solutions filles par dilution en fiole jaugée de \SI{100}{mL}.}
\end{popfigure}

\courssub{Protocole}

\textbf{Partie A — Préparation des solutions par dilutions successives au dixième}

\begin{enumerate}
\item Verser environ \SI{50}{mL} de la solution mère $S_0$ ($C_0 = \SI{0,10}{mol.L^{-1}}$) dans un bécher propre. \textbf{Rincer la pipette jaugée avec un peu de $S_0$} (jeter le rinçage).
\item À l'aide de la \textbf{pipette jaugée de \SI{10,0}{mL}} munie de la propipette, prélever $V_0 = \SI{10,0}{mL}$ de $S_0$ et les introduire dans une \textbf{fiole jaugée de \SI{100}{mL}}.
\item Compléter avec de l'eau distillée jusqu'aux deux tiers, boucher, agiter, puis ajuster au trait de jauge (le bas du ménisque affleure le trait, œil à hauteur du trait). Homogénéiser en retournant la fiole dix fois. On obtient la solution $S_1$ de concentration $C_1 = \SI{1,0e-2}{mol.L^{-1}}$.
\item Rincer la pipette avec un peu de $S_1$ (jeter le rinçage), puis répéter l'opération : prélever \SI{10,0}{mL} de $S_1$, diluer dans une fiole de \SI{100}{mL} $\Rightarrow$ solution $S_2$ ($C_2 = \SI{1,0e-3}{mol.L^{-1}}$).
\item Recommencer à partir de $S_2$ (rinçage pipette à $S_2$) $\Rightarrow$ solution $S_3$ ($C_3 = \SI{1,0e-4}{mol.L^{-1}}$).
\end{enumerate}

\textbf{Partie B — Mesures de pH}

\begin{enumerate}
\setcounter{enumi}{5}
\item Vérifier que le pH-mètre est \textbf{étalonné} (solutions tampons pH 4 et pH 7). Rincer la sonde à l'eau distillée et l'essuyer délicatement (papier absorbant non pelucheux) entre deux mesures.
\item Verser environ \SI{30}{mL} de chaque solution dans un bécher, \textbf{en commençant par la plus diluée} ($S_3$, puis $S_2$, $S_1$, $S_0$) pour limiter la pollution de la sonde par entraînement.
\item Plonger la sonde, attendre la stabilisation de l'affichage, noter le pH au centième près. Relever aussi la température de la salle.
\item \emph{Alternative papier pH :} déposer une goutte de solution sur le papier à l'aide de l'agitateur, comparer immédiatement la couleur à l'échelle de référence.
\end{enumerate}

\courssub{Observations et résultats attendus}

Les solutions sont toutes \textbf{incolores} : rien ne distingue visuellement $S_0$ de $S_3$. Seule la mesure de pH les différencie. Tableau de mesures type (température de la salle : \SI{27}{\celsius}) :

\begin{center}
\begin{tabularx}{\linewidth}{|l|X|X|X|X|}
\hline
\rowcolor{popgreen!40} \textbf{Solution} & \textbf{$C$ (\si{mol.L^{-1}})} & \textbf{$-\log C$} & \textbf{pH mesuré} & \textbf{$\Delta$pH} \\
\hline
\rowcolor{popgreenL} $S_0$ & $1,0\times 10^{-1}$ & $1,00$ & $1,02$ & --- \\
\hline
$S_1$ & $1,0\times 10^{-2}$ & $2,00$ & $2,03$ & $+1,01$ \\
\hline
\rowcolor{popgreenL} $S_2$ & $1,0\times 10^{-3}$ & $3,00$ & $2,98$ & $+0,95$ \\
\hline
$S_3$ & $1,0\times 10^{-4}$ & $4,00$ & $4,05$ & $+1,07$ \\
\hline
\end{tabularx}
\end{center}

On constate que, aux incertitudes de mesure près, $\text{pH} \approx -\log C$ et que chaque dilution au dixième augmente le pH d'\textbf{environ une unité}.

\courssub{Exploitation}

\begin{exoresolu}[Exploitation guidée du TP]
\textbf{Questions.}
\begin{enumerate}
\item Écrire l'équation de la réaction du chlorure d'hydrogène \ce{HCl} avec l'eau.
\item Pour chaque solution, comparer le pH mesuré à $-\log C$. Que peut-on en conclure sur l'avancement de la réaction de \ce{HCl} avec l'eau ?
\item Le graphe $\text{pH} = f(-\log C)$ donne une droite. Quelles sont sa pente et son ordonnée à l'origine ? Conclure.
\item Faire l'inventaire des espèces chimiques présentes dans $S_2$ ($C_2 = \SI{1,0e-3}{mol.L^{-1}}$, pH mesuré $= 2,98$) et calculer leurs concentrations à \SI{25}{\celsius} ($pK_e = 14,0$).
\item Vérifier l'électroneutralité et la conservation de la matière pour $S_2$.
\item Pourquoi la formule $\text{pH} = -\log C$ cesserait-elle d'être valable pour une solution à \SI{1,0e-8}{mol.L^{-1}} ?
\end{enumerate}
\tcblower
\textbf{Solution détaillée.}
\begin{enumerate}
\item \[ \ce{HCl + H2O -> H3O+ + Cl-} \]
\item Pour les quatre solutions, $\text{pH mesuré} \approx -\log C$ (écarts inférieurs à $0,05$ unité, imputables à la sonde, à la température et aux arrondis). Or $-\log C = -\log C_0$ correspond à $[\ce{H3O+}] = C$, c'est-à-dire à une ionisation \textbf{totale}. La réaction de \ce{HCl} avec l'eau est donc \textbf{quasi totale} : \ce{HCl} est un \cle{acide fort}.
\item Les points s'alignent sur une droite de \textbf{pente $+1$} passant quasiment par l'origine : $\text{pH} = -\log C$. C'est la signature expérimentale d'un acide fort. Chaque dilution au dixième ($C$ divisée par 10) augmente le pH d'une unité.
\item Espèces présentes : \ce{H3O+}, \ce{Cl-}, \ce{OH-} et \ce{H2O} (majoritaire). La molécule \ce{HCl} est \textbf{absente} (réaction totale).
\begin{align*}
[\ce{H3O+}] &= 10^{-\text{pH}} = 10^{-2,98} \approx \SI{1,05e-3}{mol.L^{-1}}\\
[\ce{OH-}] &= \frac{K_e}{[\ce{H3O+}]} = \frac{10^{-14}}{1,05\times 10^{-3}} \approx \SI{9,5e-12}{mol.L^{-1}}\\
[\ce{Cl-}] &\approx C_2 = \SI{1,0e-3}{mol.L^{-1}}
\end{align*}
\item Électroneutralité : $[\ce{H3O+}] = [\ce{Cl-}] + [\ce{OH-}]$ ; en effet $1,05\times 10^{-3} \approx 1,0\times 10^{-3} + 9,5\times 10^{-12}$ (l'écart de $5\%$ vient de l'incertitude sur le pH mesuré). Conservation du chlore : $[\ce{Cl-}] = C_2$ puisque tout \ce{HCl} s'est ionisé.
\item Pour $C = \SI{1,0e-8}{mol.L^{-1}}$, les ions \ce{H3O+} apportés par l'acide ($10^{-8}$) deviennent \textbf{négligeables devant ceux de l'autoprotolyse de l'eau} ($10^{-7}$) : le pH tend vers 7 et non vers 8. La formule $\text{pH} = -\log C$ n'est valable que pour $C \gtrsim \SI{1,0e-6}{mol.L^{-1}}$.
\end{enumerate}
\end{exoresolu}

\begin{attention}[Pièges fréquents lors de ce TP]
\begin{itemize}
\item \textbf{Oublier de rincer la sonde} entre deux mesures : la goutte de solution concentrée restante fausse la mesure suivante.
\item Mesurer d'abord $S_0$ puis $S_3$ sans rinçage soigneux : procéder toujours \textbf{de la plus diluée vers la plus concentrée}.
\item Confondre pipette \textbf{jaugée} (un trait, volume précis) et pipette graduée : la dilution au dixième exige la jaugée.
\item Ajuster la fiole \textbf{au-dessus} du trait de jauge : le bas du ménisque doit être tangent au trait, œil à hauteur.
\item Conclure « pH $= 4$ donc acide faible » pour $S_3$ : faux ! Un pH égal à $-\log C$ est précisément la preuve d'un acide \textbf{fort} dilué.
\end{itemize}
\end{attention}

\courssub{Conclusion}

Ce TP a permis de vérifier expérimentalement les deux propriétés fondamentales d'un acide fort. D'une part, le pH mesuré des solutions d'acide chlorhydrique coïncide avec $-\log C$ : la réaction $\ce{HCl + H2O -> H3O+ + Cl-}$ est \textbf{totale}, il ne subsiste aucune molécule \ce{HCl} en solution. D'autre part, chaque dilution au dixième fait croître le pH d'\textbf{une unité}, ce qui se traduit par une droite de pente $1$ sur le graphe $\text{pH} = f(-\log C)$. Le calcul des concentrations des espèces de $S_2$ confirme la cohérence du modèle : $[\ce{H3O+}] \approx [\ce{Cl-}] \approx C$ et $[\ce{OH-}]$ ultra-minoritaire. Ces acquis seront réutilisés de façon symétrique pour les bases fortes ($\text{pH} = pK_e + \log C$) et serviront de référence pour reconnaître, par comparaison, les acides et bases \emph{faibles} de la prochaine leçon.

\newpage

\coursec{Méthodes et exercices résolus}

\coursec{Acides forts et bases fortes}

\courssub{1. L'acide chlorhydrique, modèle d'acide fort}

\begin{aretenir}[Définition : acide fort]
Un \cle{acide fort} est une espèce chimique qui réagit \textbf{totalement} avec l'eau selon l'équation :
\[ \ce{HA + H2O -> A- + H3O+} \]
L'avancement maximal de cette transformation est égal à la quantité de matière d'acide introduite : $x_{\max} = n_{\text{HA, initial}}$. En solution, il \textbf{n'existe plus de molécules HA} : l'acide est totalement ionisé. Les acides forts usuels sont : \ce{HCl}, \ce{HBr}, \ce{HI}, \ce{HNO3}, \ce{HClO4}, \ce{H2SO4} (premier proton seulement).
\end{aretenir}

\begin{propriete}[pH d'une solution d'acide fort]
Pour une solution d'acide fort de concentration $C$ (en \SI{}{mol.L^{-1}}), la concentration en ions oxonium est :
\[ [\ce{H3O+}] = C \qquad (\text{si } C > \SI{1e-6}{mol.L^{-1}}) \]
Le pH s'en déduit directement :
\[ \boxed{\text{pH} = -\log C} \]
À \SI{25}{\celsius}, une solution d'acide fort a toujours un $\text{pH} < 7$.
\end{propriete}

\begin{methode}[Prouver le caractère fort d'un acide à partir du pH et de la concentration]
\begin{enumerate}
\item \textbf{Écrire l'hypothèse de réaction totale} : si l'acide \ce{HA} réagit totalement avec l'eau (\ce{HA + H2O -> A- + H3O+}), chaque mole d'acide introduite libère une mole d'ions \ce{H3O+}.
\item \textbf{Calculer la concentration théorique} en oxonium attendue : $[\ce{H3O+}]_{\text{théo}} = C$ (concentration apportée de l'acide).
\item \textbf{Calculer la concentration réelle} à partir du pH mesuré : $[\ce{H3O+}]_{\text{réel}} = 10^{-\text{pH}}$.
\item \textbf{Comparer} les deux valeurs : si $[\ce{H3O+}]_{\text{réel}} = [\ce{H3O+}]_{\text{théo}}$ (à la précision des mesures près, c'est-à-dire même ordre de grandeur et mêmes chiffres significatifs), la réaction est \cle{totale} : l'acide est \cle{fort}. Si $[\ce{H3O+}]_{\text{réel}} < C$, l'acide est faible.
\item \textbf{Conclure} par une phrase complète : « L'acide est totalement ionisé dans l'eau, c'est un acide fort ; on peut écrire $\text{pH} = -\log C$ ».
\end{enumerate}
\textbf{Réflexe pour une base forte} : comparer $[\ce{HO-}]_{\text{réel}} = 10^{\text{pH}-pK_e}$ à la concentration apportée $C$.
\end{methode}

\begin{exoresolu}[Le chlorure d'hydrogène est-il un acide fort ?]
On dissout du chlorure d'hydrogène gazeux \ce{HCl} dans l'eau de façon à obtenir une solution $S$ d'acide chlorhydrique de concentration $C = \SI{2,0e-3}{mol.L^{-1}}$. La mesure au pH-mètre donne $\text{pH} = 2,70$. Montrer que \ce{HCl} est un acide fort et écrire l'équation de sa réaction avec l'eau.
\tcblower
\textbf{Hypothèse de réaction totale.} Si \ce{HCl} réagit totalement avec l'eau :
\[ \ce{HCl + H2O -> Cl- + H3O+} \]
alors la concentration en ions oxonium serait égale à la concentration apportée :
\[ [\ce{H3O+}]_{\text{théo}} = C = \SI{2,0e-3}{mol.L^{-1}}. \]
\textbf{Concentration réelle déduite du pH.}
\[ [\ce{H3O+}]_{\text{réel}} = 10^{-\text{pH}} = 10^{-2,70} = \SI{1,995e-3}{mol.L^{-1}} \approx \SI{2,0e-3}{mol.L^{-1}}. \]
\textbf{Comparaison.} On constate que, aux chiffres significatifs près (deux chiffres significatifs donnés par $C = 2,0\times10^{-3}$ et $\text{pH}=2,70$) :
\[ [\ce{H3O+}]_{\text{réel}} \approx \SI{2,0e-3}{mol.L^{-1}} = C = [\ce{H3O+}]_{\text{théo}}. \]
\textbf{Conclusion.} Toutes les molécules de \ce{HCl} introduites ont réagi avec l'eau : la transformation est \cle{totale}. Le chlorure d'hydrogène est donc un \cle{acide fort}, et pour ses solutions on peut écrire directement $\text{pH} = -\log C$. Vérification : $-\log(2,0\times10^{-3}) = 2,70$.
\end{exoresolu}

\begin{experience}[Mesure du pH de solutions d'acide chlorhydrique de concentrations connues]
\textbf{Objectif.} Vérifier la relation $\text{pH} = -\log C$ pour un acide fort.
\textbf{Matériel.} pH-mètre étalonné (solutions tampons pH 4,01 et 7,00), béchers, solutions mères \ce{HCl} $C_0 = \SI{1,0e-1}{mol.L^{-1}}$, eau distillée, matériel de dilution (pipette, fiole jaugée \SI{100}{mL}).
\textbf{Protocole.}
\begin{enumerate}
\item Préparer par dilution successive 4 solutions filles : $C_1 = \SI{1,0e-2}{M}$, $C_2 = \SI{1,0e-3}{M}$, $C_3 = \SI{1,0e-4}{M}$, $C_4 = \SI{1,0e-5}{M}$.
\item Rincer l'électrode à l'eau distillée, l'essuyer délicatement (sans frotter).
\item Mesurer le pH de chaque solution en laissant stabiliser la valeur. Noter les résultats.
\end{enumerate}
\textbf{Observations et interprétation.}
\begin{center}
\begin{tabularx}{\linewidth}{|c|c|c|c|}
\hline
\rowcolor{popblueL} Solution & $C$ (\SI{}{mol.L^{-1}}) & $\text{pH}_{\text{théo}} = -\log C$ & $\text{pH}_{\text{mesuré}}$ \\
\hline
Mère $S_0$ & $1,0\times10^{-1}$ & 1,00 & $\approx 1,0$ \\
Fille $S_1$ & $1,0\times10^{-2}$ & 2,00 & $\approx 2,0$ \\
Fille $S_2$ & $1,0\times10^{-3}$ & 3,00 & $\approx 3,0$ \\
Fille $S_3$ & $1,0\times10^{-4}$ & 4,00 & $\approx 4,0$ \\
Fille $S_4$ & $1,0\times10^{-5}$ & 5,00 & $\approx 5,0$ (légèrement $< 5$) \\
\hline
\end{tabularx}
\end{center}
Pour $C \ge \SI{1e-4}{M}$, $\text{pH}_{\text{mesuré}} \approx \text{pH}_{\text{théo}}$. Pour $C = \SI{1e-5}{M}$, le pH mesuré est légèrement inférieur à 5 (ex: 4,98) car l'autoprotolyse de l'eau commence à contribuer non négligeablement aux ions \ce{H3O+}. La formule $\text{pH} = -\log C$ n'est plus valide pour $C < \SI{1e-6}{M}$.
\end{experience}

\courssub{2. L'hydroxyde de sodium, modèle de base forte}

\begin{aretenir}[Définition : base forte]
Une \cle{base forte} est une espèce qui, en solution aqueuse, fournit totalement des ions hydroxyde \ce{HO-}. Les hydroxydes des métaux alcalins (\ce{NaOH}, \ce{KOH}, \ce{LiOH}) et des métaux alcalino-terreux (\ce{Ba(OH)2}, \ce{Ca(OH)2}, \ce{Sr(OH)2}) sont des bases fortes (solides ioniques totalement dissociés). L'éthanolate de sodium \ce{NaOEt} (utilisé en synthèse organique) est aussi une base forte : \ce{NaOEt -> Na+ + ^-OEt} et \ce{^-OEt + H2O -> EtOH + HO-} (totale).
\end{aretenir}

\begin{propriete}[pH d'une solution de base forte]
Pour une solution de base forte de concentration $C$ (en \SI{}{mol.L^{-1}}), la concentration en ions hydroxyde est :
\[ [\ce{HO-}] = C \qquad (\text{si } C > \SI{1e-6}{mol.L^{-1}}) \]
On utilise le produit ionique de l'eau $K_e = [\ce{H3O+}][\ce{HO-}] = 10^{-14}$ (à \SI{25}{\celsius}) :
\[ [\ce{H3O+}] = \frac{K_e}{C} \quad\Longrightarrow\quad \boxed{\text{pH} = pK_e + \log C = 14 + \log C} \]
À \SI{25}{\celsius}, une solution de base forte a toujours un $\text{pH} > 7$.
\end{propriete}

\begin{methode}[Calculer le pH d'une solution d'acide fort ou de base forte]
\begin{enumerate}
\item \textbf{Identifier la nature du soluté} : acide fort (\ce{HCl}, \ce{HNO3}, \ce{HBr}...) ou base forte (\ce{NaOH}, \ce{KOH}, \ce{Ba(OH)2}, éthanolate de sodium...). Attention : pour \ce{Ba(OH)2}, $[\ce{HO-}] = 2C$.
\item \textbf{Acide fort} de concentration $C$ : $\text{pH} = -\log C$.
\item \textbf{Base forte} de concentration $C$ : $\text{pH} = pK_e + \log C$ (avec $pK_e = 14$ à \SI{25}{\celsius}).
\item \textbf{Vérifier le domaine de validité} : $C > \SI{1e-6}{mol.L^{-1}}$. En dessous, les ions de l'autoprotolyse de l'eau ne sont plus négligeables.
\item \textbf{Contrôle de cohérence} : acide fort $\Rightarrow \text{pH} < 7$ ; base forte $\Rightarrow \text{pH} > 7$.
\end{enumerate}
\end{methode}

\begin{exoresolu}[pH d'une soude de débouchage diluée]
Une solution d'hydroxyde de sodium (soude) a pour concentration $C_b = \SI{5,0e-2}{mol.L^{-1}}$ à \SI{25}{\celsius}. On donne $pK_e = 14$.
\begin{enumerate}
\item Écrire l'équation de la dissolution de la soude dans l'eau.
\item Calculer les concentrations $[\ce{HO-}]$ et $[\ce{H3O+}]$, puis le pH de la solution.
\end{enumerate}
\tcblower
\textbf{1. Équation de dissolution.} L'hydroxyde de sodium est un solide ionique qui se dissocie totalement ; les ions hydroxyde constituent une base forte :
\[ \ce{NaOH(s) ->[eau] Na+(aq) + HO-(aq)} \]
\textbf{2. Concentrations et pH.} La dissociation étant totale :
\[ [\ce{HO-}] = C_b = \SI{5,0e-2}{mol.L^{-1}}. \]
Par le produit ionique de l'eau :
\[ [\ce{H3O+}] = \frac{K_e}{[\ce{HO-}]} = \frac{10^{-14}}{5,0\times10^{-2}} = \SI{2,0e-13}{mol.L^{-1}}. \]
D'où le pH :
\[ \text{pH} = -\log[\ce{H3O+}] = -\log(2,0\times10^{-13}) = 12,70. \]
Autre rédaction, plus rapide, avec la formule du cours :
\[ \text{pH} = pK_e + \log C_b = 14 + \log(5,0\times10^{-2}) = 14 - 1,30 = 12,70. \]
\textbf{Conclusion.} La solution a un $\text{pH} = 12,70 > 7$ : elle est fortement basique, ce qui est cohérent pour une base forte de concentration supérieure à $\SI{1e-6}{mol.L^{-1}}$.
\end{exoresolu}

\courssub{3. Espèces chimiques en solution : inventaire et calcul des concentrations}

\begin{methode}[Dresser l'inventaire complet et calculer toutes les concentrations]
\begin{enumerate}
\item \textbf{Lister les espèces} : les ions issus du soluté (réaction totale), l'eau \ce{H2O} (espèce majoritaire, $\approx \SI{55,5}{mol.L^{-1}}$), et \emph{toujours} les ions \ce{H3O+} et \ce{HO-} présents dans toute solution aqueuse.
\item \textbf{Concentration en \ce{H3O+}} : $[\ce{H3O+}] = 10^{-\text{pH}}$.
\item \textbf{Concentration en \ce{HO-}} par le produit ionique : $[\ce{HO-}] = \dfrac{K_e}{[\ce{H3O+}]} = 10^{\text{pH}-pK_e}$.
\item \textbf{Conservation de la matière} : pour un acide fort \ce{HA} apporté à $C$, $[\ce{A-}] = C$ ; pour une base forte \ce{MOH}, $[\ce{M+}] = C$ (ions \emph{spectateurs}).
\item \textbf{Vérifier l'électroneutralité} : $\sum [\text{cations}] = \sum [\text{anions}]$. C'est un contrôle indispensable.
\end{enumerate}
\textbf{Réflexe} : dans un acide fort, $[\ce{HO-}]$ est très faible mais \emph{jamais nulle} ; dans une base forte, c'est $[\ce{H3O+}]$ qui est très faible. Ne jamais écrire « concentration nulle ».
\end{methode}

\begin{exoresolu}[Inventaire d'une solution d'acide chlorhydrique]
Une solution $S$ d'acide chlorhydrique a un $\text{pH} = 3,00$ à \SI{25}{\celsius}.
\begin{enumerate}
\item Faire l'inventaire des espèces chimiques présentes en solution.
\item Calculer la concentration de chaque espèce ionique.
\item Vérifier l'électroneutralité de la solution.
\end{enumerate}
\tcblower
\textbf{1. Inventaire.} L'acide chlorhydrique est totalement ionisé : \ce{HCl + H2O -> Cl- + H3O+}. La solution contient donc :
\begin{itemize}
\item les ions : \ce{H3O+}, \ce{Cl-}, \ce{HO-} ;
\item l'eau \ce{H2O}, espèce ultra-majoritaire ($\approx \SI{55,5}{mol.L^{-1}}$) ;
\item \textbf{aucune} molécule \ce{HCl} (réaction totale).
\end{itemize}
\textbf{2. Concentrations.}
\[ [\ce{H3O+}] = 10^{-\text{pH}} = 10^{-3,00} = \SI{1,00e-3}{mol.L^{-1}}. \]
\[ [\ce{HO-}] = \frac{K_e}{[\ce{H3O+}]} = \frac{10^{-14}}{10^{-3}} = \SI{1,00e-11}{mol.L^{-1}}. \]
L'acide étant fort, la conservation de l'élément chlore donne :
\[ [\ce{Cl-}] = C = [\ce{H3O+}] = \SI{1,00e-3}{mol.L^{-1}}. \]
\textbf{3. Électroneutralité.}
\[ \underbrace{[\ce{H3O+}]}_{\text{cation}} = \underbrace{[\ce{Cl-}] + [\ce{HO-}]}_{\text{anions}} \]
\[ \SI{1,00e-3}{mol.L^{-1}} = \SI{1,00e-3}{mol.L^{-1}} + \SI{1,00e-11}{mol.L^{-1}} \]
La relation est vérifiée à la précision des chiffres significatifs près ($10^{-11}$ est négligeable devant $10^{-3}$).
\textbf{Conclusion.} La solution est électriquement neutre ; les ions \ce{HO-}, bien que présents, y sont en concentration négligeable ($10^8$ fois moins que \ce{H3O+}).
\end{exoresolu}

\courssub{4. Effet de la dilution sur le pH}

\begin{propriete}[Loi de dilution et variation du pH]
Soit une solution d'acide fort ou de base forte de concentration $C$ et de pH initial. On la dilue $n$ fois ($n > 1$) : la nouvelle concentration est $C' = \dfrac{C}{n}$.
\begin{itemize}
\item \textbf{Acide fort} : $\text{pH}' = -\log\dfrac{C}{n} = -\log C + \log n = \text{pH} + \log n$. Le pH \textbf{augmente} de $\log n$.
\item \textbf{Base forte} : $\text{pH}' = pK_e + \log\dfrac{C}{n} = pK_e + \log C - \log n = \text{pH} - \log n$. Le pH \textbf{diminue} de $\log n$.
\end{itemize}
\textbf{Cas phare : dilution au dixième ($n=10$).} $\log 10 = 1$.
\begin{itemize}
\item Acide fort : le pH \textbf{augmente de 1 unité}.
\item Base forte : le pH \textbf{diminue de 1 unité}.
\end{itemize}
\end{propriete}

\begin{attention}[Limites de validité de la formule $\text{pH} = -\log C$ (ou $pK_e + \log C$)]
Ces formules supposent que la concentration de l'acide (ou de la base) est \textbf{beaucoup plus grande} que la concentration en ions \ce{H3O+} (ou \ce{HO-}) issue de l'autoprotolyse de l'eau ($\SI{1e-7}{M}$ à \SI{25}{\celsius}).
\textbf{Condition :} $C > \SI{1e-6}{mol.L^{-1}}$.
Si $C \le \SI{1e-6}{M}$, il faut résoudre l'électroneutralité exacte : $[\ce{H3O+}] = C + [\ce{HO-}]$ (acide) ou $[\ce{HO-}] = C + [\ce{H3O+}]$ (base).
\textbf{Conséquence majeure :} par dilution, le pH d'un acide fort \emph{tend vers 7 par valeurs inférieures} sans jamais l'atteindre ni le dépasser. On ne peut \textbf{jamais} rendre une solution acide basique par simple dilution !
\end{attention}

\begin{methode}[Prévoir le pH après dilution / Préparer une solution à pH imposé]
\begin{enumerate}
\item \textbf{Calculer le facteur de dilution} $n = \dfrac{V_{\text{final}}}{V_{\text{initial}}} = \dfrac{C_{\text{initial}}}{C_{\text{final}}}$.
\item \textbf{Appliquer la variation de pH} : $\Delta\text{pH} = \pm \log n$ ($+$ pour acide, $-$ pour base).
\item \textbf{Pour préparer une solution à pH cible} : déduire $C_{\text{cible}} = 10^{-\text{pH}}$ (acide) ou $C_{\text{cible}} = 10^{\text{pH}-pK_e}$ (base), puis utiliser $C_i V_i = C_f V_f$.
\item \textbf{Vérifier} que $C_{\text{final}} > \SI{1e-6}{mol.L^{-1}}$.
\end{enumerate}
\end{methode}

\begin{exoresolu}[Préparation d'une solution étalon par dilution]
Au laboratoire du lycée, on dispose d'une solution mère $S_0$ d'acide chlorhydrique de concentration $C_0 = \SI{1,0e-1}{mol.L^{-1}}$.
\begin{enumerate}
\item Calculer le pH de $S_0$.
\item On prélève $V_0 = \SI{10,0}{mL}$ de $S_0$ que l'on verse dans une fiole jaugée de $V_1 = \SI{100,0}{mL}$ et on complète avec de l'eau distillée. Calculer la concentration $C_1$ et le pH de la solution fille $S_1$.
\item Quel volume de $S_1$ faudrait-il diluer à \SI{100,0}{mL} pour obtenir une solution de $\text{pH} = 4,00$ ? Justifier la possibilité ou l'impossibilité.
\end{enumerate}
\tcblower
\textbf{1. pH de la solution mère.} \ce{HCl} est un acide fort :
\[ \text{pH}_0 = -\log C_0 = -\log(1,0\times10^{-1}) = 1,00. \]
\textbf{2. Solution fille.} La quantité de matière d'acide se conserve lors de la dilution :
\[ C_0 V_0 = C_1 V_1 \quad\Longrightarrow\quad C_1 = \frac{C_0 V_0}{V_1} = \frac{1,0\times10^{-1}\times 10,0}{100,0} = \SI{1,0e-2}{mol.L^{-1}}. \]
C'est une dilution au dixième ($n=10$) : le pH augmente d'une unité :
\[ \text{pH}_1 = -\log C_1 = 2,00. \]
\textbf{3. Objectif pH = 4,00.} Il faudrait une concentration $C_2 = 10^{-4,00} = \SI{1,0e-4}{mol.L^{-1}}$, soit diluer $S_1$ ($C_1 = \SI{1,0e-2}{mol.L^{-1}}$) cent fois ($n=100$). Or la fiole fait \SI{100,0}{mL} : il faudrait prélever
\[ V = \frac{C_2 V_2}{C_1} = \frac{10^{-4}\times 100,0}{10^{-2}} = \SI{1,00}{mL}, \]
ce qui est techniquement possible avec une pipette jaugée de \SI{1,00}{mL} (ou une micropipette). La concentration finale $\SI{1,0e-4}{mol.L^{-1}} > \SI{1e-6}{mol.L^{-1}}$ reste dans le domaine de validité de la formule.
\textbf{Conclusion.} $\text{pH}_0 = 1,00$ ; $\text{pH}_1 = 2,00$ ; un prélèvement de \SI{1,00}{mL} de $S_1$ dilué à \SI{100,0}{mL} donne $\text{pH} = 4,00$.
\end{exoresolu}

\begin{savaistu}[Le pH dans la vie quotidienne au Cameroun]
\begin{itemize}
\item \textbf{Vin de palme et bili-bili :} Ces boissons traditionnelles sont acides (pH $\approx$ 3,5--4,5) à cause de la fermentation lactique et acétique. Le contrôle du pH est crucial pour la conservation et le goût.
\item \textbf{Jus de bissap (foléré) :} Son acidité (pH $\approx$ 2,5--3,0) provient de l'acide hibiscus (acide organique faible) et de la vitamine C.
\item \textbf{Savonneries artisanales (région de l'Ouest, Douala) :} La saponification de l'huile de palme utilise de la soude \ce{NaOH} (base forte, pH $\approx$ 14). Le contrôle de la concentration de la soude (densimétrie ou dosage) garantit la qualité du savon et évite les excès de soude caustique irritants pour la peau.
\item \textbf{SONARA (Limbe) :} La raffinerie produit de l'acide chlorhydrique et utilise de la soude pour le traitement des eaux et la neutralisation des effluents acides.
\item \textbf{CIMENCAM (Figuil, Douala) :} Les eaux de ruissellement sur les stockages de calcaire peuvent être basiques ; leur pH est surveillé avant rejet.
\end{itemize}
\end{savaistu}

\courssub{5. Mélange d'un acide fort et d'une base forte}

\begin{methode}[Calculer le pH d'un mélange acide fort + base forte]
\begin{enumerate}
\item \textbf{Écrire la réaction} (quasi totale, $K = 1/K_e = 10^{14}$) :
\[ \ce{H3O+ + HO- -> 2 H2O} \]
\item \textbf{Calculer les quantités de matière introduites} :
$n_a = C_a V_a$ (quantité d'ions \ce{H3O+} apportée par l'acide) et $n_b = C_b V_b$ (quantité d'ions \ce{HO-} apportée par la base). \textbf{Volumes en litres !}
\item \textbf{Comparer $n_a$ et $n_b$} (tableau d'avancement recommandé) pour identifier le réactif limitant :
\begin{itemize}
\item $n_a = n_b$ : \cle{équivalence}, $\text{pH} = 7$ (à \SI{25}{\celsius}) ;
\item $n_a > n_b$ : \cle{excès d'acide}, $n(\ce{H3O+})_{\text{restant}} = n_a - n_b$ ;
\item $n_a < n_b$ : \cle{excès de base}, $n(\ce{HO-})_{\text{restant}} = n_b - n_a$.
\end{itemize}
\item \textbf{Calculer la concentration de l'espèce en excès} dans le volume total $V_{\text{total}} = V_a + V_b$ :
\begin{itemize}
\item Excès d'acide : $[\ce{H3O+}] = \dfrac{n_a - n_b}{V_a + V_b}$, puis $\text{pH} = -\log[\ce{H3O+}]$ ;
\item Excès de base : $[\ce{HO-}] = \dfrac{n_b - n_a}{V_a + V_b}$, puis $\text{pH} = pK_e + \log[\ce{HO-}]$.
\end{itemize}
\item \textbf{Contrôler la cohérence} : excès d'acide $\Rightarrow \text{pH} < 7$ ; excès de base $\Rightarrow \text{pH} > 7$.
\end{enumerate}
\end{methode}

\begin{exoresolu}[Neutralisation de l'acide d'une batterie]
On mélange $V_a = \SI{20,0}{mL}$ d'une solution d'acide chlorhydrique de concentration $C_a = \SI{1,5e-1}{mol.L^{-1}}$ avec $V_b = \SI{30,0}{mL}$ d'une solution d'hydroxyde de sodium de concentration $C_b = \SI{8,0e-2}{mol.L^{-1}}$, à \SI{25}{\celsius}.
\begin{enumerate}
\item Écrire l'équation de la réaction qui se produit.
\item Déterminer le réactif en excès à l'aide d'un tableau d'avancement.
\item Calculer le pH du mélange.
\end{enumerate}
\tcblower
\textbf{1. Équation.} Acide fort et base forte réagissent selon :
\[ \ce{H3O+ + HO- -> 2 H2O} \]
\textbf{2. Quantités initiales et tableau d'avancement.}
\[ n_a = C_a V_a = 1,5\times10^{-1}\times 20,0\times10^{-3} = \SI{3,0e-3}{mol} \]
\[ n_b = C_b V_b = 8,0\times10^{-2}\times 30,0\times10^{-3} = \SI{2,4e-3}{mol} \]
\begin{center}
\begin{tabularx}{\linewidth}{|l|X|X|X|}
\hline
\rowcolor{popblueL} État & $n(\ce{H3O+})$ (mol) & $n(\ce{HO-})$ (mol) & $n(\ce{H2O})$ \\
\hline
Initial & $3,0\times10^{-3}$ & $2,4\times10^{-3}$ & excès \\
\hline
Final ($x_{\max} = 2,4\times10^{-3}$) & $0,6\times10^{-3}$ & $\approx 0$ & excès \\
\hline
\end{tabularx}
\end{center}
Comme $n_b < n_a$, les ions \ce{HO-} sont limitants : $x_{\max} = n_b = \SI{2,4e-3}{mol}$. Il reste en excès :
\[ n(\ce{H3O+})_{\text{restant}} = n_a - n_b = 3,0\times10^{-3} - 2,4\times10^{-3} = \SI{6,0e-4}{mol}. \]
\textbf{3. pH du mélange.} Le volume total est $V_{\text{total}} = V_a + V_b = 20,0 + 30,0 = \SI{50,0}{mL} = \SI{5,00e-2}{L}$.
\[ [\ce{H3O+}] = \frac{n_a - n_b}{V_a + V_b} = \frac{6,0\times10^{-4}}{5,00\times10^{-2}} = \SI{1,2e-2}{mol.L^{-1}}. \]
\[ \text{pH} = -\log(1,2\times10^{-2}) = 1,92. \]
\textbf{Conclusion.} L'acide étant en excès, le mélange est acide : $\text{pH} = 1,92 < 7$, résultat cohérent.
\end{exoresolu}

\courssub{6. Dosage acide fort / base forte : équivalence et applications}

\begin{methode}[Déterminer le volume équivalent et la concentration inconnue]
\begin{enumerate}
\item \textbf{Équation du dosage} : \ce{H3O+ + HO- -> 2 H2O}.
\item \textbf{Condition d'équivalence} : les réactifs sont en proportions stœchiométriques :
\[ n_a = n_b \quad\Longleftrightarrow\quad C_a V_a = C_b V_{bE} \quad\Longrightarrow\quad V_{bE} = \frac{C_a V_a}{C_b} \]
\item \textbf{pH à l'équivalence} : le mélange ne contient que des ions spectateurs (\ce{Na+}, \ce{Cl-}, \ce{K+}, \ce{NO3-}...) et de l'eau : la solution est \cle{neutre}, $\text{pH} = 7$ à \SI{25}{\celsius}.
\item \textbf{Choix de l'indicateur} : sa zone de virage doit englober le pH à l'équivalence (pH 7). Exemples : bleu de bromothymol (6,0--7,6), phénolphtaléine (8,2--10,0) \emph{moins précis} car virage après l'équivalence.
\item \textbf{Exploitation} : si le volume équivalent $V_{E}$ est mesuré expérimentalement, on remonte à la concentration inconnue : $C_{\text{inconnue}} = \dfrac{C_{\text{connue}} V_{E}}{V_{\text{inconnue}}}$.
\end{enumerate}
\end{methode}

\begin{popfigure}
\begin{tikzpicture}[scale=0.8, >=Stealth]
% Support
\draw[thick, popdark] (0,0) -- (8,0);
% Burette
\draw[thick, popblue] (6.5, 0.5) -- (6.5, 6.5) -- (7.5, 6.5) -- (7.5, 0.5);
\draw[thick, popblue] (6.5, 6.5) -- (6.5, 7) node[above, popdark] {Burette};
\draw[popblue, line width=2pt] (6.5, 2) -- (7.5, 2) node[right, popdark] {$C_b, V_b$};
\draw[popblue, decorate, decoration={zigzag, segment length=2pt, amplitude=1pt}] (6.5, 1) -- (7.5, 1);
% Robinets
\draw[thick, popdark] (6.5, 0.5) -- (6.5, 0) -- (5.5, 0) -- (5.5, -0.5);
\draw[thick, popdark] (7.5, 0.5) -- (7.5, 0) -- (8.5, 0) -- (8.5, -0.5);
% Erlenmeyer
\draw[thick, poporange] (2,0) -- (1.5, 3) -- (4.5, 3) -- (4,0);
\draw[poporange, line width=2pt] (1.8, 1.2) -- (3.8, 1.2) node[right, popdark] {$C_a, V_a$};
\draw[poporange, decorate, decoration={zigzag, segment length=2pt, amplitude=1pt}] (1.8, 0.8) -- (3.8, 0.8);
% Indicator drop
\draw[poppink, fill=poppink!30] (3, 3.2) circle (0.15) node[above] {Indicateur};
% pH-meter (optional)
\draw[thick, poppurple] (0.5, 3.5) rectangle (1.5, 4.5) node[midway, align=center] {pH-mètre};
\draw[poppurple, ->] (1, 3.5) -- (1, 3);
% Labels
\node[popblue, align=center] at (7, 5) {Solution\\titrante (base)};
\node[poporange, align=center] at (3, 4.5) {Solution\\titrée (acide)};
\node[poppurple, align=center] at (1, 5) {Suivi pH};
\end{tikzpicture}
\legende{Montage pour le dosage d'un acide fort par une base forte (ou inversement) avec suivi pH-métrique ou indicateur coloré.}
\end{popfigure}

\begin{exoresolu}[Contrôle qualité d'une soude de savonnerie artisanale]
Une savonnerie artisanale de la région de l'Ouest prépare sa soude caustique pour la saponification de l'huile de palme. Pour contrôler sa concentration, un technicien dose $V_b = \SI{10,0}{mL}$ de cette solution de soude par de l'acide chlorhydrique de concentration $C_a = \SI{1,00e-1}{mol.L^{-1}}$. L'équivalence est repérée au virage du bleu de bromothymol pour $V_{aE} = \SI{12,5}{mL}$.
\begin{enumerate}
\item Définir l'équivalence et écrire la relation entre les quantités de matière.
\item Calculer la concentration $C_b$ de la soude.
\item Quel est le pH du mélange à l'équivalence ? Justifier.
\end{enumerate}
\tcblower
\textbf{1. Équivalence.} À l'équivalence, les réactifs ont été mélangés dans les proportions stœchiométriques de la réaction \ce{H3O+ + HO- -> 2 H2O} :
\[ n(\ce{H3O+})_{\text{versé}} = n(\ce{HO-})_{\text{initial}} \quad\Longleftrightarrow\quad C_a V_{aE} = C_b V_b. \]
\textbf{2. Concentration de la soude.}
\[ C_b = \frac{C_a V_{aE}}{V_b} = \frac{1,00\times10^{-1}\times 12,5}{10,0} = \SI{1,25e-1}{mol.L^{-1}}. \]
\textbf{3. pH à l'équivalence.} À l'équivalence, \ce{H3O+} et \ce{HO-} ont totalement disparu (réaction quasi totale, $K=10^{14}$). La solution ne contient que des ions spectateurs \ce{Na+} et \ce{Cl-}, sans action acido-basique sur l'eau : la solution est \cle{neutre}, donc
\[ \text{pH} = 7,00 \quad \text{à } \SI{25}{\celsius}. \]
C'est d'ailleurs pourquoi le bleu de bromothymol (zone de virage $6,0$--$7,6$) est un indicateur parfaitement adapté.
\textbf{Conclusion.} La soude artisanale a une concentration $C_b = \SI{1,25e-1}{mol.L^{-1}}$ ; le pH à l'équivalence vaut $7,00$.
\end{exoresolu}

\courssub{7. Autres acides forts, bases fortes et sécurité}

\begin{aretenir}[Acides forts et bases fortes usuels au programme]
\begin{center}
\begin{tabularx}{\linewidth}{|l|X|l|}
\hline
\rowcolor{popblueL} \textbf{Famille} & \textbf{Exemples (formule / nom)} & \textbf{Remarques} \\
\hline
Acides forts minéraux & \ce{HCl} (chlorhydrique), \ce{HBr}, \ce{HI}, \ce{HNO3} (nitrique), \ce{HClO4} (perchlorique) & Réaction totale : \ce{HA + H2O -> A- + H3O+} \\
\hline
Acide diprotique (1er proton fort) & \ce{H2SO4} (sulfurique) & \ce{H2SO4 + H2O -> HSO4- + H3O+} (total) ; \ce{HSO4-} est un acide faible ($pK_a \approx 1,9$) \\
\hline
Bases fortes (hydroxydes) & \ce{NaOH} (soude), \ce{KOH} (potasse), \ce{LiOH}, \ce{Ba(OH)2}, \ce{Sr(OH)2}, \ce{Ca(OH)2} & Dissociation totale : \ce{MOH -> M+ + HO-} ; pour \ce{Ba(OH)2} : $[\ce{HO-}] = 2C$ \\
\hline
Bases fortes organiques & Éthanolate de sodium \ce{NaOEt} (\ce{Na+ ^-OCH2CH3}) & Réaction totale avec l'eau : \ce{^-OEt + H2O -> EtOH + HO-} \\
\hline
\end{tabularx}
\end{center}
\end{aretenir}

\begin{attention}[Sécurité : pictogrammes et manipulation (norme CLP / SGH)]
\begin{itemize}
\item \textbf{Acides forts concentrés (\ce{HCl}, \ce{HNO3}, \ce{H2SO4}) et bases fortes (\ce{NaOH}, \ce{KOH}) :} \cle{Corrosif} (\text{GHS05} -- corrosion métal/peau/œil) et souvent \cle{Irritant} (\text{GHS07}).
\item \textbf{Risques :} Brûlures chimiques graves (yeux, peau, voies respiratoires), réaction exothermique violente lors de la dilution (éclaboussures).
\item \textbf{Règle d'or :} \textbf{On verse l'acide (ou la base) DANS l'eau}, jamais l'inverse. Porter lunettes de protection, gants, blouse. Manipuler sous hotte pour \ce{HCl} (vapeurs irritantes) et \ce{HNO3} (vapeurs toxiques \ce{NO2}).
\item \textbf{Premiers secours :} Rinçage immédiat et prolongé ($>$ 15 min) à grande eau. Solution de \ce{Diphoterine} si disponible. Appeler le centre antipoison (Yaoundé : 22 20 11 11 / 22 20 12 12).
\item \textbf{Élimination :} Neutralisation avant rejet (acide avec base faible type \ce{NaHCO3}, base avec acide faible type \ce{CH3COOH}), conformément au plan de gestion des déchets du lycée.
\end{itemize}
\end{attention}

\courssub{Entraînement au Baccalauréat}

\begin{exercice}[Exercice type Bac — Acide fort, dilution et mélange]
On dispose d'une solution mère $S_0$ d'acide nitrique \ce{HNO3} de concentration $C_0 = \SI{2,00e-1}{mol.L^{-1}}$.
\begin{enumerate}
\item Rappeler pourquoi \ce{HNO3} est un acide fort. Écrire l'équation de sa réaction avec l'eau.
\item Calculer le pH de $S_0$.
\item On prépare une solution $S_1$ en diluant \SI{20,0}{mL} de $S_0$ à \SI{200,0}{mL}. Calculer la concentration $C_1$ et le pH de $S_1$.
\item On mélange \SI{10,0}{mL} de $S_1$ avec \SI{15,0}{mL} d'une solution de soude \ce{NaOH} de concentration $C_b = \SI{5,0e-2}{mol.L^{-1}}$. Calculer le pH du mélange final.
\item Vérifier que la condition de validité $C > \SI{1e-6}{M}$ est respectée pour toutes les solutions utilisées.
\end{enumerate}
\end{exercice}

\begin{corrige}[Exercice type Bac]
\textbf{1.} \ce{HNO3} est un acide fort car il réagit totalement avec l'eau :
\[ \ce{HNO3 + H2O -> NO3- + H3O+} \]
\textbf{2.} $\text{pH}_0 = -\log C_0 = -\log(2,00\times10^{-1}) = 0,70$.
\textbf{3.} Dilution au dixième ($V_0=20,0$, $V_1=200,0$) : $C_1 = C_0/10 = \SI{2,00e-2}{mol.L^{-1}}$.
$\text{pH}_1 = -\log C_1 = 1,70$ (augmentation de 1 unité).
\textbf{4.} Mélange :
$n_a = C_1 V_a = 2,00\times10^{-2} \times 10,0\times10^{-3} = \SI{2,00e-4}{mol}$ (\ce{H3O+}).
$n_b = C_b V_b = 5,0\times10^{-2} \times 15,0\times10^{-3} = \SI{7,5e-4}{mol}$ (\ce{HO-}).
$n_b > n_a$ : excès de base.
$n(\ce{HO-})_{\text{restant}} = 7,5\times10^{-4} - 2,0\times10^{-4} = \SI{5,5e-4}{mol}$.
$V_{\text{tot}} = 10,0 + 15,0 = \SI{25,0}{mL} = \SI{2,50e-2}{L}$.
$[\ce{HO-}] = \frac{5,5\times10^{-4}}{2,50\times10^{-2}} = \SI{2,2e-2}{mol.L^{-1}}$.
$\text{pH} = 14 + \log(2,2\times10^{-2}) = 14 - 1,66 = 12,34$.
Cohérence : pH $> 7$ (excès de base).
\textbf{5.} $C_0 = 0,20 > 10^{-6}$ ; $C_1 = 0,020 > 10^{-6}$ ; $C_b = 0,050 > 10^{-6}$ ; $[\ce{HO-}]_{\text{final}} = 0,022 > 10^{-6}$. Condition respectée partout.
\end{corrige}

\begin{exercice}[Exercice type Bac — Dosage et concentration inconnue]
On souhaite doser une solution d'acide chlorhydrique $S_A$ de concentration inconnue $C_A$. On utilise une solution de soude $S_B$ de concentration $C_B = \SI{1,00e-1}{mol.L^{-1}}$. On verse $V_A = \SI{20,0}{mL}$ de $S_A$ dans un erlenmeyer avec quelques gouttes de bleu de bromothymol. On ajoute $S_B$ depuis une burette jusqu'au virage de l'indicateur (bleu). Le volume versé à l'équivalence est $V_{BE} = \SI{18,5}{mL}$.
\begin{enumerate}
\item Écrire l'équation de la réaction support du dosage.
\item Justifier le choix du bleu de bromothymol.
\item Calculer $C_A$.
\item On recommence le dosage avec un pH-mètre. Tracer schématiquement la courbe $\text{pH} = f(V_B)$ et indiquer les caractéristiques principales (pH initial, pH à l'équivalence, allure).
\end{enumerate}
\end{exercice}

\begin{corrige}[Exercice type Bac]
\textbf{1.} \ce{H3O+ + HO- -> 2 H2O}.
\textbf{2.} L'équivalence se fait à pH = 7. Le bleu de bromothymol a une zone de virage $[6,0 ; 7,6]$ qui contient 7. Il est donc adapté (virage jaune $\to$ bleu au passage de pH 7).
\textbf{3.} À l'équivalence : $C_A V_A = C_B V_{BE}$.
$C_A = \frac{C_B V_{BE}}{V_A} = \frac{1,00\times10^{-1} \times 18,5}{20,0} = \SI{9,25e-2}{mol.L^{-1}}$.
\textbf{4.} Courbe $\text{pH} = f(V_B)$ :
\begin{itemize}
\item Début : $V_B=0$, solution acide, $\text{pH} = -\log(9,25\times10^{-2}) \approx 1,03$.
\item Pas de zone tampon (dosage acide fort / base forte). Le pH monte très lentement au début, puis très brutalement au voisinage de l'équivalence.
\item Équivalence : $V_{BE} = \SI{18,5}{mL}$, $\text{pH} = 7,00$. Saut de pH vertical important (environ pH 4 à pH 10 pour 1 goutte $\approx \SI{0,05}{mL}$).
\item Après équivalence : excès de base, pH monte vers $\text{pH} = 14 + \log C_B \approx 13$ (asymptote).
\end{itemize}
Schéma : courbe en S très raide au point d'équivalence.
\end{corrige}

\newpage

\coursec{Exercices}

\coursec{Leçon 07 : Acides forts et bases fortes}

\courssub{L'acide chlorhydrique, modèle d'acide fort}

\begin{definition}[Acide fort]
Un \cle{acide fort} est un acide dont la réaction avec l'eau est \textbf{totale} (irréversible). Il cède la totalité de ses protons à l'eau pour former des ions oxonium \ce{H3O+}. L'équilibre est totalement déplacé vers les produits.
\end{definition}

\begin{experience}[Réaction du chlorure d'hydrogène avec l'eau]
\begin{popfigure}
\begin{tikzpicture}[scale=0.9, every node/.style={font=\small}]
% Gaz HCl
\draw[popblue, thick] (0,0) rectangle (3,2);
\node[popblue] at (1.5,1.5) {\textbf{HCl (gaz)}};
\node[popblue] at (1.5,1.0) {Molécule covalente};
\node[popblue] at (1.5,0.5) {Soluble dans l'eau};
% Flèche dissolution
\draw[->, thick, popdark] (3,1) -- (4,1) node[midway, above] {Dissolution};
% Solution aqueuse
\draw[poporange, thick] (4,0) rectangle (8,2);
\node[poporange] at (6,1.5) {\textbf{Solution aqueuse}};
% Ions
\node[poporange] at (5.5,1.0) {\ce{H3O+} (ion oxonium)};
\node[poporange] at (6.5,1.0) {\ce{Cl-} (ion chlorure)};
\node[poporange] at (6,0.5) {Conduction électrique};
\end{tikzpicture}
\legende{Dissolution du chlorure d'hydrogène gazeux dans l'eau : formation totale des ions \ce{H3O+} et \ce{Cl-}.}
\end{popfigure}

\textbf{Équation de la réaction (totale) :}
\[ \ce{HCl(g) + H2O(l) -> H3O+(aq) + Cl-(aq)} \]
\textbf{Couple acide/base :} \ce{HCl}/\ce{Cl-} (acide fort / base très faible) et \ce{H3O+}/\ce{H2O}.
\end{experience}

\begin{propriete}[pH d'une solution d'acide fort]
Pour une solution d'acide fort de concentration molaire $C$ (en \si{mol.L^{-1}}), la concentration en ions oxonium est :
\[ [\ce{H3O+}] = C \]
Le pH s'exprime alors simplement par :
\[ \mathrm{pH} = -\log C \]
Cette relation est valable pour \textbf{$C > \SI{10^{-6}}{mol.L^{-1}}$}. En dessous, l'autoprotolyse de l'eau devient prépondérante.
\end{propriete}

\begin{aretenir}[Exemples d'acides forts usuels]
\begin{itemize}
\item Acide chlorhydrique \ce{HCl} (gaz dissous, commerce : \SI{30}{\%} à \SI{35}{\%} massique, $\approx \SI{12}{mol.L^{-1}}$).
\item Acide nitrique \ce{HNO3} (liquide, commerce : \SI{65}{\%} massique, $\approx \SI{14}{mol.L^{-1}}$).
\item Acide sulfurique \ce{H2SO4} (premier proton seulement, \ce{H2SO4 -> H3O+ + HSO4-}).
\item Acide perchlorique \ce{HClO4}.
\end{itemize}
\end{aretenir}

\courssub{L'hydroxyde de sodium, modèle de base forte}

\begin{definition}[Base forte]
Une \cle{base forte} est une base qui se dissocie \textbf{totalement} dans l'eau en libérant des ions hydroxyde \ce{HO-}. C'est généralement un solide ionique (hydroxyde de métal alcalin ou alcalino-terreux) ou un sel d'alcoolate.
\end{definition}

\begin{experience}[Dissolution de l'hydroxyde de sodium]
\begin{popfigure}
\begin{tikzpicture}[scale=0.9, every node/.style={font=\small}]
% Solide NaOH
\draw[poppurple, thick] (0,0) rectangle (3,2);
\node[poppurple] at (1.5,1.5) {\textbf{NaOH (solide)}};
\node[poppurple] at (1.5,1.0) {Réseau ionique \ce{Na+} \ce{HO-}};
\node[poppurple] at (1.5,0.5) {Hygiroscopique, corrosif};
% Flèche dissolution
\draw[->, thick, popdark] (3,1) -- (4,1) node[midway, above] {Dissolution};
% Solution aqueuse
\draw[popgreen, thick] (4,0) rectangle (8,2);
\node[popgreen] at (6,1.5) {\textbf{Solution aqueuse (Soude)}};
% Ions
\node[popgreen] at (5.5,1.0) {\ce{Na+} (ion sodium)};
\node[popgreen] at (6.5,1.0) {\ce{HO-} (ion hydroxyde)};
\node[popgreen] at (6,0.5) {Conduction électrique, pH > 7};
\end{tikzpicture}
\legende{Dissociation totale de l'hydroxyde de sodium solide dans l'eau.}
\end{popfigure}

\textbf{Équation de la dissociation (totale) :}
\[ \ce{NaOH(s) ->[\text{H2O}] Na+(aq) + HO-(aq)} \]
La concentration en ions hydroxyde est égale à la concentration de la base :
\[ [\ce{HO-}] = C \]
\end{experience}

\begin{propriete}[pH d'une solution de base forte]
À \SI{25}{\celsius}, le produit ionique de l'eau est $\mathrm{Ke} = [\ce{H3O+}][\ce{HO-}] = \SI{1,0e-14}{}$.
Pour une base forte de concentration $C$ ($C > \SI{10^{-6}}{mol.L^{-1}}$) :
\[ [\ce{HO-}] = C \quad \Rightarrow \quad [\ce{H3O+}] = \frac{\mathrm{Ke}}{C} \]
\[ \mathrm{pH} = -\log\left(\frac{\mathrm{Ke}}{C}\right) = \mathrm{pKe} + \log C \]
Avec $\mathrm{pKe} = 14$ à \SI{25}{\celsius} : $\mathrm{pH} = 14 + \log C$.
\end{propriete}

\begin{aretenir}[Exemples de bases fortes usuelles]
\begin{itemize}
\item Hydroxyde de sodium \ce{NaOH} (soude, lessive, savonneries artisanales au Cameroun).
\item Hydroxyde de potassium \ce{KOH} (potasse, piles alcalines).
\item Hydroxyde de baryum \ce{Ba(OH)2} (dibasique, $[\ce{HO-}] = 2C$).
\item Éthanolate de sodium \ce{C2H5ONa} (base forte non aqueuse, utilisée en synthèse organique).
\end{itemize}
\end{aretenir}

\courssub{Espèces chimiques en solution : inventaire et calcul des concentrations}

\begin{methode}[Inventaire et concentrations dans une solution d'acide ou base fort(e)]
\begin{enumerate}
\item \textbf{Écrire l'équation de la réaction totale} (acide + eau ou dissociation base).
\item \textbf{Lister les espèces} : ions issus de la réaction totale + \ce{H2O} + \ce{H3O+} + \ce{HO-} (issus de l'autoprotolyse).
\item \textbf{Appliquer la conservation de la matière} : pour un acide fort $[\ce{A-}] = C$ ; pour une base forte $[\ce{M+}] = C$ (ou $2C$ pour \ce{Ba(OH)2}).
\item \textbf{Calculer $[\ce{H3O+}]$ et $[\ce{HO-}]$} via $\mathrm{pH} = -\log[\ce{H3O+}]$ et $\mathrm{Ke} = [\ce{H3O+}][\ce{HO-}]$.
\item \textbf{Vérifier l'électroneutralité} : $\sum [\text{cations}] = \sum [\text{anions}]$.
\end{enumerate}
\end{methode}

\begin{exemplebox}[Solution mixte NaCl + HCl]
On dissout \ce{NaCl} ($C_1$) et \ce{HCl} ($C_2$) dans l'eau.
\textbf{Inventaire :} \ce{Na+}, \ce{Cl-}, \ce{H3O+}, \ce{HO-}, \ce{H2O}.
\textbf{Concentrations :} $[\ce{Na+}] = C_1$ ; $[\ce{Cl-}] = C_1 + C_2$ ; $[\ce{H3O+}] = C_2$ ; $[\ce{HO-}] = \mathrm{Ke}/C_2$.
\textbf{Électroneutralité :} $[\ce{Na+}] + [\ce{H3O+}] = C_1 + C_2 = [\ce{Cl-}] + [\ce{HO-}] \approx [\ce{Cl-}]$.
\end{exemplebox}

\courssub{Effet de la dilution sur le pH}

\begin{propriete}[Variation du pH par dilution au dixième]
Soit une solution d'acide fort de concentration $C$ et pH $= -\log C$.
Après dilution au dixième ($C' = C/10$) :
\[ \mathrm{pH}' = -\log(C/10) = -\log C + \log 10 = \mathrm{pH} + 1 \]
\textbf{Le pH augmente de 1 unité.}

Soit une solution de base forte de concentration $C$ et pH $= \mathrm{pKe} + \log C$.
Après dilution au dixième ($C' = C/10$) :
\[ \mathrm{pH}' = \mathrm{pKe} + \log(C/10) = \mathrm{pKe} + \log C - 1 = \mathrm{pH} - 1 \]
\textbf{Le pH diminue de 1 unité.}
\end{propriete}

\begin{attention}[Limites de validité : la dilution extrême]
La relation $\mathrm{pH} = -\log C$ (acide) ou $\mathrm{pH} = \mathrm{pKe} + \log C$ (base) suppose que la contribution de l'autoprotolyse de l'eau est négligeable.
\textbf{Condition :} $C > \SI{10^{-6}}{mol.L^{-1}}$.
Si $C \le \SI{10^{-6}}{mol.L^{-1}}$, le pH tend vers 7 (neutre) et non vers des valeurs extrêmes.
\begin{itemize}
\item Acide fort $C = \SI{10^{-8}}{M}$ : $[\ce{H3O+}] \approx \SI{10^{-7}}{M}$ (eau), $\mathrm{pH} \approx 7$ (légèrement $< 7$).
\item Base forte $C = \SI{10^{-8}}{M}$ : $[\ce{HO-}] \approx \SI{10^{-7}}{M}$, $\mathrm{pH} \approx 7$ (légèrement $> 7$).
\end{itemize}
\end{attention}

\begin{popfigure}
\begin{tikzpicture}
\begin{axis}[
    width=\linewidth, height=6cm,
    xlabel={$\log_{10}(C/\mathrm{mol.L^{-1}})$}, ylabel={pH},
    xmin=-8, xmax=1, ymin=0, ymax=14,
    grid=major, grid style={popline},
    legend pos=south east,
    tick label style={font=\small}, label style={font=\small},
]
% Acide fort théorique
\addplot[popblue, thick, domain=-6:1, samples=50] {-x} node[pos=0.9, right] {Acide fort (théorique)};
% Acide fort réel (avec eau)
\addplot[popblue, dashed, thick, domain=-8:-6, samples=50] {7 - 10^(x+7)*3} node[pos=0.1, left] {Réalité (eau)};
% Base forte théorique
\addplot[poporange, thick, domain=-6:1, samples=50] {14 + x} node[pos=0.9, right] {Base forte (théorique)};
% Base forte réel
\addplot[poporange, dashed, thick, domain=-8:-6, samples=50] {7 + 10^(x+7)*3} node[pos=0.1, left] {Réalité (eau)};
% Neutre
\addplot[popmuted, dashed, domain=-8:1] {7} node[pos=0.5, above] {Neutre};
\end{axis}
\end{tikzpicture}
\legende{Évolution du pH en fonction de $\log C$ pour un acide fort et une base forte. Les pointillés montrent l'écart à la loi simple pour $C < 10^{-6}$ M.}
\end{popfigure}

\courssub{Mélange d'un acide fort et d'une base forte}

\begin{propriete}[Réaction de neutralisation]
Le mélange d'un acide fort (source de \ce{H3O+}) et d'une base forte (source de \ce{HO-}) entraîne la réaction :
\[ \ce{H3O+(aq) + HO-(aq) -> 2 H2O(l)} \]
Cette réaction est \textbf{totale} car sa constante d'équilibre est $K = 1/\mathrm{Ke} = \SI{10^{14}}{}$ (à \SI{25}{\celsius}). L'avancement maximal est atteint.
\end{propriete}

\begin{methode}[Calcul du pH d'un mélange acide fort / base forte]
\begin{enumerate}
\item Calculer les quantités de matière initiales : $n_{\text{acide}} = C_A V_A$, $n_{\text{base}} = C_B V_B$.
\item Identifier le réactif en excès via le tableau d'avancement (réaction totale, $x_{\max} = \min(n_{\text{acide}}, n_{\text{base}})$).
\item Calculer la concentration de l'ion en excès dans le volume total $V_{\text{tot}} = V_A + V_B$.
\item En déduire le pH : $\mathrm{pH} = -\log[\ce{H3O+}]$ (excès acide) ou $\mathrm{pH} = 14 + \log[\ce{HO-}]$ (excès base).
\item Si $n_{\text{acide}} = n_{\text{base}}$ (équivalence), le mélange est neutre : $\mathrm{pH} = 7$.
\end{enumerate}
\end{methode}

\courssub{Autres acides forts et bases fortes usuels ; sécurité}

\begin{aretenir}[Tableau récapitulatif des espèces fortes]
\begin{center}
\begin{tabularx}{\linewidth}{|l|X|l|X|}\hline
\rowcolor{popblueL}
\textbf{Espèce} & \textbf{Formule} & \textbf{Type} & \textbf{Équation dans l'eau} \\\hline
Chlorure d'hydrogène & \ce{HCl} & Acide fort & $\ce{HCl + H2O -> H3O+ + Cl-}$ \\\hline
Acide nitrique & \ce{HNO3} & Acide fort & $\ce{HNO3 + H2O -> H3O+ + NO3-}$ \\\hline
Acide sulfurique (1er H+) & \ce{H2SO4} & Acide fort & $\ce{H2SO4 + H2O -> H3O+ + HSO4-}$ \\\hline
Hydroxyde de sodium & \ce{NaOH} & Base forte & $\ce{NaOH -> Na+ + HO-}$ \\\hline
Hydroxyde de potassium & \ce{KOH} & Base forte & $\ce{KOH -> K+ + HO-}$ \\\hline
Éthanolate de sodium & \ce{C2H5ONa} & Base forte & $\ce{C2H5ONa -> C2H5O- + Na+}$ \\\hline
\end{tabularx}
\end{center}
\end{aretenir}

\begin{attention}[Sécurité : pictogrammes et règles (norme CLP/GHS)]
\begin{popfigure}
\begin{tikzpicture}[scale=0.8, every node/.style={font=\small}]
% GHS05 Corrosif
\begin{scope}[xshift=0cm]
\draw[thick, poppink] (0,0) circle (1.5);
\draw[poppink, thick] (-0.5,0.5) -- (0.5,-0.5);
\draw[poppink, thick] (-0.5,-0.5) -- (0.5,0.5);
\draw[poppink, thick] (0,-1.2) -- (0,1.2);
\node at (0,-2) {\textbf{GHS05} : Corrosif};
\node at (0,-2.5) {Atteint les métaux, peau, yeux};
\end{scope}
% GHS07 Danger (irritant)
\begin{scope}[xshift=4.5cm]
\draw[thick, poppink] (0,0) circle (1.5);
\draw[poppink, thick] (0,0) circle (0.5);
\node at (0,-2) {\textbf{GHS07} : Danger};
\node at (0,-2.5) {Irritant, nocif, sensibilisant};
\end{scope}
% GHS02 Inflammable (éthanolate)
\begin{scope}[xshift=9cm]
\draw[thick, poppink] (0,0) circle (1.5);
\draw[poppink, thick, fill=poppink] (-0.3,-0.5) -- (0.3,-0.5) -- (0,0.5) -- cycle;
\node at (0,-2) {\textbf{GHS02} : Inflammable};
\node at (0,-2.5) {Solide/liquide inflammable};
\end{scope}
\end{tikzpicture}
\legende{Principaux pictogrammes rencontrés sur les flacons d'acides/bases forts (HCl, HNO3, NaOH, C2H5ONa).}
\end{popfigure}

\textbf{Règles de sécurité au laboratoire (contexte camerounais) :}
\begin{enumerate}
\item Porter \textbf{obligatoirement} lunettes de protection, blouse en coton, gants (nitrile pour bases, néoprène pour acides concentrés).
\item \textbf{Jamais} verser l'eau dans l'acide concentré (éclaboussures violentes, chaleur) : \emph{« L'acide dans l'eau, c'est la règle qui sauve »}.
\item Manipuler les solutions concentrées (\ce{HCl} \SI{12}{M}, \ce{HNO3} \SI{14}{M}, \ce{NaOH} solide) sous \textbf{hotte aspirante} (vapeurs irritantes/corrosives).
\item En cas de projection : rincer \textbf{immédiatement} et \textbf{longuement} (15 min minimum) à l'eau courante (douche de sécurité, lava-yeux). Avertir l'enseignant.
\item Étiqueter toute solution préparée (nature, concentration, date, danger).
\end{enumerate}
\end{attention}

\begin{savaistu}[Le saviez-vous ? : La savonnerie artisanale au Cameroun]
La fabrication traditionnelle du savon (à Bafoussam, Foumban, ou dans les quartiers de Douala/Yaoundé) utilise la \textbf{saponification} : réaction de l'huile de palme (triglycérides) avec une solution concentrée de \textbf{soude} (\ce{NaOH}) ou de \textbf{potasse} (\ce{KOH}, issue des cendres de bois de palmier ou de bananier). La soude/potasse est une base forte corrosive (GHS05). Les artisans manipulent ces solutions très concentrées (souvent $> \SI{5}{mol.L^{-1}}$) à chaud. La maîtrise du pH (excès de base pour assurer la réaction totale, puis neutralisation/lavage) est cruciale pour obtenir un savon doux (pH final $\approx 8-10$) et non caustique. C'est une application millénaire de la chimie des bases fortes !
\end{savaistu}

\courssub{Je vérifie mes connaissances}

\begin{exercice}[QCM : acides et bases forts]
Pour chaque question, une seule réponse est exacte. Recopie la lettre correspondante.
\begin{enumerate}
\item Une solution d'acide chlorhydrique a une concentration $C = \SI{1,0e-3}{mol.L^{-1}}$. Son pH vaut :
\begin{enumerate}
\item[a)] $3$ \quad b) $11$ \quad c) $2$ \quad d) $10^{-3}$
\end{enumerate}
\item Une solution d'hydroxyde de sodium a un $\mathrm{pH} = 12$. Sa concentration en ions hydroxyde vaut :
\begin{enumerate}
\item[a)] $\SI{1,0e-12}{mol.L^{-1}}$ \quad b) $\SI{1,0e-2}{mol.L^{-1}}$ \quad c) $\SI{12}{mol.L^{-1}}$ \quad d) $\SI{1,0e-1}{mol.L^{-1}}$
\end{enumerate}
\item On dilue dix fois une solution d'acide chlorhydrique de $\mathrm{pH} = 2$. Le pH de la solution fille vaut :
\begin{enumerate}
\item[a)] $1$ \quad b) $2$ \quad c) $3$ \quad d) $12$
\end{enumerate}
\item Dans une solution d'acide chlorhydrique de concentration $C$, la concentration en ions chlorure vaut :
\begin{enumerate}
\item[a)] $[\ce{Cl-}] = C/2$ \quad b) $[\ce{Cl-}] = C$ \quad c) $[\ce{Cl-}] = 2C$ \quad d) $[\ce{Cl-}] = \mathrm{pKe}/C$
\end{enumerate}
\end{enumerate}
\end{exercice}

\begin{exercice}[Vrai ou faux, avec justification]
Pour chaque affirmation, indique si elle est \textbf{vraie} ou \textbf{fausse}, puis justifie ta réponse en une phrase ou par un calcul.
\begin{enumerate}
\item Un acide fort est un acide dont la réaction avec l'eau est totale.
\item Le pH d'une solution d'acide fort est toujours égal à $-\log C$, quelle que soit sa concentration.
\item Une solution de soude à $\SI{1,0e-4}{mol.L^{-1}}$ a un pH égal à $4$.
\item Dans une solution d'hydroxyde de sodium, on a toujours $[\ce{Na+}] = [\ce{HO-}]$.
\item Une solution d'acide chlorhydrique de concentration $C = \SI{1,0e-8}{mol.L^{-1}}$ a un pH égal à $8$.
\item L'éthanolate de sodium \ce{C2H5ONa} est une base forte.
\end{enumerate}
\end{exercice}

\begin{exercice}[Texte à trous]
Complète le texte suivant avec les mots ou expressions de la liste ci-dessous (certains mots peuvent être utilisés plusieurs fois) : \emph{totale}, \emph{partielle}, \emph{ion oxonium}, \emph{ion hydroxyde}, \emph{$-\log C$}, \emph{$\mathrm{pKe} + \log C$}, \emph{augmente}, \emph{diminue}, \emph{électroneutralité}.

Le chlorure d'hydrogène \ce{HCl} réagit avec l'eau selon une réaction \ldots\ldots\ldots : c'est donc un acide \ldots\ldots\ldots. La solution obtenue contient des ions \ldots\ldots\ldots \ce{H3O+} et des ions chlorure \ce{Cl-}. Le pH d'une telle solution de concentration $C$ (avec $C > \SI{e-6}{mol.L^{-1}}$) est donné par la relation pH = \ldots\ldots\ldots.

L'hydroxyde de sodium \ce{NaOH} se dissocie totalement dans l'eau en ions sodium \ce{Na+} et ions \ldots\ldots\ldots \ce{HO-} : c'est une base \ldots\ldots\ldots. Le pH de sa solution est donné par pH = \ldots\ldots\ldots.

Quand on dilue dix fois une solution d'acide fort, son pH \ldots\ldots\ldots d'une unité. Quand on dilue dix fois une solution de base forte, son pH \ldots\ldots\ldots d'une unité. Dans toute solution aqueuse, la somme des charges positives est égale à la somme des charges négatives : c'est la loi d'\ldots\ldots\ldots.
\end{exercice}

\begin{exercice}[Définitions et formules essentielles]
\begin{enumerate}
\item Définis un acide fort et donne deux exemples avec leurs formules.
\item Définis une base forte et donne deux exemples avec leurs formules.
\item Écris l'équation de la réaction du chlorure d'hydrogène avec l'eau, en précisant le couple acide/base mis en jeu.
\item Rappelle la relation entre $[\ce{H3O+}]$ et $[\ce{HO-}]$ dans toute solution aqueuse à \SI{25}{\celsius}, ainsi que la valeur du produit ionique de l'eau $\mathrm{Ke}$.
\end{enumerate}
\end{exercice}

\courssub{Je m'entraîne}

\begin{exoresolu}[Solution d'acide nitrique]
Une solution d'acide nitrique \ce{HNO3}, acide fort, a une concentration $C = \SI{2,0e-3}{mol.L^{-1}}$.
\begin{enumerate}
\item Écris l'équation de la réaction de l'acide nitrique avec l'eau.
\item Calcule le pH de cette solution.
\item On dilue cette solution $10$ fois, puis $100$ fois par rapport à la solution initiale. Donne le pH de chacune des deux solutions filles.
\item Que se passerait-il si on diluait cette solution $10^{7}$ fois ? Le pH vaudrait-il $9{,}7$ ? Justifie.
\end{enumerate}
\tcblower
\textbf{Solution détaillée :}
\begin{enumerate}
\item \ce{HNO3(aq) + H2O(l) -> H3O+(aq) + NO3-(aq)} (Réaction totale, acide fort).
\item $[\ce{H3O+}] = C = \SI{2,0e-3}{mol.L^{-1}}$. $\mathrm{pH} = -\log(2,0 \times 10^{-3}) = 3 - \log 2,0 = 3 - 0,30 = \mathbf{2,70}$.
\item Dilution 10 fois : $C_1 = \SI{2,0e-4}{mol.L^{-1}}$. $\mathrm{pH}_1 = -\log(2,0 \times 10^{-4}) = 4 - 0,30 = \mathbf{3,70}$ (augmentation de 1).
Dilution 100 fois : $C_2 = \SI{2,0e-5}{mol.L^{-1}}$. $\mathrm{pH}_2 = -\log(2,0 \times 10^{-5}) = 5 - 0,30 = \mathbf{4,70}$ (augmentation de 2).
\item Dilution $10^7$ fois : $C_3 = \SI{2,0e-10}{mol.L^{-1}} < \SI{10^{-6}}{mol.L^{-1}}$. La contribution de l'autoprotolyse de l'eau ($[\ce{H3O+}]_{\text{eau}} = \SI{10^{-7}}{M}$) devient majoritaire. Le pH sera proche de \textbf{7} (légèrement acide, $\approx 6,9$), \textbf{pas 9,7}. La formule $\mathrm{pH} = -\log C$ n'est plus valide.
\end{enumerate}
\end{exoresolu}

\begin{exoresolu}[Préparation d'une solution de soude]
On dissout $\SI{0,40}{g}$ d'hydroxyde de sodium \ce{NaOH} dans de l'eau distillée de façon à obtenir $\SI{500}{mL}$ de solution. On donne $M(\ce{NaOH}) = \SI{40}{g.mol^{-1}}$.
\begin{enumerate}
\item Écris l'équation de la dissolution de l'hydroxyde de sodium dans l'eau.
\item Calcule la quantité de matière de soude dissoute, puis la concentration $C$ de la solution.
\item Calcule le pH de la solution obtenue à \SI{25}{\celsius}.
\item On prélève $\SI{10}{mL}$ de cette solution et on ajoute de l'eau jusqu'à obtenir $\SI{100}{mL}$. Quel est le pH de la solution diluée ?
\end{enumerate}
\tcblower
\textbf{Solution détaillée :}
\begin{enumerate}
\item $\ce{NaOH(s) ->[\text{H2O}] Na+(aq) + HO-(aq)}$ (Dissociation totale).
\item $n(\ce{NaOH}) = \frac{m}{M} = \frac{0,40}{40} = \SI{0,010}{mol} = \SI{1,0e-2}{mol}$.
$V = \SI{500}{mL} = \SI{0,500}{L}$.
$C = \frac{n}{V} = \frac{1,0 \times 10^{-2}}{0,500} = \SI{2,0e-2}{mol.L^{-1}}$.
\item $[\ce{HO-}] = C = \SI{2,0e-2}{M}$.
$\mathrm{pOH} = -\log(2,0 \times 10^{-2}) = 2 - 0,30 = 1,70$.
$\mathrm{pH} = 14 - \mathrm{pOH} = 14 - 1,70 = \mathbf{12,30}$.
\item Dilution 10 fois ($V$ passe de 10 à 100 mL) : $C' = \SI{2,0e-3}{M}$.
$\mathrm{pOH}' = -\log(2,0 \times 10^{-3}) = 3 - 0,30 = 2,70$.
$\mathrm{pH}' = 14 - 2,70 = \mathbf{11,30}$ (diminution de 1 unité).
\end{enumerate}
\end{exoresolu}

\begin{exoresolu}[Dilution d'une soude concentrée]
On dispose d'une solution de soude (hydroxyde de sodium) de $\mathrm{pH} = 12$.
\begin{enumerate}
\item Calcule la concentration $C$ de cette solution.
\item On prélève un volume $V_0 = \SI{50}{mL}$ de cette solution. Quel volume d'eau distillée faut-il ajouter pour obtenir une solution de $\mathrm{pH} = 11$ ?
\item Quel volume d'eau faudrait-il ajouter aux mêmes $\SI{50}{mL}$ pour obtenir un pH égal à $10$ ?
\end{enumerate}
\tcblower
\textbf{Solution détaillée :}
\begin{enumerate}
\item $\mathrm{pH} = 12 \Rightarrow \mathrm{pOH} = 2 \Rightarrow [\ce{HO-}] = 10^{-2} = \SI{1,0e-2}{mol.L^{-1}}$. Donc $C = \SI{1,0e-2}{mol.L^{-1}}$.
\item Cible $\mathrm{pH} = 11 \Rightarrow \mathrm{pOH} = 3 \Rightarrow [\ce{HO-}]' = \SI{1,0e-3}{mol.L^{-1}}$.
Dilution au dixième ($C' = C/10$). $V_f = 10 \times V_0 = \SI{500}{mL}$.
Volume d'eau à ajouter : $V_{\text{eau}} = V_f - V_0 = 500 - 50 = \mathbf{\SI{450}{mL}}$.
\item Cible $\mathrm{pH} = 10 \Rightarrow \mathrm{pOH} = 4 \Rightarrow [\ce{HO-}]'' = \SI{1,0e-4}{mol.L^{-1}}$.
Dilution au centième ($C'' = C/100$). $V_f = 100 \times V_0 = \SI{5000}{mL} = \SI{5,0}{L}$.
Volume d'eau à ajouter : $V_{\text{eau}} = 5000 - 50 = \mathbf{\SI{4950}{mL}}$ (soit \SI{4,95}{L}).
\end{enumerate}
\end{exoresolu}

\begin{exoresolu}[Inventaire des espèces en solution]
On prépare $\SI{1,0}{L}$ de solution en dissolvant $\SI{5,85}{g}$ de chlorure de sodium \ce{NaCl} et $\SI{0,365}{g}$ de chlorure d'hydrogène \ce{HCl} dans l'eau. On donne : $M(\ce{NaCl}) = \SI{58,5}{g.mol^{-1}}$ et $M(\ce{HCl}) = \SI{36,5}{g.mol^{-1}}$.
\begin{enumerate}
\item Fais l'inventaire de toutes les espèces chimiques présentes dans la solution (ions et molécules).
\item Calcule la concentration de chaque ion présent en quantité notable.
\item Vérifie que la loi d'électroneutralité est satisfaite.
\item Calcule le pH de la solution.
\end{enumerate}
\tcblower
\textbf{Solution détaillée :}
\begin{enumerate}
\item \textbf{Inventaire :} \ce{Na+}, \ce{Cl-}, \ce{H3O+}, \ce{HO-}, \ce{H2O}. (Le \ce{HCl} et \ce{NaCl} sont totalement dissociés/ionisés).
\item Quantités de matière :
$n(\ce{NaCl}) = 5,85 / 58,5 = \SI{0,100}{mol} \Rightarrow [\ce{Na+}] = \SI{0,100}{M}$, $[\ce{Cl-}]_{\text{NaCl}} = \SI{0,100}{M}$.
$n(\ce{HCl}) = 0,365 / 36,5 = \SI{0,0100}{mol} \Rightarrow [\ce{H3O+}] = \SI{0,0100}{M}$, $[\ce{Cl-}]_{\text{HCl}} = \SI{0,0100}{M}$.
Concentrations totales :
$[\ce{Na+}] = \SI{0,100}{M}$.
$[\ce{Cl-}] = 0,100 + 0,0100 = \SI{0,110}{M}$.
$[\ce{H3O+}] = \SI{0,0100}{M}$.
$[\ce{HO-}] = \frac{\mathrm{Ke}}{[\ce{H3O+}]} = \frac{10^{-14}}{10^{-2}} = \SI{1,0e-12}{M}$ (négligeable).
\item \textbf{Électroneutralité :}
$\sum [\text{cations}] = [\ce{Na+}] + [\ce{H3O+}] = 0,100 + 0,0100 = \SI{0,110}{M}$.
$\sum [\text{anions}] = [\ce{Cl-}] + [\ce{HO-}] = 0,110 + 10^{-12} \approx \SI{0,110}{M}$.
$\Rightarrow$ \textbf{Électroneutralité vérifiée.}
\item $\mathrm{pH} = -\log[\ce{H3O+}] = -\log(1,0 \times 10^{-2}) = \mathbf{2,00}$.
\end{enumerate}
\end{exoresolu}

\courssub{Je me prépare au Bac}

\begin{exoresolu}[Type Bac : l'acide chlorhydrique du laboratoire]
Au laboratoire d'un lycée de Yaoundé, on trouve un flacon étiqueté « solution d'acide chlorhydrique ». Pour vérifier la nature de cet acide, on mesure le pH de la solution avec un pH-mètre étalonné : on lit $\mathrm{pH} = 2{,}7$ à \SI{25}{\celsius}. On donne $\mathrm{pKe} = 14$ à \SI{25}{\celsius}.
\begin{enumerate}
\item Calcule la concentration en ions oxonium $[\ce{H3O+}]$ de la solution.
\item En supposant que l'acide chlorhydrique est un acide fort, écris l'équation de sa réaction avec l'eau et déduis-en la concentration $C$ de la solution.
\item Fais l'inventaire complet des espèces chimiques présentes dans la solution.
\item Calcule les concentrations de toutes les espèces ioniques présentes. Justifie que $[\ce{HO-}]$ est négligeable devant $[\ce{H3O+}]$.
\item Vérifie la loi d'électroneutralité de la solution.
\item On souhaite maintenant démontrer que \ce{HCl} est bien un acide fort. Explique, à partir de la comparaison entre $[\ce{H3O+}]$ et $C$, pourquoi la réaction de \ce{HCl} avec l'eau est totale.
\item On dilue $100$ fois la solution initiale. Quel est le pH de la solution fille ? Justifie ta réponse.
\end{enumerate}
\tcblower
\textbf{Solution détaillée :}
\begin{enumerate}
\item $[\ce{H3O+}] = 10^{-\mathrm{pH}} = 10^{-2,7} = \mathbf{\SI{2,0e-3}{mol.L^{-1}}}$ (arrondi à 2 chiffres significatifs).
\item \textbf{Équation :} $\ce{HCl(aq) + H2O(l) -> H3O+(aq) + Cl-(aq)}$.
Comme la réaction est totale (hypothèse acide fort), $[\ce{H3O+}] = C$.
Donc $C = \mathbf{\SI{2,0e-3}{mol.L^{-1}}}$.
\item \textbf{Inventaire :} \ce{H3O+}, \ce{Cl-}, \ce{HO-}, \ce{H2O}. (Pas de \ce{HCl} moléculaire, réaction totale).
\item $[\ce{H3O+}] = \SI{2,0e-3}{M}$.
$[\ce{Cl-}] = C = \SI{2,0e-3}{M}$.
$[\ce{HO-}] = \frac{\mathrm{Ke}}{[\ce{H3O+}]} = \frac{10^{-14}}{2,0 \times 10^{-3}} = \mathbf{\SI{5,0e-12}{M}}$.
$[\ce{HO-}] / [\ce{H3O+}] = 2,5 \times 10^{-9} \ll 1 \Rightarrow$ \textbf{négligeable}.
\item $\sum [+] = [\ce{H3O+}] = \SI{2,0e-3}{M}$.
$\sum [-] = [\ce{Cl-}] + [\ce{HO-}] = 2,0 \times 10^{-3} + 5,0 \times 10^{-12} \approx \SI{2,0e-3}{M}$.
$\Rightarrow$ \textbf{Électroneutralité vérifiée.}
\item On a trouvé $[\ce{H3O+}] = C$. Cela signifie que \textbf{toutes} les molécules de \ce{HCl} introduites ont réagi pour donner \ce{H3O+}. L'avancement de la réaction est maximal ($x_{\max} = C$). La réaction est donc \textbf{totale} : \ce{HCl} est un acide fort.
\item Dilution 100 fois : $C' = C/100 = \SI{2,0e-5}{M}$.
$C' > \SI{10^{-6}}{M}$, la formule $\mathrm{pH} = -\log C'$ reste valide.
$\mathrm{pH}' = -\log(2,0 \times 10^{-5}) = 5 - \log 2,0 = 5 - 0,30 = \mathbf{4,70}$.
(On peut aussi dire : dilution 100 fois = 2 dilutions successives au dixième $\Rightarrow$ pH augmente de 2 unités : $2,7 + 2 = 4,7$).
\end{enumerate}
\end{exoresolu}

\begin{exoresolu}[Type Bac : mélange acide chlorhydrique -- soude]
Lors d'une séance de travaux pratiques, un élève de Terminale C mélange un volume $V_A = \SI{25}{mL}$ d'une solution d'acide chlorhydrique de concentration $C_A = \SI{0,08}{mol.L^{-1}}$ avec un volume $V_B = \SI{15}{mL}$ d'une solution d'hydroxyde de sodium de concentration $C_B = \SI{0,10}{mol.L^{-1}}$. La température est de \SI{25}{\celsius} et $\mathrm{pKe} = 14$.
\begin{enumerate}
\item Écris l'équation de la réaction qui se produit lors du mélange. Pourquoi peut-on la considérer comme totale ?
\item Calcule les quantités de matière initiales $n(\ce{H3O+})$ et $n(\ce{HO-})$ introduites dans le mélange.
\item Construis le tableau d'avancement de la réaction et détermine l'avancement maximal $x_{max}$ ainsi que le réactif en excès.
\item Calcule la concentration de l'ion en excès dans le mélange (volume total supposé égal à $V_A + V_B$).
\item Déduis-en le pH du mélange. Le mélange est-il acide, neutre ou basique ?
\item Quel volume de soude à $\SI{0,10}{mol.L^{-1}}$ aurait-il fallu verser dans les $\SI{25}{mL}$ d'acide pour obtenir un mélange de pH égal à $7$ ?
\end{enumerate}
\tcblower
\textbf{Solution détaillée :}
\begin{enumerate}
\item \textbf{Équation :} $\ce{H3O+(aq) + HO-(aq) -> 2 H2O(l)}$.
Cette réaction est l'inverse de l'autoprotolyse de l'eau. Sa constante d'équilibre est $K = 1/\mathrm{Ke} = 10^{14}$. Comme $K \gg 1$, la réaction est \textbf{totale} (quasi-complète dans le sens direct).
\item $n(\ce{H3O+})_0 = C_A \times V_A = 0,080 \times 25 \times 10^{-3} = \mathbf{\SI{2,0e-3}{mol}}$.
$n(\ce{HO-})_0 = C_B \times V_B = 0,10 \times 15 \times 10^{-3} = \mathbf{\SI{1,5e-3}{mol}}$.
\item \textbf{Tableau d'avancement :}
\begin{center}
\begin{tabular}{|c|c|c|c|}\hline
 & \ce{H3O+} & \ce{HO-} & \ce{H2O} \\\hline
État initial & $2,0 \times 10^{-3}$ & $1,5 \times 10^{-3}$ & --- \\\hline
État intermédiaire & $2,0 \times 10^{-3} - x$ & $1,5 \times 10^{-3} - x$ & --- \\\hline
État final ($x_{\max}$) & $0,5 \times 10^{-3}$ & $0$ & --- \\\hline
\end{tabular}
\end{center}
$x_{\max} = \min(2,0; 1,5) \times 10^{-3} = \mathbf{\SI{1,5e-3}{mol}}$.
Réactif limitant : \ce{HO-} (soude). Réactif en excès : \ce{H3O+} (acide).
\item Volume total $V_T = 25 + 15 = \SI{40}{mL} = \SI{4,0e-2}{L}$.
$n(\ce{H3O+})_{\text{final}} = 0,5 \times 10^{-3} = \SI{5,0e-4}{mol}$.
$[\ce{H3O+}] = \frac{5,0 \times 10^{-4}}{4,0 \times 10^{-2}} = \mathbf{\SI{1,25e-2}{mol.L^{-1}}}$.
\item $\mathrm{pH} = -\log(1,25 \times 10^{-2}) = 2 - \log 1,25 = 2 - 0,097 = \mathbf{1,90}$.
Le mélange est \textbf{acide} (pH < 7).
\item Pour pH = 7 (équivalence), il faut $n(\ce{HO-}) = n(\ce{H3O+})_0 = \SI{2,0e-3}{mol}$.
$V_B(\text{équiv}) = \frac{n(\ce{HO-})}{C_B} = \frac{2,0 \times 10^{-3}}{0,10} = \SI{2,0e-2}{L} = \mathbf{\SI{20}{mL}}$.
\end{enumerate}
\end{exoresolu}

\begin{exoresolu}[Type Bac : préparation d'une solution au laboratoire]
Un laborantin d'un lycée de Douala doit préparer $V_f = \SI{250}{mL}$ d'une solution d'acide chlorhydrique de $\mathrm{pH} = 3$ à partir d'une solution commerciale d'acide chlorhydrique de concentration $C_0 = \SI{1,0}{mol.L^{-1}}$. La température du laboratoire est de \SI{25}{\celsius} ($\mathrm{pKe} = 14$).
\begin{enumerate}
\item Calcule la concentration $C_f$ de la solution à préparer.
\item Calcule le volume $V_0$ de solution commerciale à prélever. Comment appelle-t-on cette opération ? Quel est le facteur de dilution ?
\item Décris précisément le protocole expérimental : verrerie utilisée (nomme la verrerie de précision et justifie son choix), ordre des opérations, homogénéisation.
\item La solution commerciale est corrosive. Cite deux pictogrammes de sécurité que l'on peut trouver sur le flacon et deux règles de sécurité à respecter lors de la manipulation.
\item Le laborantin vérifie la solution préparée avec un pH-mètre et lit $\mathrm{pH} = 3{,}0$. Conclus sur la réussite de la préparation.
\item Il souhaite ensuite neutraliser $\SI{20}{mL}$ de sa solution de pH $= 3$ avec une solution de soude de concentration $C_B = \SI{5,0e-4}{mol.L^{-1}}$. Calcule le volume de soude nécessaire pour atteindre l'équivalence (pH $= 7$).
\end{enumerate}
\tcblower
\textbf{Solution détaillée :}
\begin{enumerate}
\item $\mathrm{pH} = 3 \Rightarrow [\ce{H3O+}] = 10^{-3} = \SI{1,0e-3}{mol.L^{-1}}$.
Comme \ce{HCl} est fort, $C_f = [\ce{H3O+}] = \mathbf{\SI{1,0e-3}{mol.L^{-1}}}$.
\item Relation de dilution : $C_0 V_0 = C_f V_f$.
$V_0 = \frac{C_f V_f}{C_0} = \frac{1,0 \times 10^{-3} \times 250 \times 10^{-3}}{1,0} = \mathbf{\SI{0,25}{mL}}$.
\textbf{Opération :} Dilution. \textbf{Facteur de dilution :} $f = \frac{C_0}{C_f} = \frac{1,0}{1,0 \times 10^{-3}} = \mathbf{1000}$.
\textit{Note :} $0,25$ mL est un volume très petit, difficile à prélever avec précision avec une pipette jaugée classique (1 mL ou 10 mL). En pratique, on ferait une \textbf{dilution en deux étapes} (ex: 1 mL $\to$ 100 mL puis 25 mL $\to$ 250 mL) ou on utiliserait une micropipette (si disponible).
\item \textbf{Protocole rigoureux :}
\begin{enumerate}
\item \textbf{Verrerie :} Pipette jaugée de \SI{1}{mL} (ou \SI{2}{mL}) classe A (précision $\pm \SI{0,006}{mL}$) + Fiole jaugée de \SI{250}{mL} classe A + Bécher + Entonnoir + Baguette de verre.
\item \textbf{Ordre :}
    \begin{itemize}
    \item Rincer la pipette avec la solution mère (\ce{HCl} \SI{1}{M}).
    \item Prélever $V_0$ (ou volume de la dilution intermédiaire) avec la pipette (poire à pipeter \textbf{jamais} à la bouche).
    \item Verser dans la fiole jaugée (par l'entonnoir).
    \item Rincer la pipette et l'entonnoir avec de l'eau distillée (verser les rinçages dans la fiole).
    \item Compléter à l'eau distillée \textbf{jusqu'au trait de jauge} (bas du ménisque au niveau du trait, œil au niveau du ménisque).
    \item Boucher et homogénéiser (retourner la fiole plusieurs fois).
    \end{itemize}
\end{enumerate}
\item \textbf{Pictogrammes (flacon \ce{HCl} \SI{12}{M} / \SI{35}{\%}) :}
\begin{itemize}
\item \textbf{GHS05} (Corrosion) : brûle peau/yeux, corrode métaux.
\item \textbf{GHS07} (Mention d'avertissement "Danger" / Torse) : Irritant voies respiratoires (vapeurs \ce{HCl} gazeux).
\end{itemize}
\textbf{Règles de sécurité :}
\begin{enumerate}
\item Porter lunettes, blouse, gants (néoprène), manipuler sous hotte (vapeurs \ce{HCl}).
\item \textbf{Verser l'acide dans l'eau} (jamais l'inverse) pour éviter projections par ébullition locale.
\end{enumerate}
\item Mesure pH-mètre : $\mathrm{pH}_{\text{mesuré}} = 3,0 = \mathrm{pH}_{\text{théorique}}$.
La préparation est \textbf{réussie} (concentration conforme, manipulation correcte).
\item Neutralisation : $\ce{H3O+ + HO- -> 2 H2O}$.
$n(\ce{H3O+}) = C_f \times V_A = 1,0 \times 10^{-3} \times 20 \times 10^{-3} = \SI{2,0e-5}{mol}$.
À l'équivalence : $n(\ce{HO-}) = n(\ce{H3O+}) = \SI{2,0e-5}{mol}$.
$V_B = \frac{n(\ce{HO-})}{C_B} = \frac{2,0 \times 10^{-5}}{5,0 \times 10^{-4}} = \mathbf{\SI{40}{mL}}$.
\end{enumerate}
\end{exoresolu}

\newpage

\coursec{Corrigés des exercices}

\begin{corrige}[Exercice 1 — QCM : acides et bases forts]
\begin{enumerate}
    \item \textbf{Réponse a) : pH = 3.}
    \ce{HCl} est un acide fort : sa réaction avec l'eau est totale. On a donc $[\ce{H3O+}] = C = \SI{1,0e-3}{mol.L^{-1}}$.
    \[ \mathrm{pH} = -\log[\ce{H3O+}] = -\log(1,0 \times 10^{-3}) = 3 \]
    La réponse b) ($11$) correspondrait à une base forte de même concentration ; les réponses c) et d) confondent pH et concentration ou négligent l'échelle logarithmique.

    \item \textbf{Réponse b) : $[\ce{HO-}] = \SI{1,0e-2}{mol.L^{-1}}$.}
    $\mathrm{pH} = 12 \Rightarrow [\ce{H3O+}] = 10^{-12}\ \si{mol.L^{-1}}$.
    On utilise le produit ionique de l'eau $\mathrm{Ke} = [\ce{H3O+}][\ce{HO-}] = \SI{1,0e-14}{}$ (à \SI{25}{\celsius}) :
    \[ [\ce{HO-}] = \frac{\mathrm{Ke}}{[\ce{H3O+}]} = \frac{10^{-14}}{10^{-12}} = \SI{1,0e-2}{mol.L^{-1}} \]
    Autre raisonnement rapide : $\mathrm{pOH} = \mathrm{pKe} - \mathrm{pH} = 14 - 12 = 2 \Rightarrow [\ce{HO-}] = 10^{-2}\ \si{mol.L^{-1}}$.

    \item \textbf{Réponse c) : pH = 3.}
    Diluer dix fois ($C' = C/10$) une solution d'acide fort \textbf{augmente} le pH d'une unité (car $\mathrm{pH}' = -\log(C/10) = -\log C + 1 = \mathrm{pH} + 1$).
    $\mathrm{pH}' = 2 + 1 = 3$. La solution reste acide (pH < 7) : elle ne peut jamais devenir basique par simple dilution, d'où l'absurdité de la réponse d).

    \item \textbf{Réponse b) : $[\ce{Cl-}] = C$.}
    L'équation de la réaction totale est : $\ce{HCl(aq) + H2O(l) -> H3O+(aq) + Cl-(aq)}$.
    D'après la stœchiométrie (conservation de la matière), chaque mole de \ce{HCl} introduite donne exactement une mole d'ions \ce{Cl-}. D'où $[\ce{Cl-}] = C$.
\end{enumerate}
\end{corrige}

\begin{corrige}[Exercice 2 — Vrai ou faux, avec justification]
\begin{enumerate}
    \item \textbf{Vrai.} C'est la définition même d'un acide fort : sa réaction avec l'eau est totale (flèche simple \ce{->}), toutes les molécules cèdent leur proton à l'eau pour former des ions oxonium \ce{H3O+}.

    \item \textbf{Faux.} La relation $\mathrm{pH} = -\log C$ n'est valable que pour des concentrations $C > \SI{e-6}{mol.L^{-1}}$. En dessous, la contribution de l'autoprotolyse de l'eau ($[\ce{H3O+}]_{\text{eau}} = \SI{e-7}{mol.L^{-1}}$) n'est plus négligeable devant $C$ ; le pH tend alors vers 7 (sans jamais l'atteindre pour un acide).

    \item \textbf{Faux.} La soude (\ce{NaOH}) est une base forte : elle se dissocie totalement. Pour $C = \SI{e-4}{mol.L^{-1}}$ :
    \[ \mathrm{pH} = \mathrm{pKe} + \log C = 14 + \log(10^{-4}) = 14 - 4 = 10 \]
    Un pH de 4 correspondrait à un \emph{acide} fort de concentration $\SI{e-4}{mol.L^{-1}}$.

    \item \textbf{Faux.} On a bien $[\ce{Na+}] = C$ (dissociation totale du sel). En revanche, $[\ce{HO-}] = C + [\ce{HO-}]_{\text{eau}}$, car l'autoprotolyse de l'eau apporte aussi des ions \ce{HO-}. L'égalité $[\ce{Na+}] = [\ce{HO-}]$ n'est donc qu'approchée (excellente approximation si $C \gg \SI{e-7}{mol.L^{-1}}$). L'électroneutralité exacte s'écrit :
    \[ [\ce{Na+}] + [\ce{H3O+}] = [\ce{HO-}] \]

    \item \textbf{Faux.} Un acide, même très dilué, ne peut pas donner un pH basique ($> 7$). Pour $C = \SI{e-8}{mol.L^{-1}}$, la formule $\mathrm{pH} = -\log C$ n'est plus valide (condition $C > \SI{e-6}{mol.L^{-1}}$ non respectée). Les ions \ce{H3O+} issus de l'eau ($10^{-7}\ \si{mol.L^{-1}}$) dominent ; le pH est voisin de 7 (légèrement inférieur, environ 6,98).

    \item \textbf{Vrai.} L'ion éthanolate \ce{C2H5O-} est la base conjuguée de l'éthanol (acide très faible, $\mathrm{p}K_\text{a} \approx 16$). Il réagit \textbf{totalement} avec l'eau selon :
    \[ \ce{C2H5O- (aq) + H2O(l) -> C2H5OH(aq) + HO-(aq)} \]
    L'éthanolate de sodium (\ce{C2H5ONa}) est donc bien une base forte (c'est un sel qui donne une solution basique par hydrolyse totale de l'anion).
    \textbf{Attention :} c'est un solide hygroscopique, corrosif (pictogramme G05) et inflammable ; il se manipule sous atmosphère inerte.
\end{enumerate}
\end{corrige}

\begin{corrige}[Exercice 3 — Texte à trous]
Le chlorure d'hydrogène \ce{HCl} réagit avec l'eau selon une réaction \textbf{totale} : c'est donc un acide \textbf{fort}. La solution obtenue contient des ions \textbf{oxonium} \ce{H3O+} et des ions chlorure \ce{Cl-}. Le pH d'une telle solution de concentration $C$ (avec $C > \SI{e-6}{mol.L^{-1}}$) est donné par la relation $\mathrm{pH} = \textbf{$-\log C$}$.

L'hydroxyde de sodium \ce{NaOH} se dissocie totalement dans l'eau en ions sodium \ce{Na+} et ions \textbf{hydroxyde} \ce{HO-} : c'est une base \textbf{forte}. Le pH de sa solution est donné par $\mathrm{pH} = \textbf{$\mathrm{pKe} + \log C$}$.

Quand on dilue dix fois une solution d'acide fort, son pH \textbf{augmente} d'une unité. Quand on dilue dix fois une solution de base forte, son pH \textbf{diminue} d'une unité. Dans toute solution aqueuse, la somme des charges positives est égale à la somme des charges négatives : c'est la loi d'\textbf{électroneutralité}.

\textbf{Note pédagogique :} les adjectifs « fort » et « forte » (pour qualifier l'acide et la base) ne figuraient pas dans la liste de mots proposée à l'élève ; ils sont toutefois les termes attendus du cours. Les mots de la liste effectivement utilisés sont : \emph{totale}, \emph{ion oxonium}, \emph{$-\log C$}, \emph{ion hydroxyde}, \emph{$\mathrm{pKe} + \log C$}, \emph{augmente}, \emph{diminue}, \emph{électroneutralité}. Les mots \emph{partielle} et \emph{totale} (seconde occurrence dans la liste) ne sont pas utilisés : « partielle » caractérise un acide \emph{faible} (leçon suivante).
\end{corrige}

\begin{corrige}[Exercice 4 — Définitions et formules essentielles]
\begin{enumerate}
    \item \textbf{Acide fort :} espèce chimique (souvent un gaz comme \ce{HCl} ou \ce{HNO3}) dont la réaction avec l'eau est \textbf{totale} (flèche simple \ce{->}) : il cède la totalité de ses protons à l'eau pour former des ions oxonium \ce{H3O+}.
    \emph{Exemples usuels :} chlorure d'hydrogène \ce{HCl} (acide chlorhydrique), acide nitrique \ce{HNO3}, acide sulfurique \ce{H2SO4} (pour le premier proton seulement).

    \item \textbf{Base forte :} espèce chimique (souvent un solide ionique comme \ce{NaOH} ou \ce{KOH}) qui se dissocie (ou réagit) \textbf{totalement} dans l'eau en libérant des ions hydroxyde \ce{HO-}.
    \emph{Exemples usuels :} hydroxyde de sodium \ce{NaOH} (soude), hydroxyde de potassium \ce{KOH} (potasse), éthanolate de sodium \ce{C2H5ONa}.

    \item \textbf{Équation de la réaction de l'acide chlorhydrique avec l'eau (totale, flèche simple) :}
    \[ \ce{HCl(aq) + H2O(l) -> H3O+(aq) + Cl-(aq)} \]
    \textbf{Couples acide/base mis en jeu :}
    \begin{itemize}
        \item \ce{HCl}/\ce{Cl-} : \ce{HCl} cède un proton (acide), \ce{Cl-} est sa base conjuguée (très faible).
        \item \ce{H3O+}/\ce{H2O} : l'eau capte un proton (base), \ce{H3O+} est son acide conjugué.
    \end{itemize}

    \item \textbf{Produit ionique de l'eau (Ke) :} dans toute solution aqueuse à \SI{25}{\celsius}, le produit des concentrations molaires en ions oxonium et hydroxyde est constant :
    \[ \mathrm{Ke} = [\ce{H3O+}] \times [\ce{HO-}] = \SI{1,0e-14}{} \qquad (\mathrm{pKe} = -\log \mathrm{Ke} = 14) \]
    Cette relation fondamentale est \textbf{toujours vérifiée}, quelle que soit la nature de la solution (acide, neutre ou basique). Elle permet de passer du pH au pOH : $\mathrm{pH} + \mathrm{pOH} = \mathrm{pKe} = 14$.
\end{enumerate}
\end{corrige}

\newpage

\coursec{Fiche bilan}

\coursec{Acides forts et bases fortes}

\textbf{Situation déclenchante :} Au laboratoire de chimie du lycée, tu dois préparer \SI{500}{mL} d'une solution d'acide chlorhydrique \SI{0,10}{mol.L^{-1}} à partir d'une solution commerciale concentrée (\SI{12}{mol.L^{-1}}, \SI{37}{\%} massique). Quel volume prélever ? Quel sera le pH de la solution préparée ? Si tu verses par erreur \SI{10}{mL} de cette solution dans \SI{20}{mL} d'une solution de soude \SI{0,10}{mol.L^{-1}}, quel sera le pH final ? Cette leçon t'apporte les outils rigoureux pour répondre à ces questions de chimiste.

\courssub{L'acide chlorhydrique, modèle d'acide fort}

\begin{definition}[Acide fort]
Un \cle{acide fort} est un acide dont la réaction avec l'eau est \textbf{totale} (irréversible). Il se trouve entièrement ionisé en solution aqueuse : il n'existe plus de molécules d'acide non dissociées. L'acide chlorhydrique \ce{HCl} est le modèle de référence.
\end{definition}

\begin{propriete}[Équation d'ionisation et conséquence sur le pH]
L'acide chlorhydrique gazeux se dissout dans l'eau et réagit totalement avec elle :
\[ \ce{HCl(g) + H2O(l) -> H3O+(aq) + Cl-(aq)} \]
On écrit souvent de façon simplifiée : \ce{HCl -> H+ + Cl-}, mais la forme hydratée \ce{H3O+} (ion oxonium) est la rigoureuse.
\par
Conséquence immédiate : la concentration en ions oxonium est égale à la concentration molaire introduite de l'acide :
\[ [\ce{H3O+}] = C_a \quad (\text{concentration formelle de l'acide}) \]
Le pH se calcule donc directement :
\[ \mathrm{pH} = -\log_{10} [\ce{H3O+}] = -\log_{10} C_a \]
\end{propriete}

\begin{exemplebox}[Calcul du pH d'une solution d'acide chlorhydrique]
\textbf{Énoncé :} Calculer le pH d'une solution d'acide chlorhydrique de concentration molaire $C_a = \SI{2,5e-3}{mol.L^{-1}}$ à \SI{25}{\celsius}.
\par
\textbf{Solution :}
\begin{align*}
[\ce{H3O+}] &= C_a = \SI{2,5e-3}{mol.L^{-1}} \\
\mathrm{pH} &= -\log_{10}(2,5 \times 10^{-3}) \\
\mathrm{pH} &= -(\log_{10} 2,5 + \log_{10} 10^{-3}) \\
\mathrm{pH} &= -(0,3979 - 3) = 2,60
\end{align*}
\textbf{Réponse :} Le pH de cette solution est \textbf{2,60}.
\end{exemplebox}

\begin{methode}[Montrer qu'un acide est fort (critère expérimental)]
Pour affirmer qu'un acide est fort à partir de mesures :
\begin{enumerate}
  \item Mesurer le pH de la solution (pH-mètre étalonné).
  \item Connaître la concentration molaire $C_a$ de la solution (préparée par pesée/dilution rigoureuse).
  \item Calculer la valeur théorique $\mathrm{pH}_{\text{th}} = -\log_{10} C_a$.
  \item Comparer : si $\mathrm{pH}_{\text{mesuré}} \approx \mathrm{pH}_{\text{th}}$ (à \SI{0,05} près), l'acide est totalement ionisé $\Rightarrow$ \textbf{acide fort}.
\end{enumerate}
\emph{Exemple :} Pour $C_a = \SI{0,10}{mol.L^{-1}}$, $\mathrm{pH}_{\text{th}} = 1,00$. Si on mesure $\mathrm{pH} = 1,02$, c'est un acide fort. Si on mesure $\mathrm{pH} = 1,8$, c'est un acide faible.
\end{methode}

\begin{experience}[Mesure du pH de solutions d'HCl de concentrations connues]
\textbf{Matériel :} pH-mètre étalonné (tampons pH 4,01 ; 7,00 ; 10,01), béchers, solutions mères d'HCl (\SI{1}{mol.L^{-1}}), eau distillée, pipettes graduées.
\par
\textbf{Protocole :}
\begin{enumerate}
  \item Préparer par dilution au dixième successives 4 solutions : $S_1$ (\SI{0,1}{M}), $S_2$ (\SI{0,01}{M}), $S_3$ (\SI{0,001}{M}), $S_4$ (\SI{0,0001}{M}).
  \item Mesurer le pH de chaque solution après rinçage de l'électrode à l'eau distillée puis à la solution testée.
\end{enumerate}
\textbf{Observations attendues :}
\begin{center}
\begin{tabular}{|c|c|c|c|}
\hline
Solution & $C_a$ (mol/L) & $\mathrm{pH}_{\text{th}} = -\log C_a$ & $\mathrm{pH}_{\text{mesuré}}$ \\
\hline
$S_1$ & $10^{-1}$ & 1,00 & $\approx 1,0$ \\
$S_2$ & $10^{-2}$ & 2,00 & $\approx 2,0$ \\
$S_3$ & $10^{-3}$ & 3,00 & $\approx 3,0$ \\
$S_4$ & $10^{-4}$ & 4,00 & $\approx 4,0$ \\
\hline
\end{tabular}
\end{center}
\textbf{Interprétation :} L'égalité $\mathrm{pH}_{\text{mesuré}} = -\log C_a$ confirme la dissociation totale de l'HCl. On observe que \textbf{diluer 10 fois augmente le pH de 1 unité}.
\end{experience}

\courssub{L'hydroxyde de sodium, modèle de base forte}

\begin{definition}[Base forte]
Une \cle{base forte} est une base dont la réaction avec l'eau (ou la dissolution pour les hydroxydes ioniques) est \textbf{totale}. Elle fournit entièrement des ions hydroxyde \ce{HO-} en solution. L'hydroxyde de sodium \ce{NaOH} (soude) est le modèle de référence.
\end{definition}

\begin{propriete}[Dissociation et calcul du pH]
L'hydroxyde de sodium est un solide ionique qui se dissocie totalement dans l'eau :
\[ \ce{NaOH(s) ->[H2O] Na+(aq) + HO-(aq)} \]
La concentration en ions hydroxyde est égale à la concentration formelle de la base :
\[ [\ce{HO-}] = C_b \]
On utilise le produit ionique de l'eau $K_e = [\ce{H3O+}][\ce{HO-}] = 10^{-14}$ (à \SI{25}{\celsius}) pour trouver $[\ce{H3O+}]$ et le pH :
\[ [\ce{H3O+}] = \frac{K_e}{C_b} = \frac{10^{-14}}{C_b} \]
\[ \mathrm{pH} = -\log_{10} \left( \frac{10^{-14}}{C_b} \right) = 14 - (-\log_{10} C_b) = \mathbf{14 + \log_{10} C_b} \]
\end{propriete}

\begin{exemplebox}[Calcul du pH d'une solution de soude]
\textbf{Énoncé :} On a dissous \SI{0,40}{g} d'hydroxyde de sodium solide (puré à 100 \%) dans de l'eau pour obtenir \SI{250,0}{mL} de solution. $M(\ce{NaOH}) = \SI{40,0}{g.mol^{-1}}$. Calculer le pH de cette solution à \SI{25}{\celsius}.
\par
\textbf{Solution :}
\begin{enumerate}
  \item Quantité de matière : $n(\ce{NaOH}) = \frac{0,40}{40,0} = \SI{0,010}{mol}$.
  \item Concentration : $C_b = \frac{0,010}{0,250} = \SI{0,040}{mol.L^{-1}} = \SI{4,0e-2}{mol.L^{-1}}$.
  \item pH : $\mathrm{pH} = 14 + \log_{10}(4,0 \times 10^{-2}) = 14 + (\log_{10} 4,0 - 2) = 14 + (0,60 - 2) = \mathbf{12,60}$.
\end{enumerate}
\end{exemplebox}

\courssub{Espèces chimiques en solution : inventaire et calcul des concentrations}

\begin{propriete}[Relations fondamentales dans une solution aqueuse]
Pour toute solution aqueuse à \SI{25}{\celsius}, on dispose de trois relations majeures :
\begin{enumerate}
  \item \textbf{Électroneutralité :} La somme des charges positives égale la somme des charges négatives.
  \[ \sum n_i [\text{Cation}_i] = \sum m_j [\text{Anion}_j] \]
  \item \textbf{Conservation de la matière :} Pour chaque élément ou groupe d'atomes, la quantité totale introduite égale la somme des quantités présentes sous toutes les formes.
  \item \textbf{Produit ionique de l'eau :} $[\ce{H3O+}][\ce{HO-}] = K_e = 10^{-14}$.
\end{enumerate}
\end{propriete}

\begin{methode}[Inventaire des espèces et calcul des concentrations (Acide fort / Base forte)]
\begin{enumerate}
  \item Écrire l'équation de dissolution/ionisation (flèche \ce{->}).
  \item Lister les espèces \textbf{réellement présentes} en solution (éliminer les réactifs totally consommés).
  \item Écrire la relation d'électroneutralité.
  \item Écrire les relations de conservation de la matière.
  \item Utiliser $K_e$ si nécessaire.
  \item Résoudre le système (souvent trivial pour acide/base fort seul).
\end{enumerate}
\end{methode}

\begin{exemplebox}[Inventaire pour une solution d'HCl \SI{0,10}{mol.L^{-1}}]
\textbf{Espèces introduites :} \ce{HCl}, \ce{H2O}.
\textbf{Équation :} \ce{HCl + H2O -> H3O+ + Cl-} (totale).
\textbf{Espèces présentes à l'équilibre :} \ce{H2O}, \ce{H3O+}, \ce{Cl-}. (\ce{HCl} absent).
\par
\textbf{Concentrations :}
\begin{itemize}
  \item Conservation Cl : $[\ce{Cl-}] = C_a = \SI{0,10}{mol.L^{-1}}$.
  \item Électroneutralité : $[\ce{H3O+}] = [\ce{Cl-}] + [\ce{HO-}] \approx [\ce{Cl-}]$ (car $[\ce{HO-}] = 10^{-13} \ll 0,10$).
  \item Donc $[\ce{H3O+}] = \SI{0,10}{mol.L^{-1}}$, $\mathrm{pH} = 1,00$.
  \item $[\ce{HO-}] = \frac{10^{-14}}{0,10} = \SI{1e-13}{mol.L^{-1}}$.
\end{itemize}
\end{exemplebox}

\begin{exemplebox}[Inventaire pour une solution de NaOH \SI{0,010}{mol.L^{-1}}]
\textbf{Espèces présentes :} \ce{H2O}, \ce{Na+}, \ce{HO-}.
\par
\textbf{Concentrations :}
\begin{itemize}
  \item Conservation Na : $[\ce{Na+}] = C_b = \SI{0,010}{mol.L^{-1}}$.
  \item Électroneutralité : $[\ce{Na+}] + [\ce{H3O+}] = [\ce{HO-}] \Rightarrow [\ce{HO-}] \approx [\ce{Na+}] = \SI{0,010}{mol.L^{-1}}$.
  \item $[\ce{H3O+}] = \frac{10^{-14}}{0,010} = \SI{1e-12}{mol.L^{-1}}$.
  \item $\mathrm{pH} = 14 + \log(0,010) = 12,00$.
\end{itemize}
\end{exemplebox}

\courssub{Effet de la dilution sur le pH}

\begin{propriete}[Loi de dilution et variation du pH]
Soit une solution mère d'acide fort (concentration $C_i$, volume $V_i$) diluée à un volume final $V_f$ (solution fille de concentration $C_f$).
La quantité de matière en solute est conservée : $n = C_i V_i = C_f V_f$.
Le facteur de dilution est $d = \frac{V_f}{V_i} = \frac{C_i}{C_f}$.
\par
Pour un \textbf{acide fort} ($C_f > 10^{-6}$ M) :
\[ \mathrm{pH}_f = -\log C_f = -\log \left( \frac{C_i}{d} \right) = -\log C_i + \log d = \mathrm{pH}_i + \log d \]
\textbf{Règle pratique :} Une dilution \textbf{au dixième} ($d=10$) fait \textbf{augmenter le pH de 1 unité}. Une dilution au centième ($d=100$) l'augmente de 2 unités.
\par
Pour une \textbf{base forte} ($C_f > 10^{-6}$ M) :
\[ \mathrm{pH}_f = 14 + \log C_f = 14 + \log \left( \frac{C_i}{d} \right) = \mathrm{pH}_i - \log d \]
\textbf{Règle pratique :} Une dilution \textbf{au dixième} fait \textbf{diminuer le pH de 1 unité}.
\end{propriete}

\begin{attention}[Limites de validité des formules $\mathrm{pH} = -\log C$ et $\mathrm{pH} = 14 + \log C$]
Ces formules supposent que la concentration en ions \ce{H3O+} (ou \ce{HO-}) apportée par l'acide (ou la base) est \textbf{très supérieure} à celle apportée par l'autoprotolyse de l'eau ($10^{-7}$ M).
\begin{itemize}
  \item \textbf{Acide fort :} Valide si $C_a > \SI{1e-6}{mol.L^{-1}}$ (pH < 6). En dessous, le pH tend vers 7 (ne descend jamais en dessous de 7 pour un acide !).
  \item \textbf{Base forte :} Valide si $C_b > \SI{1e-6}{mol.L^{-1}}$ (pH > 8). Au-dessus, le pH tend vers 7 (ne monte jamais au-dessus de 7 pour une base !).
\end{itemize}
\textbf{Piège fréquent :} Calculer le pH d'un acide fort \SI{1e-8}{M} comme $-\log(10^{-8}) = 8$ (basique !). \textbf{Faux.} Le pH est $\approx 7$ (légèrement acide, $\approx 6,98$).
\end{attention}

\begin{popfigure}
\begin{center}
\begin{tikzpicture}
\begin{axis}[
    width=\linewidth, height=8cm,
    xlabel={Facteur de dilution $d$ (échelle log)}, ylabel={pH},
    xmode=log, log basis x=10,
    domain=1:1e6, samples=100,
    xmin=1, xmax=1e6, ymin=0, ymax=14,
    grid=major, grid style={popline},
    legend pos=south east,
    legend cell align=left,
    tick label style={font=\small},
    label style={font=\small},
]
% Acide fort initial pH=1 (C=0.1M)
\addplot[popblue, thick, popcurve] {1 + log10(x)};
\addlegendentry{Acide fort ($C_i=0,1$ M, pH$_i=1$)}

% Base forte initial pH=13 (C=0.1M)
\addplot[poporange, thick, popcurve] {13 - log10(x)};
\addlegendentry{Base forte ($C_i=0,1$ M, pH$_i=13$)}

% Ligne pH=7
\addplot[popmuted, dashed, thick] coordinates {(1,7) (1e6,7)};
\addlegendentry{Neutralité (pH=7)};

% Zone de validité
\fill[popblueL, opacity=0.2] (1,0) rectangle (1e5,7);
\fill[poporangeL, opacity=0.2] (1,7) rectangle (1e5,14);
\node[popblue, font=\tiny, anchor=south west] at (1.5, 1) {Valide ($C>10^{-6}$)};
\node[poporange, font=\tiny, anchor=north west] at (1.5, 13) {Valide ($C>10^{-6}$)};
\node[popmuted, font=\tiny, rotate=90, anchor=south] at (1e5, 3.5) {Limite};
\end{axis}
\end{tikzpicture}
\legende{Variation du pH lors de la dilution d'un acide fort (bleu) et d'une base forte (orange) initialement à 0,1 M. Au-delà d'un facteur $10^5$ ($C < 10^{-6}$ M), les courbes dévient et tendent vers 7.}
\end{center}
\end{popfigure}

\begin{exemplebox}[Préparation par dilution et prévision du pH]
\textbf{Énoncé :} On souhaite préparer \SI{250,0}{mL} d'une solution d'HCl \SI{5,0e-3}{mol.L^{-1}} à partir d'une solution mère \SI{0,100}{mol.L^{-1}}. Prévoir le pH de la solution fille.
\par
\textbf{Solution :}
\begin{enumerate}
  \item Volume à prélever : $V_i = \frac{C_f V_f}{C_i} = \frac{5,0 \times 10^{-3} \times 0,250}{0,100} = \SI{12,5}{mL}$.
  \item Opération : Pipeter \SI{12,5}{mL} de solution mère dans une fiole jaugée \SI{250}{mL}, compléter à l'eau distillée jusqu'au trait de jauge, homogénéiser.
  \item Prévision pH : Facteur de dilution $d = \frac{0,100}{5,0 \times 10^{-3}} = 20$.
  $\mathrm{pH}_i = -\log(0,100) = 1,00$.
  $\mathrm{pH}_f = 1,00 + \log_{10}(20) = 1,00 + 1,30 = \mathbf{2,30}$.
  \item Vérification validité : $C_f = \SI{5e-3}{M} > \SI{1e-6}{M}$ \checkmark.
\end{enumerate}
\end{exemplebox}

\courssub{Mélange d'un acide fort et d'une base forte}

\begin{propriete}[Réaction de neutralisation]
Lorsqu'on mélange un acide fort (source de \ce{H3O+}) et une base forte (source de \ce{HO-}), la réaction qui a lieu est :
\[ \ce{H3O+(aq) + HO-(aq) -> 2 H2O(l)} \]
C'est une réaction \textbf{quasi totale} (constante d'équilibre $K = 1/K_e = 10^{14}$). Elle va jusqu'à épuisement de l'un des réactifs.
\end{propriete}

\begin{methode}[Calcul du pH d'un mélange Acide fort + Base forte]
\begin{enumerate}
  \item Calculer les quantités de matière initiales :
  $n_{\text{init}}(\ce{H3O+}) = C_a \times V_a$ ; $n_{\text{init}}(\ce{HO-}) = C_b \times V_b$.
  \item Faire le \textbf{tableau d'avancement} de la réaction \ce{H3O+ + HO- -> 2H2O}.
  \item Identifier le \textbf{réactif en excès} (l'autre est limitant, quantité finale = 0).
  \item Calculer la concentration de l'espèce en excès dans le \textbf{volume total} $V_{\text{tot}} = V_a + V_b$.
  \item En déduire le pH :
  \begin{itemize}
    \item Excès \ce{H3O+} : $[\ce{H3O+}] = \frac{n_{\text{excès}}(\ce{H3O+})}{V_{\text{tot}}}$ ; $\mathrm{pH} = -\log [\ce{H3O+}]$.
    \item Excès \ce{HO-} : $[\ce{HO-}] = \frac{n_{\text{excès}}(\ce{HO-})}{V_{\text{tot}}}$ ; $\mathrm{pH} = 14 + \log [\ce{HO-}]$.
    \item \textbf{Équivalence stoichiométrique} ($n_{\text{init}}(\ce{H3O+}) = n_{\text{init}}(\ce{HO-})$) : $[\ce{H3O+}] = [\ce{HO-}] = 10^{-7}$ M $\Rightarrow$ $\mathrm{pH} = 7$ (à \SI{25}{\celsius}).
  \end{itemize}
\end{enumerate}
\end{methode}

\begin{exemplebox}[Mélange HCl / NaOH : détermination du pH final]
\textbf{Énoncé :} On mélange \SI{20,0}{mL} d'HCl \SI{0,10}{mol.L^{-1}} avec \SI{30,0}{mL} de NaOH \SI{0,050}{mol.L^{-1}}. Calculer le pH final à \SI{25}{\celsius}.
\par
\textbf{Solution :}
\begin{enumerate}
  \item $n_{\text{init}}(\ce{H3O+}) = 0,100 \times 0,0200 = \SI{2,00e-3}{mol}$.
  $n_{\text{init}}(\ce{HO-}) = 0,050 \times 0,0300 = \SI{1,50e-3}{mol}$.
  \item Tableau d'avancement ($x$ en mol) :
  \begin{center}
  \begin{tabular}{|c|c|c|c|}
  \hline
  & \ce{H3O+} & \ce{HO-} & \ce{H2O} \\
  \hline
  État initial & $2,00 \times 10^{-3}$ & $1,50 \times 10^{-3}$ & --- \\
  \hline
  Évolution & $-x$ & $-x$ & $+2x$ \\
  \hline
  État final & $2,00 \times 10^{-3} - x$ & $1,50 \times 10^{-3} - x$ & --- \\
  \hline
  \end{tabular}
  \end{center}
  \item Réactif limitant : \ce{HO-} (quantité moindre). $x_{\text{max}} = \SI{1,50e-3}{mol}$.
  \item Excès \ce{H3O+} : $n_{\text{excès}}(\ce{H3O+}) = 2,00 \times 10^{-3} - 1,50 \times 10^{-3} = \SI{5,00e-4}{mol}$.
  \item Volume total $V_{\text{tot}} = \SI{50,0}{mL} = \SI{0,0500}{L}$.
  $[\ce{H3O+}] = \frac{0,50 \times 10^{-3}}{0,0500} = \SI{0,010}{mol.L^{-1}} = \SI{1,0e-2}{mol.L^{-1}}$.
  \item $\mathrm{pH} = -\log(1,0 \times 10^{-2}) = \mathbf{2,00}$.
\end{enumerate}
\end{exemplebox}

\begin{experience}[Suivi pH-métrique du dosage d'un acide fort par une base forte]
\textbf{Dispositif :} Bécher avec \SI{20,0}{mL} d'HCl \SI{0,10}{M} + eau distillée + barreau magnétique. pH-mètre immergé. Burette remplie de NaOH \SI{0,10}{M}.
\par
\textbf{Protocole :} Ajouter la soude par fractions (\SI{2}{mL} puis \SI{0,5}{mL} près de l'équivalence), noter le pH après chaque ajout et homogénéisation.
\par
\textbf{Observations / Courbe attendue :}
\begin{itemize}
  \item Zone 1 (Avant équivalence, $V_b < 20$ mL) : Excès \ce{H3O+}, pH acide, monte lentement.
  \item Zone 2 (Au voisinage de l'équivalence $V_b \approx 20$ mL) : Saut de pH brutal (pH 4 $\to$ 10).
  \item Zone 3 (Après équivalence, $V_b > 20$ mL) : Excès \ce{HO-}, pH basique, monte lentement vers pH de la soude.
  \item Point d'équivalence : $V_{E} = \SI{20,0}{mL}$, $\mathrm{pH}_E = 7$.
\end{itemize}
\emph{Application :} Ce principe est utilisé en industrie (SONARA, brasseries, savonneries) pour titrer l'acidité ou la basicité des produits.
\end{experience}

\courssub{Autres acides forts et bases fortes usuels ; sécurité}

\begin{savaistu}[Acides et bases forts dans l'industrie et le quotidien camerounais]
\begin{itemize}
  \item \textbf{Acide nitrique \ce{HNO3}} : Utilisé à la SONARA (raffinerie de Limbé) pour la nitration, fabrication d'engrais (nitrate d'ammonium), gravure sur métaux. Fumées brunes toxiques (\ce{NO2}) au contact de la lumière/air.
  \item \textbf{Acide sulfurique \ce{H2SO4}} : \textbf{1ère acidité forte} (\ce{H2SO4 -> H3O+ + HSO4-}). 2ème acidité faible (\ce{HSO4- <=> H3O+ + SO4^2-}, pKa $\approx$ 1,9). Utilisé dans les batteries de voiture (acide de batterie), fabrication d'engrais (CIMENCAM pour le gypse), lessivage minerais.
  \item \textbf{Potasse \ce{KOH}} : Base forte, plus soluble que NaOH. Utilisée pour la fabrication de savons mous (savon noir, savon liquide artisanaux au Cameroun), électrolytes alcalins.
  \item \textbf{Éthanolate de sodium \ce{C2H5ONa}} : Base forte non aqueuse (dans l'éthanol). Réagit violemment avec l'eau : \ce{C2H5O- + H2O -> C2H5OH + HO-}. Utilisé en synthèse organique (transestérification pour biodiesel, condensation de Claisen).
  \item \textbf{Contexte local :} Le \emph{bili-bili} (bière de mil) et le \emph{vin de palme} contiennent de l'acide acétique (acide \textbf{faible}), pas de l'HCl ! L'acidité du \emph{bissap} (hibiscus) est due à des acides organiques faibles (citrique, malique, hibiscus). Seuls les produits industriels (nettoyants WC, déboucheurs, acide de batterie) contiennent des acides forts.
\end{itemize}
\end{savaistu}

\begin{attention}[Règles de sécurité impératives (Pictogrammes GHS)]
\begin{center}
\begin{tabular}{|c|c|p{8cm}|}
\hline
\textbf{Pictogramme} & \textbf{Désignation} & \textbf{Produits concernés \& Consignes} \\
\hline
\textbf{GHS05} & \textbf{Corrosif (GHS05)} & \ce{HCl} conc., \ce{HNO3} conc., \ce{H2SO4} conc., \ce{NaOH} solide/conc., \ce{KOH}. \\
& & \textbf{Risque :} Brûlures graves peau/yeux, atteinte voies respiratoires. \\
& & \textbf{Protection :} Blouse, \textbf{gants nitrile}, \textbf{lunettes de protection} obligatoires. \\
& & \textbf{Manipulation :} \textbf{Toujours verser l'acide (ou la base) dans l'eau}, jamais l'inverse (éclaboussures, projection chaleur). \\
\hline
\textbf{GHS06} & \textbf{Toxique (GHS06)} & \ce{HNO3} fumant, \ce{NaCN}, \ce{KCN}. Danger mortel par ingestion/inhalation. Manipulation sous hotte. \\
\hline
\textbf{GHS02} & \textbf{Inflammable (GHS02)} & \ce{C2H5ONa} (éthanolate), solutions alcooliques. Tenir loin flammes/étincelles. \\
\hline
\end{tabular}
\end{center}
\textbf{En cas de projection :} Rincer \textbf{immédiatement} et \textbf{longuement} (min 15 min) à l'eau courante (douche oculaire / douche de sécurité). Avertir l'enseignant. Ne pas neutraliser sur la peau (réaction exothermique).
\end{attention}

\begin{popfigure}
\begin{center}
\begin{tikzpicture}[
  node distance=2.5cm and 1.5cm,
  concept/.style={draw, thick, rounded corners, fill=popblue!20, text width=2.8cm, align=center, font=\small\bfseries, minimum height=1.5cm},
  central/.style={concept, fill=popink, text=white, text width=3.5cm, minimum height=2cm, font=\normalsize\bfseries},
  link/.style={-Stealth, thick, popmuted}
]
\node[central] (center) {Acides forts \& Bases fortes};
\node[concept, fill=popblue!30] (acide) at (135:3.5cm) {Acide fort\\ réaction totale\\ pH $=-\log C_a$\\ Ex: HCl, HNO$_3$};
\node[concept, fill=poporange!30] (base) at (45:3.5cm) {Base forte\\ dissociation totale\\ pH $=14+\log C_b$\\ Ex: NaOH, KOH};
\node[concept, fill=popgreen!30] (especes) at (-45:3.5cm) {Espèces en solution\\ électroneutralité\\ conservation matière\\ $K_e=10^{-14}$};
\node[concept, fill=poppurple!30] (dilution) at (-135:3.5cm) {Dilution\\ $\div 10 \Rightarrow \Delta$pH $=\pm 1$\\ valable si $C>10^{-6}$ M};
\node[concept, fill=poppink!30] (melange) at (90:4.5cm) {Mélange acide+base\\ $\ce{H3O+ + HO- -> 2H2O}$\\ réactif en excès\\ équivalence $\Rightarrow$ pH=7};
\node[concept, fill=popgold!30] (securite) at (-90:4.5cm) {Sécurité\\ GHS05 Corrosif\\ Acide dans l'eau\\ EPI: gants, lunettes};

\draw[link] (center) -- (acide);
\draw[link] (center) -- (base);
\draw[link] (center) -- (especes);
\draw[link] (center) -- (dilution);
\draw[link] (center) -- (melange);
\draw[link] (center) -- (securite);
\end{tikzpicture}
\legende{Carte mentale de la leçon : les six idées forces à retenir pour le Bac.}
\end{center}
\end{popfigure}

\begin{definition}[Rappel des définitions clés]
\begin{itemize}
  \item \cle{Acide fort} : Acide réagissant \textbf{totalement} avec l'eau : \ce{AH + H2O -> H3O+ + A-}. Exemples : \ce{HCl, HNO3, H2SO4} (1ère acidité).
  \item \cle{Base forte} : Base se dissociant/réagissant \textbf{totalement} pour donner \ce{HO-}. Exemples : \ce{NaOH, KOH, C2H5ONa}.
  \item \cle{Ion oxonium} \ce{H3O+} : Espèce responsable du caractère acide. \cle{Ion hydroxyde} \ce{HO-} : Espèce responsable du caractère basique.
  \item \cle{Produit ionique de l'eau} $K_e = [\ce{H3O+}][\ce{HO-}] = 10^{-14}$ à \SI{25}{\celsius} $\Rightarrow$ $\mathrm{p}K_e = 14$.
\end{itemize}
\end{definition}

\begin{propriete}[Équations-types et formules fondamentales à connaître par cœur]
\begin{center}
\begin{tabularx}{\linewidth}{|l|X|X|}
\hline
\rowcolor{popblueL} \textbf{Grandeur} & \textbf{Acide fort (concentration $C_a$)} & \textbf{Base forte (concentration $C_b$)} \\
\hline
Équation & \ce{HCl + H2O -> H3O+ + Cl-} & \ce{NaOH -> Na+ + HO-} \\
\hline
$[\ce{H3O+}]$ & $C_a$ & $\dfrac{K_e}{C_b}$ \\
\hline
$[\ce{HO-}]$ & $\dfrac{K_e}{C_a}$ & $C_b$ \\
\hline
pH & $\mathrm{pH} = -\log C_a$ & $\mathrm{pH} = 14 + \log C_b$ \\
\hline
$C$ à partir du pH & $C_a = 10^{-\mathrm{pH}}$ & $C_b = 10^{\mathrm{pH}-14}$ \\
\hline
\end{tabularx}
\end{center}
\textbf{Neutralisation :} \ce{H3O+ + HO- -> 2 H2O} (totale). À l'équivalence : $\mathrm{pH} = 7$ (\SI{25}{\celsius}).
\end{propriete}

\begin{methode}[Méthodes express à maîtriser pour le Bac]
\begin{enumerate}
  \item \textbf{Prouver qu'un acide/base est fort :} Comparer pH mesuré et pH théorique ($-\log C_a$ ou $14+\log C_b$). Égalité $\Rightarrow$ fort.
  \item \textbf{Inventaire espèces + concentrations :} Équation $\to$ Liste espèces $\to$ Électroneutralité + Conservation + $K_e$ $\to$ Calcul.
  \item \textbf{Dilution :} $C_i V_i = C_f V_f$. Prévision pH : Acide $\mathrm{pH}_f = \mathrm{pH}_i + \log(V_f/V_i)$ ; Base $\mathrm{pH}_f = \mathrm{pH}_i - \log(V_f/V_i)$. \textbf{Valide si $C_f > 10^{-6}$ M.}
  \item \textbf{Mélange Acide fort + Base forte :} $n(\ce{H3O+})$ vs $n(\ce{HO-})$ $\to$ Tableau d'avancement $\to$ Excès $\to$ Concentration dans $V_{\text{tot}}$ $\to$ pH.
\end{enumerate}
\end{methode}

\begin{aretenir}[Checklist avant le Bac]
\begin{itemize}
  \item $\square$ Je sais définir un acide fort et une base forte (réaction \textbf{totale} avec l'eau).
  \item $\square$ Je sais écrire les équations \ce{HCl + H2O -> H3O+ + Cl-} et \ce{NaOH -> Na+ + HO-}.
  \item $\square$ Je connais par cœur $\mathrm{pH} = -\log C_a$ pour un acide fort.
  \item $\square$ Je connais par cœur $\mathrm{pH} = 14 + \log C_b$ pour une base forte.
  \item $\square$ Je sais retrouver $C$ à partir du pH : $C_a = 10^{-\mathrm{pH}}$ et $C_b = 10^{\mathrm{pH}-14}$.
  \item $\square$ Je sais montrer qu'un acide est fort en comparant pH mesuré et $-\log C$.
  \item $\square$ Je sais faire l'inventaire des espèces en solution avec électroneutralité et conservation de la matière.
  \item $\square$ Je sais qu'une dilution au dixième fait varier le pH d'une unité (acide : $+1$ ; base : $-1$).
  \item $\square$ Je connais la limite de validité des formules : $C > \SI{1e-6}{mol.L^{-1}}$.
  \item $\square$ Je sais calculer le pH d'un mélange acide fort + base forte en identifiant le réactif en excès.
  \item $\square$ Je sais qu'à l'équivalence acide fort / base forte, $\mathrm{pH} = 7$ à \SI{25}{\celsius}.
  \item $\square$ Je connais les règles de sécurité et les pictogrammes des produits corrosifs (GHS05).
\end{itemize}
\end{aretenir}

\begin{exoresolu}[Exercice type Bac : Dilution et pH d'un acide fort]
\textbf{Énoncé :} On dispose d'une solution mère $S_0$ d'acide chlorhydrique de concentration molaire $C_0 = \SI{1,00}{mol.L^{-1}}$.
\begin{enumerate}
  \item On prélève \SI{10,0}{mL} de $S_0$ que l'on introduit dans une fiole jaugée de \SI{250,0}{mL} que l'on complète à l'eau distillée. On obtient la solution $S_1$. Calculer la concentration $C_1$ de $S_1$ et prévoir son pH.
  \item On prélève \SI{10,0}{mL} de $S_1$ que l'on dilue à \SI{250,0}{mL} pour obtenir $S_2$. Calculer $C_2$ et prévoir le pH de $S_2$.
  \item On prélève \SI{10,0}{mL} de $S_2$ que l'on dilue à \SI{250,0}{mL} pour obtenir $S_3$. Calculer $C_3$. La formule $\mathrm{pH} = -\log C_3$ est-elle encore valide ? Justifier. Donner la valeur approchée du pH de $S_3$.
\end{enumerate}
\tcblower
\textbf{Solution détaillée :}
\begin{enumerate}
  \item Dilution au 25ème ($d_1 = 25$). $C_1 = \frac{C_0}{25} = \frac{1,00}{25} = \SI{0,0400}{mol.L^{-1}} = \SI{4,00e-2}{mol.L^{-1}}$.
  $\mathrm{pH}_1 = -\log(4,00 \times 10^{-2}) = 2 - \log 4,00 = 2 - 0,60 = \mathbf{1,40}$.
  \item Dilution au 25ème à partir de $S_1$. Facteur total depuis $S_0$ : $d_{\text{tot}} = 25 \times 25 = 625$.
  $C_2 = \frac{C_0}{625} = \frac{1,00}{625} = \SI{1,60e-3}{mol.L^{-1}}$.
  $\mathrm{pH}_2 = -\log(1,60 \times 10^{-3}) = 3 - \log 1,60 = 3 - 0,20 = \mathbf{2,80}$.
  (Vérification : $\mathrm{pH}_1 + \log 25 = 1,40 + 1,40 = 2,80$).
  \item Dilution au 25ème à partir de $S_2$. Facteur total $d_{\text{tot}} = 25^3 = 15625$.
  $C_3 = \frac{1,00}{15625} = \SI{6,40e-5}{mol.L^{-1}}$.
  \textbf{Validité :} $C_3 = \SI{6,4e-5}{M} > \SI{1e-6}{M}$. La formule est \textbf{encore valide}.
  $\mathrm{pH}_3 = -\log(6,40 \times 10^{-5}) = 5 - \log 6,40 = 5 - 0,81 = \mathbf{4,19}$.
  \emph{Note :} Si on faisait une 4ème dilution au 25ème, $C_4 = \SI{2,56e-6}{M} > 10^{-6}$ (limite). Une 5ème donnerait $C_5 \approx \SI{1e-7}{M}$, formule invalide, pH $\approx 7$.
\end{enumerate}
\end{exoresolu}

\begin{exoresolu}[Exercice type Bac : Mélange acide fort / base forte et équivalence]
\textbf{Énoncé :} On verse \SI{20,0}{mL} d'une solution d'hydroxyde de sodium de concentration molaire $C_b = \SI{0,100}{mol.L^{-1}}$ dans un bécher. On ajoute progressivement une solution d'acide chlorhydrique de concentration molaire $C_a = \SI{0,100}{mol.L^{-1}}$ tout en mesurant le pH.
\begin{enumerate}
  \item Écrire l'équation de la réaction qui a lieu. Préciser si elle est totale.
  \item Calculer le volume $V_{E}$ d'acide chlorhydrique à verser pour atteindre l'équivalence.
  \item Calculer le pH du mélange pour $V_a = \SI{10,0}{mL}$ (avant équivalence).
  \item Calculer le pH du mélange pour $V_a = \SI{20,0}{mL}$ (à l'équivalence).
  \item Calculer le pH du mélange pour $V_a = \SI{30,0}{mL}$ (après équivalence).
  \item Tracer schématiquement la courbe $\mathrm{pH} = f(V_a)$.
\end{enumerate}
\tcblower
\textbf{Solution détaillée :}
\begin{enumerate}
  \item \ce{H3O+(aq) + HO-(aq) -> 2 H2O(l)}. Réaction \textbf{totale} ($K = 10^{14}$).
  \item À l'équivalence : $n(\ce{H3O+}) = n(\ce{HO-}) \Rightarrow C_a V_E = C_b V_b$.
  $V_E = \frac{C_b V_b}{C_a} = \frac{0,100 \times 20,0}{0,100} = \mathbf{20,0\ mL}$.
  \item $V_a = \SI{10,0}{mL} < V_E$. Excès de \ce{HO-}.
  $n_{\text{init}}(\ce{HO-}) = 0,100 \times 0,0200 = \SI{2,00e-3}{mol}$.
  $n_{\text{init}}(\ce{H3O+}) = 0,100 \times 0,0100 = \SI{1,00e-3}{mol}$.
  $n_{\text{excès}}(\ce{HO-}) = 2,00 \times 10^{-3} - 1,00 \times 10^{-3} = \SI{1,00e-3}{mol}$.
  $V_{\text{tot}} = 20,0 + 10,0 = \SI{30,0}{mL} = \SI{0,0300}{L}$.
  $[\ce{HO-}] = \frac{1,00 \times 10^{-3}}{0,0300} = \SI{3,33e-2}{mol.L^{-1}}$.
  $\mathrm{pH} = 14 + \log(3,33 \times 10^{-2}) = 14 - 1,48 = \mathbf{12,52}$.
  \item $V_a = \SI{20,0}{mL} = V_E$. Équivalence stoichiométrique.
  Solution ne contient que \ce{Na+}, \ce{Cl-}, \ce{H2O}. Neutre.
  $\mathrm{pH} = \mathbf{7,00}$ (à \SI{25}{\celsius}).
  \item $V_a = \SI{30,0}{mL} > V_E$. Excès de \ce{H3O+}.
  $n_{\text{init}}(\ce{H3O+}) = 0,100 \times 0,0300 = \SI{3,00e-3}{mol}$.
  $n_{\text{excès}}(\ce{H3O+}) = 3,00 \times 10^{-3} - 2,00 \times 10^{-3} = \SI{1,00e-3}{mol}$.
  $V_{\text{tot}} = 20,0 + 30,0 = \SI{50,0}{mL} = \SI{0,0500}{L}$.
  $[\ce{H3O+}] = \frac{1,00 \times 10^{-3}}{0,0500} = \SI{2,00e-2}{mol.L^{-1}}$.
  $\mathrm{pH} = -\log(2,00 \times 10^{-2}) = 2 - 0,30 = \mathbf{1,70}$.
  \item \textbf{Courbe schématique :} Partir de pH $\approx 13$ ($V_a=0$), descente lente, saut vertical brutal au voisinage de $V_E=20$ mL (pH 4 $\to$ 10), puis descente lente vers pH $\approx 1$ ($V_a$ grand).
\end{enumerate}
\end{exoresolu}

\begin{exercice}[Entraînement - Préparation au Bac]
\begin{enumerate}
  \item \textbf{QCM :} On dilue 10 fois une solution d'acide fort de pH = 3,00 (valide). Le nouveau pH est :
  A) 2,00 \quad B) 3,00 \quad C) 4,00 \quad D) 13,00
  \item \textbf{Calcul :} Quelle masse d'hydroxyde de sodium solide ($M = \SI{40,0}{g.mol^{-1}}$) faut-il dissoudre pour préparer \SI{500,0}{mL} de solution de pH = 13,00 à \SI{25}{\celsius} ?
  \item \textbf{Mélange :} On mélange \SI{50,0}{mL} d'HCl \SI{0,020}{mol.L^{-1}} et \SI{50,0}{mL} de NaOH \SI{0,030}{mol.L^{-1}}. Calculer le pH final.
  \item \textbf{Preuve :} Une solution d'acide nitrique \ce{HNO3} a une concentration $C = \SI{5,0e-4}{mol.L^{-1}}$. On mesure son pH : $\mathrm{pH}_{\text{mes}} = 3,30$. L'acide nitrique est-il un acide fort dans ces conditions ? Justifier par un calcul.
  \item \textbf{Dilution limite :} On dilue \SI{100}{fois} une solution de NaOH \SI{1e-4}{mol.L^{-1}}. Quel est le pH approché de la solution fille ? Expliquer pourquoi la formule $\mathrm{pH} = 14 + \log C$ ne s'applique pas.
\end{enumerate}
\end{exercice}

\begin{corrige}[Exercice n°1]
\begin{enumerate}
  \item \textbf{C} : Dilution au dixième $\Rightarrow$ pH augmente de 1 $\Rightarrow$ 3,00 + 1 = 4,00.
  \item $\mathrm{pH} = 13,00 \Rightarrow [\ce{HO-}] = 10^{\mathrm{pH}-14} = 10^{-1} = \SI{0,10}{mol.L^{-1}} = C_b$.
  $n(\ce{NaOH}) = C_b \times V = 0,10 \times 0,500 = \SI{0,050}{mol}$.
  $m = n \times M = 0,050 \times 40,0 = \mathbf{2,00\ g}$.
  \item $n(\ce{H3O+}) = 0,020 \times 0,050 = \SI{1,0e-3}{mol}$.
  $n(\ce{HO-}) = 0,030 \times 0,050 = \SI{1,5e-3}{mol}$.
  Excès \ce{HO-} = $0,5 \times 10^{-3}$ mol. $V_{\text{tot}} = \SI{0,100}{L}$.
  $[\ce{HO-}] = \SI{5,0e-3}{M}$. $\mathrm{pH} = 14 + \log(5,0 \times 10^{-3}) = 14 - 2,30 = \mathbf{11,70}$.
  \item $\mathrm{pH}_{\text{th}} = -\log(5,0 \times 10^{-4}) = 4 - \log 5,0 = 4 - 0,70 = 3,30$.
  $\mathrm{pH}_{\text{mes}} = 3,30$. Égalité parfaite $\Rightarrow$ \textbf{Oui, acide fort}.
  \item $C_{\text{mère}} = \SI{1e-4}{M} > \SI{1e-6}{M}$ : la relation $\mathrm{pH} = 14 + \log C$ s'applique à la solution mère (pH $=10$).
  Dilution 100 fois $\Rightarrow C_{\text{fille}} = \SI{1e-6}{M}$. C'est la \textbf{limite}.
  $[\ce{HO-}]_{\text{apportée}} = 10^{-6}$ M. $[\ce{HO-}]_{\text{eau}} = 10^{-7}$ M. Contribution eau non négligeable (10\%).
  $[\ce{HO-}]_{\text{total}} \approx 1,1 \times 10^{-6}$ M. $\mathrm{pOH} \approx 5,96 \Rightarrow \mathrm{pH} \approx \mathbf{8,04}$.
  Formule $\mathrm{pH} = 14 + \log(10^{-6}) = 8,00$ donne une approximation acceptable mais rigoureusement fausse car $C \not > 10^{-6}$ (égalité). Pour $C < 10^{-6}$, pH $\to$ 7.
\end{enumerate}
\end{corrige}

\findecours

\end{document}
